% Newton's Laws of Motion — Pure Mathematics with Mechanics Upper Sixth — Cours signé OnBuch+
% Official MINESEC syllabus. Compiler avec : tectonic PMM16-newton-s-laws-of-motion.tex
\newcommand{\DOCMATIERE}{Pure Mathematics with Mechanics}
\newcommand{\DOCNIVEAU}{Upper Sixth}
\newcommand{\DOCMODULE}{Upper Sixth — Pure mathematics with mechanics}
\newcommand{\DOCLECON}{16}
\newcommand{\DOCTITRE}{Newton's Laws of Motion}
% OnBuch+ — préambule « cours signés » ENRICHI (compilé avec tectonic / XeLaTeX)
% Système visuel « Pop » d'OnBuch+ : fond crème (jamais blanc), bordures encre
% épaisses, ombres « dures » décalées, titres Archivo Black en capitales.
% Version MATHÉMATIQUES : graphes (pgfplots/tikz), boîtes pédagogiques
% (définition, propriété, méthode, attention, le savais-tu, exercice résolu,
% exercice, TP, bilan), page de garde et signature OnBuch+ ancrée en bas.
%
% Macros attendues dans main.tex AVANT \input{preamble} :
%   \DOCMATIERE  \DOCNIVEAU  \DOCMODULE  \DOCLECON (numéro)  \DOCTITRE
\documentclass[11pt]{article}

\usepackage{fontspec}
\usepackage[english]{babel}
\usepackage[a4paper,margin=2cm,top=3.4cm,bottom=2.8cm,headheight=1.4cm,footskip=1.3cm]{geometry}
\usepackage{amsmath,amssymb,amsfonts}
\usepackage{mathtools} % \xmapsto, \coloneqq... souvent utilisés par les modèles
\usepackage{cancel} % simplifications barrées (bilans de forces)
% \abs{z} (module, valeur absolue) et \norm{u} (norme), utilisés par les modèles
\providecommand{\abs}{}\let\abs\relax\DeclarePairedDelimiter{\abs}{\lvert}{\rvert}
\providecommand{\norm}{}\let\norm\relax\DeclarePairedDelimiter{\norm}{\lVert}{\rVert}
\usepackage[table]{xcolor} % [table] : fournit aussi \rowcolors (alternance de lignes)
\usepackage{pagecolor}
\usepackage{tikz}
\usetikzlibrary{arrows.meta,positioning,calc,shapes.geometric,shapes.misc,shapes.arrows,decorations.pathmorphing,decorations.pathreplacing,decorations.markings,patterns,shadows,mindmap,trees,fit,backgrounds,matrix,intersections,angles,quotes}
\usepackage{pgfplots}
\usepackage[version=4]{mhchem} % équations nucléaires \ce{^{235}_{92}U}
\usepackage[european,siunitx]{circuitikz} % schémas électriques
\pgfplotsset{compat=1.18}
\usepgfplotslibrary{fillbetween,groupplots} % aires entre courbes (calcul intégral)
\usepackage{enumitem}
\usepackage{fancyhdr}
% [most] charge toutes les bibliothèques tcolorbox (listings, minted, raster,
% documentation...) et gonfle inutilement la mémoire TeX sur un document long
% (~300 boîtes) : on ne charge que ce qui est réellement utilisé.
\usepackage[skins,breakable]{tcolorbox}
\usepackage{graphicx}
\usepackage{colortbl}
\usepackage{booktabs}
\usepackage{tabularx}
\usepackage{diagbox} % case d'en-tête diagonale des tableaux à double entrée
% Colonnes extensibles centrée / gauche / droite (C, L, R), souvent utilisées par les modèles
\newcolumntype{C}{>{\centering\arraybackslash}X}
\newcolumntype{L}{>{\raggedright\arraybackslash}X}
\newcolumntype{R}{>{\raggedleft\arraybackslash}X}
\usepackage{array}
\usepackage{multicol}
\usepackage{siunitx}
% parse-numbers=false : accepte aussi \SI{2,5 \times 10^{-3}}{...}
\sisetup{locale=UK,output-decimal-marker={.},group-minimum-digits=4,parse-numbers=false}
\DeclareSIUnit\ohm{\ensuremath{\symup{\Omega}}} % U+2126 absent de la police
\DeclareSIUnit\division{div} % réglages d'oscilloscope (ms/div, V/div)
\DeclareSIUnit\tr{tr} % tours (vitesse de rotation, ex. \SI{1500}{\tr\per\minute})
\usepackage{qrcode}
% \degree : commande fréquemment utilisée par les modèles pour un angle ou une
% température en dehors de \SI{}{} (non fournie par les paquets chargés).
\providecommand{\degree}{^{\circ}}
% \qed (fin de démonstration), utilisé par les modèles sans amsthm
\providecommand{\qed}{\hfill$\square$}
% \barre : double barre des valeurs interdites dans un tableau de variations (tkz-tab)
\providecommand{\barre}{\ensuremath{\|}}
% environnement proof (amsthm non chargé) utilisé par les modèles
\ifdefined\proof\else\newenvironment{proof}[1][Proof]{\par\noindent\textit{#1.}\ }{\qed\par}\fi
\usepackage{lastpage}
\usepackage{hyperref}
\hypersetup{hidelinks}

% ============================================================
% Palette « Pop » OnBuch+ — mêmes hex que la classe P (plus_ui.dart)
% ============================================================
\definecolor{poporange}{HTML}{F59321}
\definecolor{popdark}{HTML}{E07A0C}
\definecolor{popcream}{HTML}{FFF6E8}
\definecolor{popink}{HTML}{1C1714}
\definecolor{poppurple}{HTML}{7A5AE0}
\definecolor{popgreen}{HTML}{2E9E6A}
\definecolor{poppink}{HTML}{E0457B}
\definecolor{popblue}{HTML}{2F7DE1}
\definecolor{popgold}{HTML}{C98A12}
\definecolor{popmuted}{HTML}{8A7F76}
\definecolor{popline}{HTML}{EADFD2}
\definecolor{popgray}{HTML}{8A7F76}
\colorlet{popgrey}{popgray}
% Alias tolérants (le modèle nomme parfois les couleurs différemment)
\colorlet{popwhite}{white}
\colorlet{popblack}{popink}
\colorlet{popred}{poppink}
\colorlet{popyellow}{popgold}
% Teintes claires pour fonds de tableaux / zones
\colorlet{popblueL}{popblue!10!white}
\colorlet{poporangeL}{poporange!14!white}
\colorlet{popgreenL}{popgreen!12!white}
\colorlet{poppurpleL}{poppurple!12!white}
\colorlet{poppinkL}{poppink!12!white}
\colorlet{popgoldL}{popgold!14!white}

\pagecolor{popcream}

% ============================================================
% Typographie
% ============================================================
\defaultfontfeatures{Ligatures=TeX}
\setmainfont{PlusJakartaSans-Medium.ttf}[Path=../../../fonts/,BoldFont=PlusJakartaSans-Bold.ttf]
% Maths en sans-serif (Fira Math) avec chiffres et lettres droites de Plus
% Jakarta, pour que les formules se fondent dans le texte.
\usepackage[math-style=ISO,bold-style=ISO]{unicode-math}
\setmathfont{FiraMath-Regular.otf}[Path=../../../fonts/]
\setmathfont{PlusJakartaSans-Medium.ttf}[Path=../../../fonts/,range={up/num,up/{Latin,latin},\mathup}]
% (icomma retiré : décimales avec point en anglais)
% Caractères mathématiques Unicode absents des polices de texte (Plus Jakarta,
% Archivo Black) : rendus par leur équivalent mathématique, en texte comme en
% formule — sinon ils s'affichent en carré vide (ex. « Arithmétique dans ℤ »).
\usepackage{newunicodechar}
\newunicodechar{ℝ}{\ensuremath{\mathbb{R}}}
\newunicodechar{ℤ}{\ensuremath{\mathbb{Z}}}
\newunicodechar{ℂ}{\ensuremath{\mathbb{C}}}
\newunicodechar{ℕ}{\ensuremath{\mathbb{N}}}
\newunicodechar{ℚ}{\ensuremath{\mathbb{Q}}}
\newunicodechar{𝔻}{\ensuremath{\mathbb{D}}}
\newunicodechar{α}{\ensuremath{\alpha}}
\newunicodechar{β}{\ensuremath{\beta}}
\newunicodechar{γ}{\ensuremath{\gamma}}
\newunicodechar{δ}{\ensuremath{\delta}}
\newunicodechar{ε}{\ensuremath{\varepsilon}}
\newunicodechar{θ}{\ensuremath{\theta}}
\newunicodechar{λ}{\ensuremath{\lambda}}
\newunicodechar{μ}{\ensuremath{\mu}}
\newunicodechar{σ}{\ensuremath{\sigma}}
\newunicodechar{τ}{\ensuremath{\tau}}
\newunicodechar{φ}{\ensuremath{\varphi}}
\newunicodechar{ψ}{\ensuremath{\psi}}
\newunicodechar{ω}{\ensuremath{\omega}}
\newunicodechar{Σ}{\ensuremath{\Sigma}}
% majuscules produites par \MakeUppercase dans les titres (α-aminés -> Α)
\newunicodechar{Α}{\ensuremath{\alpha}}
\newunicodechar{Β}{\ensuremath{\beta}}
\newunicodechar{Γ}{\ensuremath{\Gamma}}
\newunicodechar{Δ}{\ensuremath{\Delta}}
\newunicodechar{Ε}{\ensuremath{\varepsilon}}
\newunicodechar{Θ}{\ensuremath{\Theta}}
\newunicodechar{Λ}{\ensuremath{\Lambda}}
\newunicodechar{Μ}{\ensuremath{\mu}}
\newunicodechar{Τ}{\ensuremath{\tau}}
\newunicodechar{Φ}{\ensuremath{\Phi}}
\newunicodechar{Ψ}{\ensuremath{\Psi}}
\newunicodechar{Ω}{\ensuremath{\Omega}}
\newunicodechar{↦}{\ensuremath{\mapsto}}
\newunicodechar{∈}{\ensuremath{\in}}
\newunicodechar{∉}{\ensuremath{\notin}}
\newunicodechar{∧}{\ensuremath{\wedge}}
\newunicodechar{⇔}{\ensuremath{\Leftrightarrow}}
\newunicodechar{⟺}{\ensuremath{\Longleftrightarrow}}
\newunicodechar{⇒}{\ensuremath{\Rightarrow}}
\newunicodechar{⟹}{\ensuremath{\Longrightarrow}}
\newunicodechar{∘}{\ensuremath{\circ}}
\newunicodechar{∪}{\ensuremath{\cup}}
\newunicodechar{∩}{\ensuremath{\cap}}
\newunicodechar{≡}{\ensuremath{\equiv}}
\newunicodechar{⊂}{\ensuremath{\subset}}
\newunicodechar{⊥}{\ensuremath{\perp}}
\newunicodechar{∃}{\ensuremath{\exists}}
\newunicodechar{∀}{\ensuremath{\forall}}
\newunicodechar{∛}{\ensuremath{\sqrt[3]{\,}}}
\newunicodechar{∜}{\ensuremath{\sqrt[4]{\,}}}
\newunicodechar{⃗}{\ensuremath{{}^{\to}}}
\newunicodechar{ⁿ}{\ifmmode^{n}\else\textsuperscript{\ensuremath{n}}\fi}
\newunicodechar{ˣ}{\ifmmode^{x}\else\textsuperscript{\ensuremath{x}}\fi}
\newunicodechar{ᵇ}{\ifmmode^{b}\else\textsuperscript{\ensuremath{b}}\fi}
\newunicodechar{ʳ}{\ifmmode^{r}\else\textsuperscript{\ensuremath{r}}\fi}
\newunicodechar{ᵘ}{\ifmmode^{u}\else\textsuperscript{\ensuremath{u}}\fi}
\newunicodechar{ʸ}{\ifmmode^{y}\else\textsuperscript{\ensuremath{y}}\fi}
\newunicodechar{ᵃ}{\ifmmode^{a}\else\textsuperscript{\ensuremath{a}}\fi}
\newunicodechar{ᵉ}{\ifmmode^{e}\else\textsuperscript{\ensuremath{e}}\fi}
\newunicodechar{ᵏ}{\ifmmode^{k}\else\textsuperscript{\ensuremath{k}}\fi}
\newunicodechar{ᵖ}{\ifmmode^{p}\else\textsuperscript{\ensuremath{p}}\fi}
\newunicodechar{ᴺ}{\ifmmode^{N}\else\textsuperscript{\ensuremath{N}}\fi}
\newunicodechar{ᶜ}{\ifmmode^{c}\else\textsuperscript{\ensuremath{c}}\fi}
\newunicodechar{ʰ}{\ifmmode^{h}\else\textsuperscript{\ensuremath{h}}\fi}
\newunicodechar{ᵠ}{\ifmmode^{q}\else\textsuperscript{\ensuremath{q}}\fi}
\newunicodechar{ₙ}{\ifmmode_{n}\else\textsubscript{\ensuremath{n}}\fi}
\newunicodechar{ₐ}{\ifmmode_{a}\else\textsubscript{\ensuremath{a}}\fi}
\newunicodechar{ᵢ}{\ifmmode_{i}\else\textsubscript{\ensuremath{i}}\fi}
\newunicodechar{ᵣ}{\ifmmode_{r}\else\textsubscript{\ensuremath{r}}\fi}
\newunicodechar{ₖ}{\ifmmode_{k}\else\textsubscript{\ensuremath{k}}\fi}
\newunicodechar{ᵦ}{\ifmmode_{\beta}\else\textsubscript{\ensuremath{\beta}}\fi}
\newunicodechar{ₓ}{\ifmmode_{x}\else\textsubscript{\ensuremath{x}}\fi}
\newfontfamily\popdisplay{ArchivoBlack-Regular.ttf}[Path=../../../fonts/]
\newfontfamily\poplogo{SpaceGrotesk-Bold.ttf}[Path=../../../fonts/]
\newfontfamily\popsemi{PlusJakartaSans-SemiBold.ttf}[Path=../../../fonts/]
\newfontfamily\popxbold{PlusJakartaSans-ExtraBold.ttf}[Path=../../../fonts/]
\newfontfamily\popxboldls{PlusJakartaSans-ExtraBold.ttf}[Path=../../../fonts/,LetterSpace=3.0]

\setlist[enumerate]{leftmargin=1.7em,itemsep=5pt,topsep=4pt}
\setlist[itemize]{leftmargin=1.7em,itemsep=4pt,topsep=4pt,label={\color{poporange}\textbullet}}
\setlength{\parindent}{0pt}
\setlength{\parskip}{5pt}
\renewcommand{\arraystretch}{1.25}

% pgfplots : style Pop par défaut
\pgfplotsset{
  every axis/.append style={axis line style={popink,line width=1pt},tick style={popink},
    grid style={popline,line width=0.5pt},label style={font=\small},tick label style={font=\footnotesize}},
  popcurve/.style={poporange,line width=1.8pt,no markers,smooth},
}
% Même style disponible dans un \draw TikZ simple (hors pgfplots)
\tikzset{popcurve/.style={poporange,line width=1.8pt}}
\tikzset{popgrid/.style={popline,very thin}}
\colorlet{popcurve}{poporange} % le nom du style est parfois utilisé comme couleur

% ============================================================
% Pill / squiggle
% ============================================================
\newcommand{\poppill}[3]{%
  \tikz[baseline=-0.55ex]\node[draw=popink,line width=1.2pt,rounded corners=2.6pt,%
    fill=#2,inner xsep=6.5pt,inner ysep=3pt]{%
    \popxboldls\fontsize{7}{7}\selectfont\color{#3}\MakeUppercase{#1}};%
}
\newcommand{\popsquiggle}[1]{%
  \tikz[baseline=-0.4ex]{
    \draw[#1,line width=2.6pt,line cap=round]
      plot [smooth] coordinates {(0,0.09) (0.34,0.01) (0.68,0.21) (1.02,0.01) (1.36,0.10)};
  }%
}
% Mot-clé surligné (marqueur orange sous le texte)
\newcommand{\cle}[1]{{\popxbold\color{popdark}#1}}

% ============================================================
% En-tête / pied
% ============================================================
\pagestyle{fancy}
\fancyhf{}
\renewcommand{\headrulewidth}{0pt}
\renewcommand{\footrulewidth}{0pt}
\lhead{\raisebox{-0.15ex}{\poplogo\fontsize{13}{13}\selectfont\color{popink}OnBuch\color{poporange}+}}
\rhead{\poppill{\DOCMATIERE\ \textbullet\ \DOCNIVEAU\ \textbullet\ Chapter \DOCLECON}{white}{popink}}
\lfoot{\popxbold\fontsize{7.6}{7.6}\selectfont\color{popmuted}\MakeUppercase{Signed course OnBuch+}\ \textbullet\ {\popsemi\DOCTITRE}}
\rfoot{\tikz[baseline=-0.6ex]\node[rounded corners=8.5pt,fill=popink,minimum height=17pt,minimum width=17pt,inner xsep=5pt,inner sep=0]{\popxbold\fontsize{8}{8}\selectfont\color{popcream}\thepage\,/\,\pageref*{LastPage}};}

% ============================================================
% Titres
% ============================================================
\newcommand{\chaptitre}[2]{%
  \begin{tcolorbox}[enhanced,colback=poporange,colframe=popink,boxrule=2.6pt,arc=11pt,
      drop shadow=popink,left=15pt,right=15pt,top=12pt,bottom=11pt,before skip=3pt,after skip=18pt]
    \poppill{#2}{popink}{popcream}\par\vspace{9pt}
    {\raggedright\hyphenpenalty=10000\exhyphenpenalty=10000\popdisplay\fontsize{22}{24}\selectfont\color{popcream}\MakeUppercase{#1}\par}
  \end{tcolorbox}
}
% Section numérotée : pastille orange avec le numéro + titre Archivo
\newcounter{popsec}
\newcommand{\coursec}[1]{%
  \stepcounter{popsec}\setcounter{popsub}{0}%
  \par\vspace{16pt}\noindent
  \begin{minipage}{\linewidth}
  \tikz[baseline=-0.7ex]\node[draw=popink,line width=1.6pt,fill=poporange,rounded corners=4pt,minimum size=22pt,inner sep=0]{\popdisplay\fontsize{12}{12}\selectfont\color{popcream}\thepopsec};%
  \hspace{8pt}{\popdisplay\fontsize{15}{17}\selectfont\color{popink}\MakeUppercase{#1}}\par
  \vspace{4pt}\noindent{\color{poporange}\rule{36pt}{3.6pt}}
  \end{minipage}\par\nopagebreak\vspace{8pt}
}
\newcounter{popsub}
\newcommand{\courssub}[1]{%
  \stepcounter{popsub}%
  \par\vspace{9pt}\noindent{\popxbold\fontsize{11.5}{13}\selectfont\color{popink}{\color{poporange}\thepopsec.\thepopsub}\hspace{6pt}#1}\par\nopagebreak\vspace{4pt}
}

% ============================================================
% Boîtes pédagogiques — même recette : fond blanc, bordure épaisse colorée,
% ombre dure encre, badge plein. Titre optionnel [#1].
% ============================================================
\tcbset{popbase/.style={breakable,enhanced,colback=white,boxrule=2.3pt,arc=9pt,
  shadow={3pt}{-3pt}{0pt}{popink},left=14pt,right=14pt,top=11pt,bottom=11pt,before skip=11pt,after skip=15pt,
  fontupper=\popsemi\fontsize{10.6}{14.5}\selectfont\color{popink},
  fontlower=\popsemi\fontsize{10.6}{14.5}\selectfont\color{popink}}}
% Coche dessinée (le glyphe ✓ n'existe pas dans les polices du document)
\newcommand{\popcheck}[1]{\tikz[baseline=-0.5ex]\draw[#1,line width=1.6pt,line cap=round,line join=round](-0.13,0.0)--(-0.04,-0.09)--(0.13,0.1);}
\providecommand{\coche}{\popcheck{popgreen}}
\providecommand{\resultat}[1]{\par\smallskip\noindent\fcolorbox{popgreen}{popgreenL}{\parbox{\dimexpr\linewidth-2\fboxsep-2\fboxrule\relax}{\textbf{Result.} #1}}\par\smallskip}
\newcommand{\popboxhead}[3]{\poppill{#1}{#2}{white}\ifx\relax#3\relax\else\hspace{6pt}{\popxbold\fontsize{10.6}{12}\selectfont\color{popink}#3}\fi\par\vspace{7pt}}

\newtcolorbox{aretenir}[1][]{popbase,colframe=popgreen,before upper={\popboxhead{Key points}{popgreen}{#1}}}
\newtcolorbox{exemplebox}[1][]{popbase,colframe=poporange,before upper={\popboxhead{Exemple}{poporange}{#1}}}
\newtcolorbox{definition}[1][]{popbase,colframe=popblue,before upper={\popboxhead{Definition}{popblue}{#1}}}
\newtcolorbox{propriete}[1][]{popbase,colframe=poppurple,before upper={\popboxhead{Property}{poppurple}{#1}}}
\newtcolorbox{methode}[1][]{popbase,colframe=popgold,before upper={\popboxhead{Method}{popgold}{#1}}}
\newtcolorbox{attention}[1][]{popbase,colframe=poppink,before upper={\popboxhead{Warning}{poppink}{#1}}}
\newtcolorbox{experience}[1][]{popbase,colframe=popblue,colback=popblueL,before upper={\popboxhead{Experiment}{popblue}{#1}}}
% « Le savais-tu ? » : carte violette pleine (contexte, culture, Cameroun)
\newtcolorbox{savaistu}[1][]{popbase,colback=poppurple,colframe=popink,
  fontupper=\popsemi\fontsize{10.4}{14.2}\selectfont\color{white},
  before upper={\poppill{Did you know?}{white}{poppurple}\ifx\relax#1\relax\else\hspace{6pt}{\popxbold\color{white}#1}\fi\par\vspace{7pt}}}
% Exercice résolu : énoncé en haut, \tcblower puis la solution sur fond crème
\newcounter{popexo}\newcounter{popexr}
\newtcolorbox{exoresolu}[1][]{popbase,colframe=poporange,colbacklower=popcream,
  segmentation style={popink,line width=1.2pt,dashed},
  before upper={\stepcounter{popexr}\popboxhead{Worked example \thepopexr}{poporange}{#1}},
  before lower={\poppill{Solution}{popink}{popcream}\par\vspace{6pt}}}
% Exercice (non résolu en place) — la correction est dans la partie Corrigés
\newtcolorbox{exercice}[1][]{popbase,colframe=popink,
  before upper={\stepcounter{popexo}\popboxhead{Exercise \thepopexo}{popink}{#1}}}
% Corrigé
\newtcolorbox{corrige}[1][]{popbase,colframe=popgreen,colback=popgreenL,
  before upper={\popboxhead{Solution}{popgreen}{#1}}}
% Objectifs / prérequis / bilan
\newtcolorbox{objectifs}{popbase,colframe=popink,colback=poporangeL,
  before upper={\popboxhead{Chapter objectives}{popink}{}}}
\newtcolorbox{prerequis}{popbase,colframe=popink,colback=popblueL,
  before upper={\popboxhead{Prerequisites}{popblue}{}}}
\newtcolorbox{bilan}{popbase,colframe=popink,colback=white,boxrule=2.8pt,
  before upper={\popboxhead{Fiche bilan}{poporange}{L'essentiel à connaître pour l'examen}}}
% Figure encadrée avec légende
\newtcolorbox{popfigure}[1][]{enhanced,breakable=false,colback=white,colframe=popline,boxrule=1.4pt,arc=8pt,
  left=10pt,right=10pt,top=10pt,bottom=8pt,before skip=10pt,after skip=12pt,halign=center,
  lower separated=false,halign lower=center,
  fontlower=\popsemi\fontsize{9}{11.5}\selectfont\color{popmuted},
  #1}
\newcommand{\legende}[1]{\par\vspace{6pt}{\popsemi\fontsize{9}{11.5}\selectfont\color{popmuted}#1}}

% ============================================================
% Signature OnBuch+
% ============================================================
\newcommand{\poplogoinline}[1]{{\poplogo\fontsize{#1}{#1}\selectfont\color{popink}OnBuch\color{poporange}+}}
\newcommand{\signatureonbuch}{%
  \begin{tcolorbox}[enhanced,colback=popink,colframe=popink,boxrule=2.2pt,arc=10pt,
      drop shadow=poporange,left=14pt,right=14pt,top=10pt,bottom=10pt,before skip=0pt,after skip=0pt]
    \tikz[baseline=-0.7ex]\node[circle,fill=popgreen,draw=popcream,line width=1.3pt,minimum size=22pt,inner sep=0]{\popcheck{white}};%
    \hspace{9pt}\begin{minipage}[c]{0.62\linewidth}
      {\popxbold\fontsize{12}{13}\selectfont\color{popcream}Signed\ \textbullet\ {\poplogo OnBuch\color{poporange}+}}\par\vspace{2pt}
      {\raggedright\popsemi\fontsize{8}{10.5}\selectfont\color{popline}\MakeUppercase{OnBuch+ teaching team}\\\MakeUppercase{Follows the official MINESEC syllabus}\par}
    \end{minipage}\hfill
    {\poplogo\fontsize{22}{22}\selectfont\color{popcream}OnBuch\color{poporange}+}
  \end{tcolorbox}%
}

% ============================================================
% Page de garde
% #1 titre · #2 matière · #3 niveau · #4 module · #5 n° leçon · #6 durée ·
% #7 description / objectifs (texte ou itemize)
% ============================================================
\newcommand{\pagegarde}[7]{%
  \thispagestyle{empty}
  \newgeometry{a4paper,margin=1.7cm,top=1.6cm,bottom=1.4cm}%
  \begin{tikzpicture}[remember picture,overlay]
    \fill[poporange!18] ([xshift=-3.2cm,yshift=-3.6cm]current page.north east) circle (4.6cm);
    \fill[poppurple!14] ([xshift=2.4cm,yshift=7.6cm]current page.south west) circle (3.8cm);
  \end{tikzpicture}%
  {\poplogo\fontsize{30}{30}\selectfont\color{popink}OnBuch\color{poporange}+}\hfill
  \raisebox{0.6ex}{\poppill{Signed course \textbullet\ Premium}{popink}{popcream}}\par
  \vspace{1pt}\popsquiggle{poppurple}\par
  \vspace{8pt}
  \begin{tcolorbox}[enhanced,colback=poporange,colframe=popink,boxrule=2.8pt,arc=13pt,
      drop shadow=popink,left=18pt,right=18pt,top=12pt,bottom=13pt,after skip=10pt]
    \poppill{#2 \textbullet\ #3}{popink}{popcream}\hspace{5pt}\poppill{#4}{white}{popink}\par\vspace{8pt}
    {\popdisplay\fontsize{14}{14}\selectfont\color{popink}CHAPTER #5}\par\vspace{4pt}
    {\raggedright\hyphenpenalty=10000\exhyphenpenalty=10000\popdisplay\fontsize{25}{27}\selectfont\color{popcream}\MakeUppercase{#1}\par}
    \vspace{8pt}
    {\popxbold\fontsize{9.5}{9.5}\selectfont\color{popink}\MakeUppercase{Suggested time: #6}}
  \end{tcolorbox}
  \begin{tcolorbox}[enhanced,colback=white,colframe=popink,boxrule=2.2pt,arc=10pt,
      drop shadow=popink,left=16pt,right=16pt,top=10pt,bottom=10pt,after skip=8pt]
    \poppill{In this chapter}{poppurple}{white}\par\vspace{5pt}
    \popsemi\fontsize{9.3}{12.6}\selectfont\color{popink}#7
  \end{tcolorbox}
  \noindent\begin{minipage}[t]{0.487\linewidth}
    \begin{tcolorbox}[enhanced,colback=popgreen,colframe=popink,boxrule=2.2pt,arc=10pt,
        drop shadow=popink,left=11pt,right=11pt,top=10pt,bottom=10pt,height=3.1cm,valign=center]
      \tcbox[colback=white,colframe=popink,boxrule=1.3pt,arc=5pt,boxsep=0pt,nobeforeafter,
        left=3pt,right=3pt,top=3pt,bottom=3pt,valign=center]{\qrcode[height=1.9cm]{https://play.google.com/store/apps/details?id=com.luvvix.onbuch&hl=en}}%
      \hspace{8pt}\begin{minipage}[c]{0.5\linewidth}
        {\popdisplay\fontsize{11}{12.5}\selectfont\color{white}\MakeUppercase{Download OnBuch}}\par\vspace{4pt}
        {\raggedright\popsemi\fontsize{8.4}{11}\selectfont\color{white}Free \textbullet\ Google Play\\AI tutor, past papers, results\par}
      \end{minipage}
    \end{tcolorbox}
  \end{minipage}\hfill
  \begin{minipage}[t]{0.487\linewidth}
    \begin{tcolorbox}[enhanced,colback=poppurple,colframe=popink,boxrule=2.2pt,arc=10pt,
        drop shadow=popink,left=11pt,right=11pt,top=10pt,bottom=10pt,height=3.1cm,valign=center]
      \tikz[baseline=-0.7ex]\node[circle,fill=white,draw=popink,line width=1.3pt,minimum size=1.6cm,inner sep=0]{\popdisplay\fontsize{24}{24}\selectfont\color{poppurple}+};%
      \hspace{8pt}\begin{minipage}[c]{0.6\linewidth}
        {\popdisplay\fontsize{11}{12.5}\selectfont\color{white}\MakeUppercase{Go OnBuch+}}\par\vspace{4pt}
        {\raggedright\popsemi\fontsize{8.4}{11}\selectfont\color{white}All signed courses, unlimited AI tutor, PDF sheets — OnBuch+ tab in the app\par}
      \end{minipage}
    \end{tcolorbox}
  \end{minipage}
  \vspace{8pt}
  \popcontenu
  \vfill
  \signatureonbuch
  \restoregeometry
  \newpage
}

% Rangée « Ce cours contient » (page de garde)
\newcommand{\poptile}[3]{% #1 couleur #2 gros texte #3 légende
  \begin{tcolorbox}[enhanced,colback=#1,colframe=popink,boxrule=2pt,arc=9pt,shadow={2.5pt}{-2.5pt}{0pt}{popink},
      left=6pt,right=6pt,top=9pt,bottom=9pt,halign=center,valign=center,height=1.75cm,width=\linewidth,nobeforeafter]
    {\popdisplay\fontsize{12.5}{13}\selectfont\color{white}\MakeUppercase{#2}}\par\vspace{3pt}
    {\centering\popsemi\fontsize{7.8}{9.5}\selectfont\color{white}#3\par}
  \end{tcolorbox}}
\newcommand{\popcontenu}{%
  \par\noindent{\popxboldls\fontsize{7.5}{7.5}\selectfont\color{popmuted}THIS SIGNED COURSE CONTAINS}\par\vspace{7pt}
  \noindent\begin{minipage}[t]{0.235\linewidth}\poptile{popblue}{Course}{complete, illustrated, step by step}\end{minipage}\hfill
  \begin{minipage}[t]{0.235\linewidth}\poptile{poporange}{Methods}{and worked examples}\end{minipage}\hfill
  \begin{minipage}[t]{0.235\linewidth}\poptile{poppink}{Exercises}{exam style + solutions}\end{minipage}\hfill
  \begin{minipage}[t]{0.235\linewidth}\poptile{popgreen}{Summary sheet}{the essentials to remember}\end{minipage}\par}

% Sommaire
\newtcolorbox{sommaire}{enhanced,colback=white,colframe=popink,boxrule=2.2pt,arc=10pt,
  drop shadow=popink,left=18pt,right=18pt,top=13pt,bottom=13pt,before skip=6pt,after skip=20pt,
  before upper={\poppill{Contents}{poppurple}{white}\par\vspace{9pt}\popsemi\fontsize{10.6}{14.5}\selectfont\color{popink}}}

% Signature de fin de cours — poussée tout en bas de la dernière page
\newcommand{\findecours}{%
  \par\vfill\nopagebreak
  \begin{center}\popsquiggle{poporange}\end{center}\vspace{4pt}
  \signatureonbuch
}
\AtBeginDocument{\def\llbracket{\mathopen{[\mkern-3.5mu[}}\def\rrbracket{\mathclose{]\mkern-3.5mu]}}} % intervalles d'entiers (glyphe ⟦ absent de la police)
% Glyphes absents de Fira Math : redéfinis à partir de glyphes disponibles
\AtBeginDocument{%
  \def\setminus{\mathbin{\backslash}}\let\smallsetminus\setminus
  \def\vdots{\mathord{\vbox{\baselineskip3.2pt\lineskiplimit0pt\kern4pt\hbox{.}\hbox{.}\hbox{.}}}}%
  \def\ddots{\mathinner{\mkern1mu\raise6.5pt\hbox{.}\mkern2mu\raise3.6pt\hbox{.}\mkern2mu\raise0.7pt\hbox{.}\mkern1mu}}%
  \def\triangle{\mathord{\scalebox{0.9}{$\symup{\Delta}$}}}%
  \def\nabla{\mathord{\raisebox{\depth}{\scalebox{1}[-1]{$\symup{\Delta}$}}}}%
  \def\checkmark{\popcheck{popdark}}%
  \def\star{\mathbin{\ast}}\def\bigstar{\mathord{\ast}}%
  \def\diamond{\mathbin{\scalebox{0.6}{$\Diamond$}}}%
  \def\bigcup{\mathop{\mathchoice{\raisebox{-0.3em}{\scalebox{1.7}{$\cup$}}}{\scalebox{1.25}{$\cup$}}{\cup}{\cup}}\displaylimits}%
  \def\bigcap{\mathop{\mathchoice{\raisebox{-0.3em}{\scalebox{1.7}{$\cap$}}}{\scalebox{1.25}{$\cap$}}{\cap}{\cap}}\displaylimits}%
  \def\lesseqqgtr{\mathrel{\lessgtr}}\def\bigoplus{\mathop{\oplus}\displaylimits}%
}
\providecommand{\ssi}{if and only if} % abréviation fréquente
\AtBeginDocument{\providecommand{\wideparen}{\overparen}} % arc de cercle

%% FIGFIX-BEGIN v4 — correctifs globaux des figures (docs/onbuch-generation/tools/figfix)
\usepackage{adjustbox}\usepackage{etoolbox}
% toute figure TikZ/pgfplots plus large que la ligne est réduite à la largeur disponible
\BeforeBeginEnvironment{tikzpicture}{\begin{adjustbox}{max width=\linewidth}}
\AfterEndEnvironment{tikzpicture}{\end{adjustbox}}
% légendes pgfplots lisibles par-dessus les courbes
\pgfplotsset{every axis/.append style={legend style={fill=white,fill opacity=0.92,text opacity=1,draw=black!15,inner sep=3pt}}}
% étiquettes d'arêtes (syntaxe edge["..."]) avec fond blanc
\tikzset{every edge quotes/.append style={fill=white,inner sep=1.2pt}}
%% FIGFIX-END
\begin{document}

\pagegarde{Newton's Laws of Motion}{Pure Mathematics with Mechanics}{Upper Sixth}{Upper Sixth — Pure mathematics with mechanics}{16}{5 periods}{This chapter introduces Newton's three laws of motion, the cornerstone of classical mechanics. Students will learn to analyse forces, predict accelerations, and understand the limits of these laws in modern physics.
\vspace{4pt}
\begin{multicols}{2}\small
\begin{itemize}
\item State and explain Newton's first law (law of inertia) and identify inertial reference frames.
\item Apply Newton's second law (F = ma) to solve problems involving constant and variable forces.
\item Use the concept of momentum and impulse to analyse collisions and impacts.
\item State Newton's third law and correctly draw action–reaction pairs in free‑body diagrams.
\item Construct and interpret free‑body diagrams for particles and rigid bodies in equilibrium and motion.
\item Solve real‑life problems (e.g., vehicle braking, rocket thrust, tension in ropes) using Newton's laws.
\end{itemize}\end{multicols}}

\begin{sommaire}
\begin{multicols}{2}
\begin{enumerate}
\item 1. Inertia and the First Law
\item 2. Force, Acceleration and the Second Law
\item 3. Action–Reaction and the Third Law
\item 4. Applications, Limitations and Extensions
\item Integration activity
\item Methods and worked examples
\item Exercises
\item Solutions to the exercises
\item Summary sheet
\end{enumerate}
\end{multicols}
\end{sommaire}

\begin{objectifs}
\begin{itemize}
\item State and explain Newton's first law (law of inertia) and identify inertial reference frames.
\item Apply Newton's second law (F = ma) to solve problems involving constant and variable forces.
\item Use the concept of momentum and impulse to analyse collisions and impacts.
\item State Newton's third law and correctly draw action–reaction pairs in free‑body diagrams.
\item Construct and interpret free‑body diagrams for particles and rigid bodies in equilibrium and motion.
\item Solve real‑life problems (e.g., vehicle braking, rocket thrust, tension in ropes) using Newton's laws.
\item Discuss the limitations of Newtonian mechanics (high speeds, quantum scale) and outline its extensions.
\item Communicate mechanical reasoning clearly with appropriate mathematical notation and diagrams.
\end{itemize}
\end{objectifs}

\begin{prerequis}
\begin{itemize}
\item Vectors: addition, resolution, and scalar product (Lower Sixth).
\item Kinematics: displacement, velocity, acceleration, and equations of uniform acceleration.
\item Concept of force as a push or pull and its measurement in newtons.
\item Basic algebra and calculus (differentiation of simple functions).
\item Understanding of inertial and non‑inertial frames from introductory physics.
\end{itemize}
\end{prerequis}

\newpage

\coursec{Inertia and the First Law}

Imagine a loaded taxi speeding down the steep hills of Buea; when the driver brakes suddenly, the passengers lurch forward. Why do they keep moving? This everyday scene hides the fundamental principles that govern all motion, from a football kick in Yaoundé to the launch of a satellite. The taxi slows because the brakes act on \emph{it}; but nothing acts directly on the passengers, so they continue with the velocity they had. This simple observation, which every Cameroonian has felt in a \emph{bendskin} or a coaster bus, took humanity more than two thousand years to understand properly. In this section we trace that intellectual journey, define \cle{inertia} precisely, state \cle{Newton's First Law}, and learn to recognise the frames of reference in which it holds.

\courssub{Historical background: Aristotle, Galileo, Newton}

\textbf{Aristotle's view (4th century BC).} Aristotle taught that the ``natural state'' of a body is rest, and that a force is needed to \emph{maintain} any motion. Push a chair and it moves; stop pushing and it stops. This agrees with naive everyday experience, and for nearly twenty centuries it was accepted as truth.

\textbf{Galileo's breakthrough (early 17th century).} Galileo Galilei asked a subtler question: \emph{why} does the chair stop? His answer: because of friction, not because motion ``naturally dies''. He imagined (and performed) experiments with balls rolling down one inclined plane and up another. The ball always rises to almost the same height it started from; and as the second plane is made less and less steep, the ball travels farther and farther to regain that height. In the limit of a horizontal, perfectly smooth plane, the ball would never regain its height — it would roll on forever at constant speed. Galileo concluded:

\begin{quote}
In the absence of friction, a body in motion on a horizontal surface continues moving uniformly, with no force required to sustain the motion.
\end{quote}

Force is therefore needed to \emph{change} motion, not to maintain it. This was a revolution in thinking: the ``natural state'' includes uniform motion, not only rest.

\textbf{Newton's synthesis (1687).} Isaac Newton, in his \emph{Philosophi\ae\ Naturalis Principia Mathematica} (1687), elevated Galileo's insight to the status of a fundamental law — his First Law of Motion — and built upon it the whole edifice of classical mechanics. Rest and uniform rectilinear motion are dynamically equivalent: neither requires a cause; only \emph{changes} of velocity do.

\begin{savaistu}[A thought experiment that changed the world]
Galileo's inclined-plane argument is one of the first great ``thought experiments'' in science: he could not eliminate friction in his laboratory, so he eliminated it \emph{in his mind}, and reasoned about the limiting case. This method — idealise, then reason rigorously — remains at the heart of theoretical physics. Newton later wrote that he had ``stood on the shoulders of giants''; Galileo was the first of them.
\end{savaistu}

\courssub{Inertia and mass}

The taxi passengers lurch forward because their bodies ``resist'' the change of velocity imposed by the braking vehicle. This resistance to any change in the state of motion is called \cle{inertia}.

\begin{definition}[Inertia]
\cle{Inertia} is the property of every material body by which it resists any change in its velocity — whether a change of speed, of direction, or the start of motion from rest.
\end{definition}

Experience shows that some bodies have more inertia than others. It is far easier to stop a rolling football than a rolling water drum moving at the same speed; easier to push-start a bicycle than a lorry. The physical quantity that measures inertia is the \cle{mass}.

\begin{definition}[Mass as a measure of inertia]
The \cle{mass} $m$ of a body is the quantitative measure of its inertia: the greater the mass, the more the body resists changes in its velocity. Mass is a scalar quantity, measured in kilograms (\SI{}{kg}) in the SI system.
\end{definition}

\begin{attention}[Mass is not weight]
Mass (inertia) must not be confused with \emph{weight}, which is the gravitational force exerted on the body. A bag of cement of mass \SI{50}{kg} has the same inertia on Earth and on the Moon, but its weight is about six times smaller on the Moon (since $g_{\text{Moon}} \approx \SI{1.6}{m.s^{-2}}$ against $g_{\text{Earth}} \approx \SI{9.81}{m.s^{-2}}$). Inertia is intrinsic; weight depends on the local gravity. In the GCE examination, never write ``the mass of the bag is \SI{50}{N}'' — newtons measure force, not mass.
\end{attention}

\courssub{Statement of Newton's First Law}

\begin{propriete}[Newton's First Law (Law of Inertia)]
A body remains at rest, or continues to move in a straight line with constant velocity, unless it is compelled to change that state by a net external force acting upon it. Equivalently:
\[
\vec{F}_{\text{net}} = \sum \vec{F} = \vec{0} \quad \Longleftrightarrow \quad \vec{v} = \text{constant (possibly } \vec{0}\text{).}
\]
\end{propriete}

Several remarks are essential for a rigorous understanding:

\begin{itemize}
\item The law concerns the \emph{net} (resultant) force. A body may be acted upon by several forces; if they balance, the velocity does not change.
\item ``Constant velocity'' means constant \emph{vector} velocity: constant speed \emph{and} constant direction. A car going round a roundabout at steady speed is \emph{not} covered by the First Law — its direction changes, so a net force must act on it (we shall see in the chapter on circular motion that this force points towards the centre).
\item Rest is simply the special case $\vec{v}=\vec{0}$. The First Law thus unifies ``bodies at rest stay at rest'' and ``bodies in uniform motion stay in uniform motion''.
\item The law is an \emph{equivalence}: zero net force implies constant velocity, and conversely, observed constant velocity implies zero net force. Both directions are used in problems.
\end{itemize}

\begin{aretenir}[What the First Law really says]
\begin{itemize}
\item No net force $\Rightarrow$ no acceleration: $\vec{a}=\vec{0}$.
\item Force is the cause of \emph{change} of velocity, not of velocity itself.
\item A body in uniform rectilinear motion is dynamically indistinguishable from a body at rest.
\end{itemize}
\end{aretenir}

\courssub{Inertial reference frames}

Here lies a subtlety that the GCE examiners appreciate. Is the First Law true in \emph{every} frame of reference? Consider again the braking taxi. Relative to the \emph{road}, the passengers simply continue at constant velocity — the First Law holds. But relative to the \emph{taxi} (a decelerating frame), the passengers suddenly accelerate forward \emph{with no visible force pushing them}. The First Law appears violated. The taxi, while braking, is not a suitable frame for Newton's mechanics.

\begin{definition}[Inertial frame of reference]
An \cle{inertial frame} (or Galilean frame) is a frame of reference in which a body subject to no net force moves with constant velocity (in particular, remains at rest if initially at rest). Newton's laws are valid only in inertial frames.
\end{definition}

The Earth is, strictly speaking, not inertial: it rotates on its axis and revolves around the Sun, so it is permanently accelerating. However, for the durations and distances of ordinary laboratory experiments (and of GCE problems!), these accelerations are so small that \textbf{a frame fixed to the ground may be treated as inertial} to an excellent approximation. Moreover, any frame moving at constant velocity relative to an inertial frame is itself inertial: a train travelling smoothly at constant speed on straight rails is just as good a frame as the platform.

\begin{methode}[Testing whether a frame is inertial]
\begin{enumerate}
\item \textbf{Free-body test.} Isolate a body from all interactions (or arrange balanced forces). Observe its motion relative to the frame. If it stays at rest or moves uniformly in a straight line, the frame passes the test.
\item \textbf{Foucault pendulum test.} Suspend a long pendulum and set it swinging. The plane of oscillation is fixed relative to inertial space. If the plane appears to rotate relative to your frame (as Foucault demonstrated publicly in Paris in 1851), your frame is rotating — hence non-inertial. On Earth the plane rotates slowly, revealing that the Earth frame is only approximately inertial.
\item \textbf{Plumb-line test.} Hang a plumb line in the frame. If it hangs vertically and steadily while no horizontal force acts on the bob, the frame has no horizontal acceleration; if it tilts or swings without cause, the frame is accelerating.
\end{enumerate}
\end{methode}

\courssub{Identifying forces on a body at rest or in uniform motion}

The First Law gives us a powerful diagnostic tool: \textbf{if a body is at rest or in uniform rectilinear motion (in an inertial frame), then the vector sum of the forces on it is zero.} We use this to identify and relate forces. The systematic procedure is: isolate the body, list \emph{only} the forces acting \emph{on} it, draw the free-body diagram, then write $\sum\vec{F}=\vec{0}$.

\textbf{Example 1: the book on the table.} A book lies motionless on a horizontal table. Two forces act on it:
\begin{itemize}
\item its \cle{weight} $\vec{P}=m\vec{g}$, directed vertically downward (the Earth's pull);
\item the \cle{normal reaction} $\vec{R}$ of the table, directed vertically upward, perpendicular to the contact surface (the table's support).
\end{itemize}
Since the book is at rest, the First Law gives $\vec{P}+\vec{R}=\vec{0}$, hence $R=P=mg$.

\begin{popfigure}
\begin{center}
\begin{tikzpicture}[scale=1.0, >=stealth]
% table
\fill[popblueL] (-3,0) rectangle (3,0.25);
\draw[popline] (-3,0) -- (3,0);
\draw[popline] (-2.6,0) -- (-2.6,-0.9);
\draw[popline] (2.6,0) -- (2.6,-0.9);
\node[popmuted, font=\small] at (0,-0.65) {table (ground, inertial frame)};
% book
\fill[poporange!25] (-1.1,0.25) rectangle (1.1,0.95);
\draw[poporange, thick] (-1.1,0.25) rectangle (1.1,0.95);
\node[popdark, font=\small] at (0,0.6) {book, mass $m$};
% forces
\draw[->, very thick, poppurple] (0,0.6) -- (0,2.0) node[above, font=\small] {$\vec{R}$ (normal reaction)};
\draw[->, very thick, popgreen] (0,0.6) -- (0,-0.8) node[below, font=\small] {$\vec{P}=m\vec{g}$ (weight)};
\end{tikzpicture}
\end{center}
\legende{Free-body diagram of a book at rest on a horizontal table: the two forces balance, so the net force is zero and the book remains at rest (First Law).}
\end{popfigure}

\textbf{Example 2: the hockey puck on ice.} A puck struck on very smooth ice slides in a straight line at nearly constant speed long after the stick has left it. Horizontally, almost no force acts (friction is tiny); vertically, weight and the ice's reaction cancel. The net force is nearly zero, so the motion is nearly uniform — a beautiful, direct illustration of the First Law, and of why Galileo insisted on eliminating friction in his reasoning.

\textbf{Example 3: the passenger in the braking taxi.} Before braking, the passenger moves with the taxi at velocity $\vec{v}_0$; the forces on him (weight, seat reaction, friction with the seat) are balanced. When the taxi brakes, the \emph{taxi} experiences a backward force, but the passenger does not (at first). By the First Law, he continues at $\vec{v}_0$ relative to the road — which looks, from inside the taxi, like a forward lurch. The seat belt then provides the real backward force that decelerates him with the vehicle.

\begin{popfigure}
\begin{center}
\begin{tikzpicture}[scale=1.0, >=stealth]
% road
\draw[popline, thick] (-4,0) -- (4.6,0);
\foreach \x in {-3.5,-2.5,...,4.0} \draw[popmuted] (\x,0) -- (\x+0.5,0);
% taxi body
\fill[popgold!30] (-2.2,0.35) rectangle (1.6,1.35);
\draw[popdark, thick] (-2.2,0.35) rectangle (1.6,1.35);
\draw[popdark, thick] (-1.4,1.35) -- (-0.9,1.95) -- (0.7,1.95) -- (1.2,1.35);
\fill[popblueL] (-0.85,1.42) rectangle (0.65,1.88);
\node[popdark, font=\small] at (-0.3,0.85) {taxi braking};
% wheels
\fill[popink] (-1.5,0.35) circle (0.32);
\fill[popink] (0.9,0.35) circle (0.32);
% braking force on taxi
\draw[->, very thick, poppink] (-0.3,0.15) -- (-2.6,0.15) node[below left, font=\small, poppink] {braking force on taxi};
% passenger
\fill[popblue] (0.1,1.15) circle (0.16);
\draw[popblue, thick] (0.1,0.99) -- (0.1,0.55);
\draw[popblue, thick] (0.1,0.85) -- (-0.2,0.65);
\draw[popblue, thick] (0.1,0.85) -- (0.4,0.65);
\draw[popblue, thick] (0.1,0.55) -- (-0.12,0.38);
\draw[popblue, thick] (0.1,0.55) -- (0.32,0.38);
% inertia continuation arrow
\draw[->, very thick, popgreen] (0.35,1.15) -- (2.4,1.15) node[right, font=\small, align=left] {passenger keeps\\ velocity $\vec{v}_0$};
\node[popmuted, font=\small] at (-3.1,1.7) {road frame (inertial)};
\end{tikzpicture}
\end{center}
\legende{When the taxi brakes, a backward force acts on the taxi but not directly on the passenger. By inertia, the passenger continues forward at the initial velocity $\vec{v}_0$ until the seat belt or seat exerts a retarding force.}
\end{popfigure}

\begin{attention}[``No force'' versus ``balanced forces'']
A classic trap. The First Law is satisfied in \emph{both} situations — a body deep in space with \textbf{no force at all}, and the book on the table with \textbf{balanced forces} — because in both cases the \emph{net} force is zero and the acceleration is zero. But physically the situations differ: only the truly force-free body can serve to \emph{define} an inertial frame. A frame in which a genuinely force-free body accelerates is non-inertial, even if other bodies with balanced forces happen to move uniformly in it. In examinations, always say ``the \emph{resultant} (net) force is zero'', never ``there is no force'', when forces are present but balanced.
\end{attention}

\begin{exoresolu}[Forces on a crate in a moving lorry]
A crate of mass $m = \SI{40}{kg}$ rests on the flat floor of a lorry travelling along a straight horizontal road at a constant speed of $\SI{60}{km.h^{-1}}$. The crate does not slide.
\begin{enumerate}
\item List the forces acting on the crate and give the value of the resultant force.
\item The driver then brakes. Explain, using the First Law, what happens to the crate if the floor is very smooth.
\end{enumerate}
\tcblower
\textbf{Solution.}
\begin{enumerate}
\item The lorry moves uniformly in a straight line, so it constitutes (to a very good approximation) an inertial frame, and the crate is at rest in it. The forces on the crate are: its weight $\vec{P}=m\vec{g}$ (vertically downward) and the floor's normal reaction $\vec{R}$ (vertically upward). Since the velocity is constant, the First Law gives
\[
\vec{P}+\vec{R}=\vec{0} \quad\Longrightarrow\quad R = P = mg = 40 \times 9.81 = \SI{392.4}{N},
\]
and the resultant force is $\vec{F}_{\text{net}}=\vec{0}$. Note carefully: \emph{no horizontal force is needed}, because uniform motion requires none. (If the floor were rough, a tiny static friction could exist, but it is not required and we take it as zero here.)
\item When the lorry brakes, the lorry slows but no horizontal force acts on the crate (smooth floor, negligible friction). By the First Law the crate keeps its velocity of $\SI{60}{km.h^{-1}}$ relative to the road: it \emph{slides forward} relative to the lorry, towards the cabin, until it meets an obstacle that exerts a retarding force on it. This is exactly why loads in lorries must be strapped down.
\end{enumerate}
\end{exoresolu}

\begin{exercice}[Thinking with the First Law]
\begin{enumerate}
\item A ball of mass $\SI{0.45}{kg}$ rests on the pitch at the Ahmadou Ahidjo stadium. Draw its free-body diagram and compute the reaction of the ground (take $g=\SI{9.81}{m.s^{-2}}$).
\item A parachutist falls at \emph{constant} vertical speed. Is the resultant force on him zero? Name the forces involved.
\item Passengers in a coaster bus are thrown to the \emph{right} when the bus turns. In which direction did the bus turn? Justify with the First Law.
\item A frame of reference is fixed in a lift. In which of the following cases is this frame (approximately) inertial: (a) the lift is at rest; (b) the lift rises at constant speed; (c) the lift accelerates upward? Justify each answer.
\end{enumerate}
\end{exercice}

\begin{corrige}[Exercise 1]
\begin{enumerate}
\item Two forces act on the ball: its weight $\vec{P}=m\vec{g}$ downward and the ground's normal reaction $\vec{R}$ upward. The ball is at rest, so by the First Law $\vec{P}+\vec{R}=\vec{0}$, hence
\[
R = P = mg = 0.45 \times 9.81 \approx \SI{4.41}{N}, \quad \text{directed vertically upward.}
\]
\item Yes. Constant velocity (uniform rectilinear motion) implies, by the First Law, that the net force is zero. The forces are his weight $\vec{P}$ downward and the air resistance (drag) $\vec{f}$ upward, with $f = P$: the two forces balance. (This constant falling speed is called the \emph{terminal velocity}.)
\item The passengers' bodies tend to continue in a straight line (inertia) while the bus turns beneath them. If they are thrown to the right, the bus turned to the \emph{left}: relative to the bus, their rectilinear motion carries them towards the outside of the bend.
\item (a) At rest: inertial — a frame at rest relative to the ground is inertial to the usual approximation. (b) Rising at constant speed: inertial — it moves with constant velocity relative to an inertial frame. (c) Accelerating upward: \emph{not} inertial — a body released inside would appear to accelerate downward relative to the lift with no corresponding real force, violating the First Law.
\end{enumerate}
\end{corrige}

\begin{aretenir}[Key points of this section]
\begin{itemize}
\item Aristotle: force maintains motion (wrong). Galileo: friction stops motion; without it, uniform motion persists. Newton: this becomes the First Law.
\item \cle{Inertia} is the resistance of a body to changes of velocity; \cle{mass} measures inertia (scalar, in \SI{}{kg}); it must not be confused with weight (a force, in \SI{}{N}).
\item \textbf{First Law:} $\sum \vec{F}=\vec{0} \iff \vec{v}=\text{constant}$. Force changes velocity; it does not maintain it.
\item Newton's laws hold only in \cle{inertial frames}; the ground frame is inertial to an excellent approximation for everyday motions, and any frame in uniform rectilinear motion relative to it is also inertial.
\item A body at rest or in uniform motion has \emph{balanced} forces (net force zero) — which is not the same as having no forces at all.
\end{itemize}
\end{aretenir}

\coursec{Force, Acceleration and the Second Law}

In the previous section we met Newton's First Law: a body keeps its state of rest or of uniform motion in a straight line unless a \emph{net external force} compels it to change. That law tells us \emph{when} motion changes; it does not tell us \emph{by how much}. If you push a wheelbarrow full of sand at the market in Bamenda and then push the same wheelbarrow when it is empty, you feel immediately that the same push produces very different changes of motion. Newton's Second Law makes this observation precise: it is the quantitative heart of dynamics, linking \cle{force}, \cle{mass} and \cle{acceleration} in one compact equation.

\courssub{From intuition to the law: momentum}

Experience suggests two facts. First, the harder you push (the greater the force), the faster the velocity changes. Second, the heavier the object (the greater the mass), the harder it is to change its velocity. Newton captured both ideas by introducing a new quantity.

\begin{definition}[Momentum]
The \cle{momentum} of a body of mass $m$ moving with velocity $\vec{v}$ is the vector quantity
\[
\vec{p} = m\vec{v}.
\]
Its SI unit is the kilogram metre per second, $\mathrm{kg\,m\,s^{-1}}$ (equivalently the newton second, $\mathrm{N\,s}$). Momentum points in the same direction as the velocity.
\end{definition}

Momentum measures the ``quantity of motion'' of a body. A $1200\ \mathrm{kg}$ car moving at $10\ \mathrm{m\,s^{-1}}$ has the same magnitude of momentum as a $60\ \mathrm{kg}$ athlete sprinting at $20\ \mathrm{m\,s^{-1}}$ — both are equally ``hard to stop'' in this precise sense.

\begin{propriete}[Newton's Second Law — momentum form]
The \cle{rate of change of momentum} of a body equals the \cle{net external force} acting on it, and takes place in the direction of that force:
\[
\sum \vec{F} = \frac{d\vec{p}}{dt}.
\]
\emph{(This statement, and its derivation of $\sum\vec{F}=m\vec{a}$ for constant mass, are examinable.)}
\end{propriete}

This is the form Newton actually gave, and it is the most general one: it applies even when the mass of the system changes.

\courssub{Constant mass: the familiar equation $\sum \vec{F} = m\vec{a}$}

For a body whose mass does not change, the derivation is immediate:
\[
\sum \vec{F} = \frac{d\vec{p}}{dt} = \frac{d}{dt}(m\vec{v}) = m\frac{d\vec{v}}{dt} = m\vec{a}.
\]

\begin{propriete}[Second Law for constant mass]
If the mass $m$ of the system is constant, then
\[
\sum \vec{F} = m\vec{a},
\]
where $\sum \vec{F}$ is the vector sum of \emph{all} external forces and $\vec{a}$ is the acceleration of the centre of mass. For variable mass systems one must use the general form $\sum \vec{F} = \dfrac{d(m\vec{v})}{dt}$.
\end{propriete}

Several consequences deserve emphasis:

\begin{itemize}
\item The equation is \textbf{vectorial}: it is really three scalar equations, one per axis. In two dimensions, $\sum F_x = ma_x$ and $\sum F_y = ma_y$.
\item The SI unit of force follows from the law: $1\ \mathrm{N} = 1\ \mathrm{kg\,m\,s^{-2}}$. A force of $1\ \mathrm{N}$ gives a mass of $1\ \mathrm{kg}$ an acceleration of $1\ \mathrm{m\,s^{-2}}$.
\item Weight is just the Second Law applied to gravity: $\vec{W}=m\vec{g}$, so $W = mg$ with $g \approx 9.81\ \mathrm{m\,s^{-2}}$.
\item If $\sum\vec{F}=\vec{0}$ then $\vec{a}=\vec{0}$: the First Law is recovered as a special case.
\end{itemize}

\begin{attention}[Vectors before numbers]
Do not mix scalar and vector quantities. $\sum\vec{F}=m\vec{a}$ is a \emph{vector} equation. \textbf{Always resolve forces into components along chosen axes before substituting numbers.} Writing ``$F = ma$'' with magnitudes of forces that are not parallel to the motion is the most common source of error at GCE. Also remember: $m\vec{a}$ is \emph{not} an extra force to be drawn on the diagram — it is the \emph{result} of the real forces.
\end{attention}

\courssub{Impulse and change of momentum}

Integrating the Second Law with respect to time gives a result of great practical value, especially for collisions and blows, where forces act briefly but intensely.

\begin{definition}[Impulse]
The \cle{impulse} of a force $\vec{F}$ acting between times $t_1$ and $t_2$ is
\[
\vec{J} = \int_{t_1}^{t_2} \vec{F}\,dt.
\]
Its unit is the newton second, $\mathrm{N\,s}$.
\end{definition}

\begin{propriete}[Impulse–momentum theorem]
The impulse of the net force equals the change of momentum of the body:
\[
\vec{J} = \int_{t_1}^{t_2} \sum\vec{F}\,dt = \Delta\vec{p} = m\vec{v}_2 - m\vec{v}_1 \quad (\text{constant mass}).
\]
\end{propriete}

The proof is one line: since $\sum\vec{F} = d\vec{p}/dt$, integrating both sides from $t_1$ to $t_2$ gives $\int \sum\vec{F}\,dt = \vec{p}(t_2)-\vec{p}(t_1)$. Geometrically, the impulse is the \emph{area under the force–time graph}.

\begin{popfigure}
\begin{center}
\begin{tikzpicture}
\begin{axis}[
  width=10cm, height=6cm,
  axis lines=left,
  xlabel={$t$ (s)}, ylabel={$F$ (N)},
  xmin=0, xmax=1.2, ymin=0, ymax=1.3,
  xtick={0.2,0.7}, xticklabels={$t_1$,$t_2$},
  ytick=\empty,
  every axis x label/.style={at={(ticklabel* cs:1)}, anchor=west},
]
\addplot[draw=none, fill=popblueL, domain=0.2:0.7] {0.9*sin(deg(3.14159*(x-0.2)/0.5))} \closedcycle;
\addplot[very thick, popblue, domain=0.2:0.7, samples=60] {0.9*sin(deg(3.14159*(x-0.2)/0.5))};
\addplot[very thick, popblue] coordinates {(0,0) (0.2,0)};
\addplot[very thick, popblue] coordinates {(0.7,0) (1.1,0)};
\node[popink] at (axis cs:0.45,0.45) {$J=\displaystyle\int_{t_1}^{t_2} F\,dt = \Delta p$};
\end{axis}
\end{tikzpicture}
\end{center}
\legende{The impulse of a force is the area under the $F$--$t$ graph; it equals the change of momentum.}
\end{popfigure}

This explains everyday experience: a fielder catching a fast cricket ball draws his hands back, lengthening the contact time $\Delta t$; since $J=\Delta p$ is fixed by the ball, a longer $\Delta t$ means a smaller average force on the hands.

\courssub{A systematic method for dynamics problems}

\begin{methode}[Solving a dynamics problem]
\begin{enumerate}
\item \textbf{Choose the system}: decide which body (or bodies) you analyse, and isolate it mentally from its surroundings.
\item \textbf{Draw the free-body diagram}: represent \emph{all} external forces on the system (weight, normal reaction, tension, friction, applied forces) — nothing else.
\item \textbf{Apply $\sum\vec{F}=m\vec{a}$ in components}: choose axes wisely (often parallel and perpendicular to the motion), resolve every force, and write one scalar equation per axis.
\item \textbf{Solve the kinematics}: once $\vec{a}$ is known, use the constant-acceleration formulae ($v=u+at$, $s=ut+\tfrac12 at^2$, $v^2=u^2+2as$) to answer the question.
\end{enumerate}
\end{methode}

\courssub{Worked model 1: block on an inclined plane}

Consider a block of mass $m$ on a rough plane inclined at angle $\theta$ to the horizontal, with coefficient of kinetic friction $\mu_k$. The forces are the weight $m\vec{g}$, the normal reaction $\vec{N}$, and friction $\vec{f}$ opposing the motion.

\begin{popfigure}
\begin{center}
\begin{tikzpicture}[scale=1.1, >=stealth]
% inclined plane
\draw[thick, popink] (0,0) -- (5,0);
\draw[thick, popink] (0,0) -- (4.6,2.3);
\draw[popink] (1.1,0) arc (0:26.5:1.1) node[midway, right, xshift=2pt] {$\theta$};
% block (rotated frame)
\begin{scope}[rotate around={26.5:(2.9,1.45)}]
\draw[fill=popblueL, draw=popblue, thick] (2.4,1.45) rectangle (3.4,2.15);
\node at (2.9,1.8) {$m$};
% forces from centre
\draw[->, very thick, poporange] (2.9,1.8) -- ++(0,-1.9) node[below] {$m\vec{g}$};
\draw[->, very thick, popgreen] (2.9,1.8) -- ++(0,1.3) node[above] {$\vec{N}$};
\draw[->, very thick, poppurple] (2.9,1.8) -- ++(1.5,0) node[right] {$\vec{f}$};
% components of weight
\draw[->, thick, poppink, dashed] (2.9,1.8) -- ++(0,-1.7) node[pos=0.6, left, xshift=-2pt] {\scriptsize $mg\cos\theta$};
\draw[->, thick, poppink, dashed] (2.9,1.8) -- ++(-0.85,0) node[pos=0.5, above] {\scriptsize $mg\sin\theta$};
\end{scope}
\end{tikzpicture}
\end{center}
\legende{Free-body diagram on an inclined plane: the weight is resolved into components $mg\sin\theta$ (down the slope) and $mg\cos\theta$ (into the plane).}
\end{popfigure}

Take axes $x$ down the slope and $y$ perpendicular to it. Resolving:
\begin{align*}
y:\quad & N - mg\cos\theta = 0 \quad\Rightarrow\quad N = mg\cos\theta,\\
x:\quad & mg\sin\theta - f = ma, \qquad f = \mu_k N = \mu_k mg\cos\theta.
\end{align*}
Hence
\[
a = g\left(\sin\theta - \mu_k\cos\theta\right).
\]
Note the cancellation of $m$: all blocks slide down a given rough plane with the same acceleration, a fact Galileo would have appreciated.

\courssub{Worked model 2: the Atwood machine}

Two masses $m_1 > m_2$ hang from a light inextensible string passing over a smooth pulley. The string tension $T$ is the same on both sides; both masses share the same magnitude of acceleration $a$ (the string is inextensible). Applying the Second Law to each mass separately, taking the direction of motion as positive:
\begin{align*}
m_1 g - T &= m_1 a,\\
T - m_2 g &= m_2 a.
\end{align*}
Adding eliminates $T$:
\[
a = \frac{(m_1 - m_2)}{m_1 + m_2}\,g, \qquad
T = \frac{2 m_1 m_2}{m_1 + m_2}\,g.
\]

\begin{attention}[One system or two?]
A frequent mistake is to write $m_1 g - m_2 g = m_1 a$ (forgetting that $m_2$ also accelerates). When two bodies move together, either treat each body separately with its own equation, or treat the whole system with total mass $m_1+m_2$ and \emph{external} driving force $(m_1-m_2)g$. Never mix the two viewpoints in a single equation.
\end{attention}

\courssub{Worked model 3: braking distance of a car}

A car of mass $m$ travelling at speed $u$ brakes with a constant total resistive force $F$ (braking force plus friction). The Second Law gives the deceleration $a = F/m$. Using $v^2 = u^2 + 2as$ with $v=0$:
\[
0 = u^2 - 2\frac{F}{m}s \quad\Rightarrow\quad s = \frac{m u^2}{2F}.
\]
The stopping distance grows with the \emph{square} of the speed: doubling your speed quadruples the braking distance. This is why speed limits near schools in Yaoundé or Douala are a matter of physics, not bureaucracy.

\begin{exoresolu}[Impulse of a football kick]
A football of mass $0.43\ \mathrm{kg}$, initially at rest, is kicked and leaves the foot at $24\ \mathrm{m\,s^{-1}}$. The foot is in contact with the ball for $0.012\ \mathrm{s}$. Find (a) the impulse delivered, (b) the average force exerted on the ball.
\tcblower
\textbf{Solution.}
(a) By the impulse--momentum theorem, with the ball initially at rest:
\[
J = \Delta p = m v - 0 = 0.43 \times 24 = 10.32\ \mathrm{N\,s}.
\]
(b) Assuming a constant average force over the contact time,
\[
F_{\text{avg}} = \frac{J}{\Delta t} = \frac{10.32}{0.012} = 860\ \mathrm{N}.
\]
That is roughly the weight of an $88\ \mathrm{kg}$ person, applied for a hundredth of a second — enough to explain why a poorly struck ball stings the foot.
\end{exoresolu}

\begin{exoresolu}[Block pulled on a rough horizontal table]
A crate of mass $25\ \mathrm{kg}$ is pulled along a horizontal floor by a horizontal force of $120\ \mathrm{N}$. The coefficient of kinetic friction is $\mu_k = 0.30$. Take $g = 9.81\ \mathrm{m\,s^{-2}}$. Find the acceleration of the crate and the distance covered in $4.0\ \mathrm{s}$ starting from rest.
\tcblower
\textbf{Solution.} Free-body diagram: weight $mg$ down, normal reaction $N$ up, applied force $F=120\ \mathrm{N}$, friction $f$ opposing motion.

Vertically: $N = mg = 25 \times 9.81 = 245.25\ \mathrm{N}$.

Friction: $f = \mu_k N = 0.30 \times 245.25 = 73.6\ \mathrm{N}$.

Horizontally (Second Law):
\[
F - f = ma \quad\Rightarrow\quad a = \frac{120 - 73.6}{25} = 1.86\ \mathrm{m\,s^{-2}}.
\]
Kinematics from rest:
\[
s = \tfrac12 a t^2 = \tfrac12 \times 1.86 \times (4.0)^2 = 14.9\ \mathrm{m}.
\]
\end{exoresolu}

\courssub{A glimpse of variable mass: rocket propulsion}

So far the mass of the system was constant. A rocket is different: it ejects burnt fuel backwards, so its mass decreases as it accelerates. The simple formula $\sum\vec{F}=m\vec{a}$ then fails, and one must return to the general statement
\[
\sum \vec{F} = \frac{d(m\vec{v})}{dt}.
\]
Applying this to a rocket ejecting gas at speed $v_e$ relative to the rocket leads (the full derivation is beyond GCE, but the result is instructive) to the \cle{rocket equation}:
\[
m\frac{dv}{dt} = v_e\left(-\frac{dm}{dt}\right) + F_{\text{ext}},
\]
where $-dm/dt > 0$ is the rate at which fuel mass is expelled. The thrust $v_e(-dm/dt)$ pushes the rocket forward even in the vacuum of space, where there is nothing to ``push against'' — a beautiful confirmation that momentum, not air, is what matters. For the GCE, you should simply know that the Second Law in momentum form handles such systems, while $F=ma$ does not.

\begin{aretenir}[Key points of this section]
\begin{itemize}
\item Momentum: $\vec{p}=m\vec{v}$; the Second Law states $\sum\vec{F} = d\vec{p}/dt$.
\item Constant mass: $\sum\vec{F}=m\vec{a}$, applied \emph{in components} after resolving all forces.
\item Impulse: $\vec{J}=\int\vec{F}\,dt=\Delta\vec{p}$, the area under the $F$--$t$ graph.
\item Method: choose the system, draw the free-body diagram, apply $\sum\vec{F}=m\vec{a}$ per axis, then use kinematics.
\item Classic models: inclined plane gives $a=g(\sin\theta-\mu_k\cos\theta)$; Atwood machine gives $a=\dfrac{(m_1-m_2)g}{m_1+m_2}$; braking distance scales as $u^2$.
\item Variable mass (rockets): use $\sum\vec{F}=d(m\vec{v})/dt$, never $m\vec{a}$ alone.
\end{itemize}
\end{aretenir}

\begin{exercice}[Practice]
\begin{enumerate}
\item A minibus of mass $1800\ \mathrm{kg}$ accelerates from rest to $15\ \mathrm{m\,s^{-1}}$ in $10\ \mathrm{s}$. Find the net force acting on it, and the impulse delivered by the engine thrust (assumed to be the only horizontal force).
\item A block of mass $4.0\ \mathrm{kg}$ rests on a plane inclined at $30^{\circ}$ to the horizontal. The coefficient of kinetic friction is $0.20$. Find its acceleration down the plane. (Take $g=9.81\ \mathrm{m\,s^{-2}}$.)
\item In an Atwood machine, $m_1 = 5.0\ \mathrm{kg}$ and $m_2 = 3.0\ \mathrm{kg}$. Find the acceleration and the tension in the string.
\item A car travelling at $20\ \mathrm{m\,s^{-1}}$ brakes with a constant deceleration of $5.0\ \mathrm{m\,s^{-2}}$. Find the braking distance, and comment on the effect of doubling the initial speed.
\end{enumerate}
\end{exercice}

\begin{corrige}[Exercise 1]
\begin{enumerate}
\item $a = \dfrac{15-0}{10} = 1.5\ \mathrm{m\,s^{-2}}$; $F = ma = 1800 \times 1.5 = 2700\ \mathrm{N}$. Impulse $J = \Delta p = 1800 \times 15 = 27\,000\ \mathrm{N\,s}$ (check: $J = F\Delta t = 2700 \times 10$).
\item $a = g(\sin 30^{\circ} - 0.20\cos 30^{\circ}) = 9.81(0.5 - 0.20 \times 0.866) = 3.21\ \mathrm{m\,s^{-2}}$.
\item $a = \dfrac{(5-3)}{5+3}\times 9.81 = 2.45\ \mathrm{m\,s^{-2}}$; $T = \dfrac{2 \times 5 \times 3}{8}\times 9.81 = 36.8\ \mathrm{N}$.
\item $s = \dfrac{u^2}{2a} = \dfrac{400}{10} = 40\ \mathrm{m}$. Doubling $u$ to $40\ \mathrm{m\,s^{-1}}$ gives $s = 160\ \mathrm{m}$: four times farther.
\end{enumerate}
\end{corrige}

\begin{savaistu}[Why seatbelts are Newton's Second Law in action]
In a crash at $50\ \mathrm{km\,h^{-1}}$, an unbelted passenger stops in roughly $0.01\ \mathrm{s}$ on the dashboard, while a belted passenger stops over about $0.15\ \mathrm{s}$ as the belt stretches. The change of momentum $\Delta p$ is the same, but the average force $F = \Delta p/\Delta t$ is about fifteen times smaller with the belt. The Second Law, quietly, saves lives every day on the roads of Cameroon.
\end{savaistu}

\coursec{Action--Reaction and the Third Law}

In the previous section we learnt that a \emph{net} external force acting on a body produces an acceleration in proportion to that force: $\vec{F} = m\vec{a}$. But where do forces come from in the first place? A force is never a property of a single isolated body; it is always the symptom of an \emph{interaction} between two bodies. When you push a wheelbarrow, the wheelbarrow pushes back on your hands. When a pirogue's paddle pushes water backwards on the Wouri river, the water pushes the pirogue forwards. Newton's third law is the precise statement of this remarkable symmetry of nature, and it is the key to understanding propulsion, contact forces, and the behaviour of systems of connected bodies.

\courssub{Observation: forces always come in pairs}

\begin{experience}[A simple observation]
Sit on a chair with wheels (or stand on a skateboard) facing a wall, and push the wall away from you. You and the chair roll \emph{backwards}, away from the wall. Yet you pushed the wall \emph{forwards}. The only way to explain your motion is that the wall exerted a force on \emph{you}, directed away from it. Similarly, place two trolleys with spring-loaded bumpers in contact on a smooth bench and release them: \emph{both} trolleys move apart, each accelerating away from the other. Neither trolley can push the other without being pushed itself.
\end{experience}

The lesson of these observations is universal: whenever body $A$ exerts a force on body $B$, body $B$ simultaneously exerts a force on body $A$. There is no such thing as a one-sided force.

\courssub{Statement of the Third Law}

\begin{propriete}[Newton's Third Law of Motion]
If body $A$ exerts a force $\vec{F}_{AB}$ on body $B$, then body $B$ exerts a force $\vec{F}_{BA}$ on body $A$ such that
\[
\vec{F}_{BA} = -\vec{F}_{AB}.
\]
The two forces are equal in magnitude, opposite in direction, and act along the same line. This law holds for all types of interaction (contact, gravitational, electric, magnetic) and whether the bodies are at rest or accelerating. (The statement itself is examinable; no formal proof is required, but you must be able to \emph{apply} it rigorously.)
\end{propriete}

\begin{definition}[Action--reaction pair]
An \cle{action--reaction pair} is a pair of forces that:
\begin{enumerate}
\item are equal in magnitude;
\item are opposite in direction;
\item act on \emph{different} bodies; and
\item arise from the \emph{same} interaction, so they are of the same physical type (both gravitational, both contact, both frictional, etc.).
\end{enumerate}
If one force of the pair is called the \emph{action}, the other is the \emph{reaction}; the choice of which is which is purely conventional.
\end{definition}

\begin{attention}[The four characteristics are all essential]
Two forces of equal magnitude and opposite direction are \emph{not} necessarily an action--reaction pair. The decisive tests are: \textbf{(i)} do they act on different bodies? \textbf{(ii)} do they come from the same interaction? If either test fails, the forces are not a third-law pair.
\end{attention}

\courssub{Identifying action--reaction pairs}

\begin{methode}[Identifying action--reaction pairs]
\begin{enumerate}
\item Choose one force acting in the situation, and write it with two subscripts: $\vec{F}_{\text{exerting},\text{receiving}}$. For example $\vec{F}_{\text{Earth,book}}$ is the gravitational pull of the Earth on the book (the book's weight).
\item Identify the two bodies involved in that single interaction.
\item The reaction is the force exchanged in the same interaction, with the subscripts reversed: $\vec{F}_{\text{book,Earth}}$, the gravitational pull of the book on the Earth.
\item Check: the two forces act on different bodies, have the same magnitude, opposite directions, and the same physical nature.
\end{enumerate}
\end{methode}

\begin{exemplebox}[Everyday action--reaction pairs]
\begin{itemize}
\item \textbf{Book on a table.} The table pushes the book upwards (normal contact force on the book); the book pushes the table downwards (normal contact force on the table). Same interaction, different bodies.
\item \textbf{Your weight.} The Earth pulls you downwards with force $mg$; you pull the Earth \emph{upwards} with a force of magnitude $mg$. Because the Earth's mass is enormous, its acceleration is utterly negligible.
\item \textbf{Walking.} Your foot pushes the ground backwards and downwards; the ground pushes your foot forwards and upwards. It is the ground's forward push on you that accelerates you along the road.
\end{itemize}
\end{exemplebox}

\begin{attention}[Weight and normal reaction are NOT an action--reaction pair]
This is the most frequent error at the GCE. Consider a book resting on a table. Its weight $\vec{W}$ (Earth pulls book, downwards) and the normal reaction $\vec{R}$ of the table (table pushes book, upwards) are equal and opposite \emph{when the book is in equilibrium} --- but they both act on the \emph{same body} (the book) and they arise from \emph{different} interactions (gravitational for $\vec{W}$, contact for $\vec{R}$). Proof that they are not a pair: if you suddenly lift the table, $\vec{R}$ changes instantly while $\vec{W}$ does not; and in a lift accelerating upwards, $R > mg$, so they are not even equal. The true reaction to $\vec{W}$ is the book's gravitational pull on the Earth; the true reaction to $\vec{R}$ is the book's downward push on the table.
\end{attention}

\courssub{Free-body diagrams for interacting bodies}

To apply the second law correctly to each body of a system, we draw a separate \cle{free-body diagram} for each body, showing \emph{only} the forces acting \emph{on} that body. The third law then tells us how a force appearing in one diagram is related to a force in the other.

\begin{popfigure}
\begin{tikzpicture}[scale=1.0, >=stealth]
% Ground
\draw[thick, popmuted] (-0.5,0) -- (11.5,0);
\foreach \x in {0,0.6,...,11.4}{\draw[popmuted] (\x,0) -- (\x-0.25,-0.2);}
% Horse (stylised)
\draw[fill=popblueL, draw=popdark, thick] (1.2,0.55) ellipse (0.9 and 0.5);
\draw[fill=popblueL, draw=popdark, thick] (2.25,1.0) circle (0.32);
\draw[popdark, thick] (0.6,0.35) -- (0.55,0.02);
\draw[popdark, thick] (1.0,0.3) -- (0.95,0.02);
\draw[popdark, thick] (1.6,0.3) -- (1.65,0.02);
\draw[popdark, thick] (1.9,0.35) -- (2.0,0.02);
\node[popdark] at (1.2,0.55) {\small Horse};
% Rope
\draw[thick, poporange] (2.1,0.6) -- (4.6,0.6);
\node[poporange, above] at (3.35,0.62) {\small rope};
% Cart
\draw[fill=popgreenL, draw=popdark, thick] (4.6,0.45) rectangle (6.6,1.35);
\node[popdark] at (5.6,0.9) {\small Cart};
\draw[fill=white, draw=popdark, thick] (5.0,0.25) circle (0.25);
\draw[fill=white, draw=popdark, thick] (6.2,0.25) circle (0.25);
% Forces on horse
\draw[->, very thick, popblue] (1.2,0.55) -- (2.9,0.55) node[midway, above left=2pt, xshift=-6pt]{\small $\vec{F}_{\text{grd,horse}}$};
\draw[->, very thick, poppink] (1.2,0.55) -- (-0.1,0.55) node[midway, below left=1pt]{\small $\vec{T}_{\text{cart,horse}}$};
% Forces on cart
\draw[->, very thick, popblue] (5.6,1.35) -- (6.9,1.35) node[midway, above]{\small $\vec{T}_{\text{horse,cart}}$};
\draw[->, very thick, poppink] (5.6,1.35) -- (4.3,1.35) node[midway, above]{\small $\vec{f}_{\text{grd,cart}}$};
% Reaction of horse on ground
\draw[->, very thick, poppurple] (1.0,0.02) -- (-0.3,0.02) node[midway, below]{\small $\vec{F}_{\text{horse,grd}}$};
\end{tikzpicture}
\legende{Horse and cart. The horse pushes the ground backwards ($\vec{F}_{\text{horse,grd}}$); the ground pushes the horse forwards ($\vec{F}_{\text{grd,horse}}$): a third-law pair. The rope tensions $\vec{T}_{\text{cart,horse}}$ and $\vec{T}_{\text{horse,cart}}$ form another pair.}
\end{popfigure}

\begin{exemplebox}[The horse--cart paradox]
A student objects: ``If the cart pulls the horse backwards with a force equal and opposite to the force with which the horse pulls the cart forwards, the two forces cancel and the system can never move!'' Resolve the paradox.
\par\smallskip
\emph{Resolution.} The two forces of an action--reaction pair act on \emph{different} bodies, so they can never cancel each other. $\vec{T}_{\text{horse,cart}}$ acts on the \emph{cart} alone; $\vec{T}_{\text{cart,horse}}$ acts on the \emph{horse} alone. Whether the horse accelerates forwards depends on the forces \emph{on the horse}: the ground pushes it forwards with $\vec{F}_{\text{grd,horse}}$ while the rope pulls it backwards with $\vec{T}_{\text{cart,horse}}$. The horse--cart system accelerates because the ground exerts a net forward external force on it. The horse moves forward \emph{because it pushes the ground backwards}, not merely because it pulls the cart.
\end{exemplebox}

\courssub{Propulsion: the rocket and the swimmer}

Propulsion is the most spectacular illustration of the third law. A rocket does \emph{not} ``push against the air''; it works perfectly in the vacuum of space. The engine expels hot gas backwards at high speed: the rocket exerts a backward force on the gas, and the gas exerts an equal forward force --- the \cle{thrust} --- on the rocket.

\begin{popfigure}
\begin{tikzpicture}[scale=1.0, >=stealth]
% Exhaust plume
\fill[poporange!40] (0.2,0.5) -- (-2.6,1.3) -- (-2.6,-0.3) -- cycle;
\draw[->, very thick, poporange] (0.2,0.5) -- (-2.3,0.5) node[midway, below]{\small force on gas $\vec{F}_{\text{rock,gas}}$};
% Rocket body
\draw[fill=popblueL, draw=popdark, thick] (0.2,0.1) rectangle (2.6,0.9);
\draw[fill=popblueL, draw=popdark, thick] (2.6,0.1) -- (3.4,0.5) -- (2.6,0.9) -- cycle;
\draw[fill=popblueL, draw=popdark, thick] (0.2,0.1) -- (-0.3,-0.35) -- (0.6,0.1) -- cycle;
\draw[fill=popblueL, draw=popdark, thick] (0.2,0.9) -- (-0.3,1.35) -- (0.6,0.9) -- cycle;
\node[popdark] at (1.4,0.5) {\small Rocket};
% Thrust
\draw[->, very thick, popblue] (1.4,0.5) -- (4.6,0.5) node[midway, above]{\small thrust $\vec{F}_{\text{gas,rock}}$};
% Gas particles
\foreach \x/\y in {-1.0/0.9,-1.5/0.4,-2.0/0.8,-1.2/0.15,-1.8/0.2}{\fill[popmuted] (\x,\y) circle (0.05);}
\node[popmuted, align=center, text width=2.6cm] at (-1.6,-0.9) {\small exhaust gas expelled backwards};
\end{tikzpicture}
\legende{Free-body reasoning for a rocket. The rocket pushes the exhaust gas backwards; the gas pushes the rocket forwards. Thrust is the reaction force of the expelled gas on the rocket.}
\end{popfigure}

\begin{exemplebox}[The swimmer]
A swimmer pushes the water backwards and downwards with her hands and feet (action on the water). The water pushes her forwards and upwards with an equal and opposite force (reaction on the swimmer). It is this reaction that propels her along the pool. The same principle explains rowing, the flight of birds (wings push air down and back; air pushes the bird up and forward) and the recoil of a gun.
\end{exemplebox}

\courssub{Internal and external forces; motion of the centre of mass}

Consider a \emph{system} of several bodies (for example horse $+$ cart, or rocket $+$ exhaust gas). Forces acting within the system divide into two classes. This analysis addresses the syllabus requirement on extensions of Newton's laws (Lesson 181).

\begin{definition}[Internal and external forces]
A force is \cle{internal} to a system if it is exerted by one part of the system on another part of the same system. A force is \cle{external} if it is exerted on the system by a body outside it.
\end{definition}

By the third law, internal forces occur in equal and opposite pairs acting within the system, so their vector sum is zero. Applying the second law to the whole system of total mass $M$:
\[
\sum \vec{F}_{\text{ext}} + \underbrace{\sum \vec{F}_{\text{int}}}_{=\,\vec{0}} = M\vec{a}_{\text{cm}},
\qquad\text{hence}\qquad
\[\boxed{\;\sum \vec{F}_{\text{ext}} = M\vec{a}_{\text{cm}}\;}\]
\]
where $\vec{a}_{\text{cm}}$ is the acceleration of the \cle{centre of mass} of the system.

\begin{propriete}[Motion of the centre of mass]
The centre of mass of a system moves as if the whole mass of the system were concentrated there and all external forces were applied there. Internal forces cannot change the motion of the centre of mass. (This result follows directly from the second and third laws as shown above; the derivation is instructive and may be asked in a guided form.)
\end{propriete}

\begin{exemplebox}[Consequences]
\begin{itemize}
\item A man standing on a perfectly smooth (frictionless) frozen lake cannot start walking: there is no external horizontal force, so his centre of mass stays fixed. He can only move by throwing something (his shoe!) in the opposite direction.
\item For the horse--cart system, the internal tensions cancel; the only external horizontal force is the friction of the ground on the horse's hooves, and it is this force that accelerates the centre of mass of the system.
\item For the rocket $+$ gas system in deep space, the total momentum is conserved: as gas flies backwards, the rocket flies forwards, the centre of mass of the whole continuing undisturbed.
\end{itemize}
\end{exemplebox}

\courssub{Applications: tension and normal reaction}

\textbf{Tension in a rope.} When a light (massless) inextensible rope connects two bodies, the rope pulls each body towards itself with a force called the \cle{tension} $T$. Each body, in turn, pulls the rope with an equal opposite force (third law). For a \emph{light} rope the tension has the same magnitude throughout: if the tensions at the two ends differed, the net force on the massless rope would be non-zero, giving it an infinite acceleration, which is impossible.

\textbf{Normal reaction.} When two surfaces are in contact, each exerts on the other a contact force perpendicular to the surface: the \cle{normal reaction}. The table pushes the book up; the book pushes the table down with an equal and opposite force. These two normal reactions --- not the weight and the reaction --- form the third-law pair.

\begin{exoresolu}[Two connected bodies on a smooth table]
Two blocks, $A$ of mass $m_A = 2\ \mathrm{kg}$ and $B$ of mass $m_B = 3\ \mathrm{kg}$, rest in contact on a smooth horizontal table. A horizontal force $P = 10\ \mathrm{N}$ is applied to $A$, pushing it against $B$. Find (a) the acceleration of the blocks; (b) the contact force between them; (c) state the reaction to this contact force.
\tcblower
\textbf{(a) Acceleration.} Treat $A$ and $B$ as one system of mass $m_A + m_B = 5\ \mathrm{kg}$. The contact forces between the blocks are internal and cancel (third law). The only external horizontal force is $P$:
\[
a = \frac{P}{m_A + m_B} = \frac{10}{5} = 2\ \mathrm{m\,s^{-2}}.
\]
\textbf{(b) Contact force.} Isolate block $B$. The only horizontal force on $B$ is the push $F_{A,B}$ of $A$ on $B$. By the second law:
\[
F_{A,B} = m_B\, a = 3 \times 2 = 6\ \mathrm{N}.
\]
\emph{Check with block $A$:} the forces on $A$ are $P$ forwards and $F_{B,A}$ backwards, so $P - F_{B,A} = m_A a$, giving $F_{B,A} = 10 - 2\times 2 = 6\ \mathrm{N}$. The two contact forces are equal and opposite, as the third law demands.
\par\smallskip
\textbf{(c) Reaction.} The reaction to the force $F_{A,B}$ (block $A$ pushing block $B$ forwards) is the force $F_{B,A}$ of block $B$ pushing block $A$ backwards, of magnitude $6\ \mathrm{N}$. Note that $F_{A,B} \neq P$: the applied force and the contact force are \emph{not} an action--reaction pair, since they act on different bodies but arise from different interactions.
\end{exoresolu}

\begin{exoresolu}[Normal reaction in an accelerating lift]
A woman of mass $60\ \mathrm{kg}$ stands on the floor of a lift accelerating upwards at $1.5\ \mathrm{m\,s^{-2}}$. Find the normal reaction of the floor on her, and the force she exerts on the floor. Take $g = 9.81\ \mathrm{m\,s^{-2}}$.
\tcblower
The forces \emph{on the woman} are her weight $mg$ (downwards) and the normal reaction $R$ of the floor (upwards). Taking upwards as positive and applying the second law:
\[
R - mg = ma \quad\Longrightarrow\quad R = m(g+a) = 60\,(9.81 + 1.5) = 678.6\ \mathrm{N}.
\]
By the third law, the woman exerts on the floor a downward force of the same magnitude, $678.6\ \mathrm{N}$: this is the reaction to $R$, and it is what a scale placed between her feet and the floor would read. She feels ``heavier'' while the lift accelerates upwards. Note that $R \neq mg$ here: another proof that weight and normal reaction are not a third-law pair.
\end{exoresolu}

\courssub{Consolidation}

\begin{aretenir}[Key points]
\begin{itemize}
\item Newton's third law: $\vec{F}_{BA} = -\vec{F}_{AB}$; forces always come in pairs from a single interaction.
\item An action--reaction pair: equal magnitude, opposite direction, \textbf{different bodies}, same type of force.
\item The two members of a pair never cancel, because they act on different bodies.
\item Weight and normal reaction on the same body are \emph{not} an action--reaction pair.
\item Internal forces of a system cancel in pairs; the centre of mass obeys $\sum \vec{F}_{\text{ext}} = M\vec{a}_{\text{cm}}$.
\item Propulsion (rocket, swimmer, walking) is explained by the reaction of the expelled or pushed medium on the body.
\item In a light inextensible rope the tension is the same throughout; contact between surfaces produces equal and opposite normal reactions on each body.
\end{itemize}
\end{aretenir}

\begin{exercice}[Action--reaction pairs]
For each force below, state the body it acts on, and describe its third-law reaction: (a) the tension pulling a bucket upwards from a rope; (b) the friction force of the road on the tyre of a lorry accelerating along the Douala--Yaoundé highway; (c) the gravitational attraction of the Moon on the Earth.
\end{exercice}

\begin{exercice}[Connected bodies]
A tractor of mass $1200\ \mathrm{kg}$ pulls a trailer of mass $800\ \mathrm{kg}$ along a horizontal road with acceleration $0.5\ \mathrm{m\,s^{-2}}$. Resistances are negligible. Find (a) the tension in the coupling; (b) the forward driving force exerted by the ground on the tractor; (c) explain, using the third law, why the tension does not prevent the tractor from accelerating.
\end{exercice}

\begin{corrige}[Exercise 1]
(a) Acts on the bucket (upwards). Reaction: the bucket pulls the rope \emph{downwards} with an equal force. (b) Acts on the tyre (forwards). Reaction: the tyre pushes the road \emph{backwards} with an equal friction force. (c) Acts on the Earth (towards the Moon). Reaction: the Earth attracts the Moon with an equal force directed towards the Earth --- this is precisely the centripetal force keeping the Moon in orbit.
\end{corrige}

\begin{corrige}[Exercise 2]
(a) Isolate the trailer: $T = m_{\text{trailer}}\,a = 800 \times 0.5 = 400\ \mathrm{N}$. (b) For the whole system: $F = (1200+800)\times 0.5 = 1000\ \mathrm{N}$. (c) The trailer pulls the tractor backwards with $400\ \mathrm{N}$ (reaction to the tension on the trailer), but this acts \emph{on the tractor}, while the tension acts \emph{on the trailer}; they do not cancel. The tractor accelerates because the ground pushes it forwards with $1000\ \mathrm{N} > 400\ \mathrm{N}$.
\end{corrige}

\begin{savaistu}[Did you know?]
Newton stated his third law in the \emph{Principia} (1687) with the example of a horse pulling a rope: ``If a horse draws a stone tied to a rope, the horse will be equally drawn back towards the stone.'' More than three centuries later, the same law governs the launches of every satellite: a rocket climbing above the equator gains its thrust purely from the reaction of its own exhaust gases, needing nothing external to push against. Engineers speak of ``specific impulse'', but at heart it is all Newton's third law.
\end{savaistu}

\coursec{Applications, Limitations and Extensions}

Having established the three fundamental laws of motion, we now turn to their power in explaining the world around us, the boundaries beyond which they cease to apply, and the sophisticated mathematical frameworks that extend them. This section bridges the gap between textbook mechanics and the reality of engineering, modern physics, and examination technique.

\courssub{Real-Life Applications of Newton's Laws}

Newton's laws are not abstract curiosities; they are the operating system of the macroscopic world. Every vehicle you ride in, every bridge you cross, and every ball kicked in a football match obeys these principles.

\textbf{Vehicle Safety: Impulse and Crumple Zones}
Consider a car collision. Newton's Second Law in its original form, $\vec{F} = \frac{d\vec{p}}{dt}$, tells us that the force experienced during a crash is the rate of change of momentum. For a given change in momentum $\Delta \vec{p}$ (determined by the car's mass and change in velocity), the average force $\vec{F}_{\text{avg}}$ is inversely proportional to the duration $\Delta t$ of the impact:
\[
\vec{F}_{\text{avg}} = \frac{\Delta \vec{p}}{\Delta t}.
\]
The quantity $\int \vec{F}\,dt = \Delta \vec{p}$ is called the \cle{impulse} of the force. Engineers design \cle{crumple zones} at the front and rear of vehicles to deliberately collapse in a controlled manner during a crash. This increases $\Delta t$, thereby drastically reducing the peak force transmitted to the passenger compartment. Seat belts and airbags further increase the stopping time for the occupants relative to the car interior.

\begin{popfigure}
\begin{tikzpicture}[scale=0.9]
% Car before crash
\draw[thick, popblue] (0,0) rectangle (4,1.5);
\draw[thick, popblue] (1,1.5) -- (1.5,2.5) -- (3,2.5) -- (3,1.5);
\draw[fill=black] (0.7,0) circle (0.2);
\draw[fill=black] (3.3,0) circle (0.2);
\node[popblue] at (2, -0.5) {Before impact: $v = u$};

% Car during crash (crumple zone compressed)
\begin{scope}[xshift=6cm]
\draw[thick, poporange] (0,0) rectangle (3,1.5); % shortened chassis
\draw[thick, poporange, decorate, decoration={zigzag, segment length=4pt, amplitude=2pt}] (3,0) -- (4,0); % crumpled zone
\draw[thick, poporange] (3,0) -- (3,1.5);
\draw[thick, poporange] (1,1.5) -- (1.5,2.5) -- (2.5,2.5) -- (2.5,1.5);
\draw[fill=black] (0.7,0) circle (0.2);
\draw[fill=black] (2.3,0) circle (0.2);
\draw[->, thick, popdark] (2, 2.8) -- (2, 1.8) node[midway, right] {$F_{\text{wall}}$};
\node[poporange] at (1.5, -0.5) {During impact: Crumple zone absorbs energy};
\end{scope}

% Force-Time Graphs
\begin{scope}[yshift=-5.5cm, xshift=1cm]
\begin{axis}[
    width=12cm, height=4.5cm,
    axis lines=middle,
    xlabel={$t$}, ylabel={$F$},
    xmin=0, xmax=10, ymin=0, ymax=5,
    xtick={3,7}, xticklabels={$t_1$, $t_2$},
    ytick=\empty,
    clip=false,
    title style={yshift=-1.2cm},
    title={Force vs Time: Rigid Body vs Crumple Zone}
]
% Rigid body: short, high peak
\addplot[thick, poporange, samples=100, domain=2.5:3.5] {4*exp(-100*(x-3)^2)};
\node[poporange, anchor=south] at (axis cs:3,4.2) {Rigid (no crumple)};

% Crumple zone: longer, lower peak
\addplot[thick, popblue, samples=100, domain=5:9] {2*exp(-5*(x-7)^2)};
\node[popblue, anchor=south] at (axis cs:7,2.2) {Crumple zone};

% Area labels
\node[align=center] at (axis cs:3,1.5) {$\Delta p = \int F\,dt$};
\node[align=center] at (axis cs:7,0.8) {Same $\Delta p$ (equal area)};
\end{axis}
\end{scope}
\end{tikzpicture}
\legende{Left: Crumple zone increases impact duration. Right: Force--time graphs showing reduced peak force for the same impulse (area under curve).}
\end{popfigure}

\begin{exemplebox}[Seat Belt Force Calculation]
A car of mass $\SI{1200}{kg}$ travelling at $\SI{30}{m/s}$ hits a rigid wall and stops in $\SI{0.1}{s}$ without a crumple zone. With a crumple zone, the stopping time increases to $\SI{0.5}{s}$. Calculate the average force on the car in both cases.
\tcblower
\textbf{Without crumple zone:}
$\Delta p = m\Delta v = 1200 \times (0 - 30) = -36000\ \mathrm{kg\,m/s}$.
$F_{\text{avg}} = \frac{\Delta p}{\Delta t} = \frac{-36000}{0.1} = -\SI{360000}{N} = -\SI{360}{kN}$.

\textbf{With crumple zone ($\Delta t = 0.5\ \mathrm{s}$):}
$F_{\text{avg}} = \frac{-36000}{0.5} = -\SI{72000}{N} = -\SI{72}{kN}$.

The crumple zone reduces the average force by a factor of 5. This is the difference between a survivable accident and a fatal one.
\end{exemplebox}

\textbf{Sports: Ball Impact and the Coefficient of Restitution}
When a footballer kicks a ball, or a tennis player serves, Newton's Third Law acts at the contact surface: the force exerted by the foot on the ball equals the force exerted by the ball on the foot. The resulting acceleration of the ball ($a = F/m$) depends on the contact time and the average force. The \cle{coefficient of restitution} $e$ quantifies the "bounciness" of the collision:
\[
e = \frac{\text{relative speed of separation}}{\text{relative speed of approach}}, \qquad 0 \le e \le 1.
\]
For a perfectly elastic collision $e=1$ (kinetic energy conserved); for a perfectly inelastic collision $e=0$ (bodies stick together). This parameter governs the deformation forces during impact, a direct application of Newton's laws to deformable bodies.

\textbf{Engineering: Bridge Forces and Static Equilibrium}
For a bridge at rest, $\vec{a} = 0$, so Newton's First Law gives $\sum \vec{F} = 0$ and $\sum \vec{\tau} = 0$. Engineers resolve forces in trusses, cables, and arches using these equilibrium conditions. The design of the Wouri Bridge in Douala or the Bamenda Ring Road bridges relies on calculating tension, compression, and shear forces in every member to ensure $\sum F_x = 0$, $\sum F_y = 0$ for every joint (method of joints) and for every section (method of sections).

\courssub{Limitations of Newton's Laws}

Newton's laws are supremely successful for "everyday" objects: sizes $\gg$ atoms, speeds $\ll c$ (speed of light), and in inertial frames. They fail in three major regimes.

\begin{attention}[Domain of Validity]
Newton's Laws \textbf{do not apply} in the following situations:
\begin{enumerate}
    \item \textbf{Relativistic speeds:} When $v \gtrsim 0.1c$ ($c \approx \SI{3e8}{m/s}$), use Special/General Relativity.
    \item \textbf{Quantum scales:} When action $\sim h$ (Planck's constant, $h \approx \SI{6.63e-34}{J\,s}$), use Quantum Mechanics.
    \item \textbf{Non-inertial frames:} In accelerating/rotating frames, fictitious forces must be added to use $\sum \vec{F} = m\vec{a}$.
\end{enumerate}
\end{attention}

\textbf{1. High Speeds: Special Relativity}
At speeds approaching $c$, mass is not constant. The relativistic momentum is $\vec{p} = \gamma m_0 \vec{v}$, where $m_0$ is rest mass and $\gamma = \frac{1}{\sqrt{1 - v^2/c^2}}$ is the \cle{Lorentz factor}. Newton's Second Law becomes $\vec{F} = \frac{d\vec{p}}{dt} = \frac{d}{dt}(\gamma m_0 \vec{v})$. As $v \to c$, $\gamma \to \infty$, so an infinite force would be required to reach $c$. No massive particle can reach the speed of light.

\begin{propriete}[Relativistic Momentum and Energy]
For a particle of rest mass $m_0$ moving at velocity $\vec{v}$:
\[
\vec{p} = \gamma m_0 \vec{v}, \quad E = \gamma m_0 c^2, \quad \gamma = \frac{1}{\sqrt{1 - v^2/c^2}}.
\]
The Newtonian expressions $p = mv$ and $E_k = \frac{1}{2}mv^2$ are low-speed approximations ($v \ll c$) obtained by expanding $\gamma \approx 1 + \frac{v^2}{2c^2}$.
\end{propriete}

\textbf{2. Atomic Scales: Quantum Mechanics}
At the scale of atoms ($\sim \SI{e-10}{m}$), particles exhibit \cle{wave-particle duality}. The concept of a definite trajectory $(x(t), v(t))$ breaks down (Heisenberg Uncertainty Principle: $\Delta x \Delta p \ge \hbar/2$). Forces are replaced by potential energy functions $V(\vec{r})$ in the Schrödinger equation. Newton's deterministic laws are replaced by probabilistic wave mechanics.

\begin{savaistu}[Cameroon Context]
The \cle{Laser} (Light Amplification by Stimulated Emission of Radiation), used in telecommunications (fibre optics linking Yaoundé and Douala), medical surgery, and barcode scanners in supermarkets, is a pure quantum device. Its operation cannot be explained by Newton's laws; it requires quantum transitions between atomic energy levels.
\end{savaistu}

\textbf{3. Non-Inertial Frames: Fictitious Forces}
Newton's laws are formulated for \cle{inertial frames} (frames at rest or moving with constant velocity relative to the "fixed stars"). In a frame accelerating with $\vec{a}_0$ or rotating with angular velocity $\vec{\omega}$, the equation $\sum \vec{F}_{\text{real}} = m\vec{a}'$ (where $\vec{a}'$ is acceleration measured in the non-inertial frame) \textbf{fails}. To salvage the form $\sum \vec{F} = m\vec{a}'$, we must add \cle{fictitious forces} (also called inertial forces or pseudo-forces). This is the essence of d'Alembert's principle.

\begin{propriete}[Fictitious Forces in Non-Inertial Frames]
In a frame with translational acceleration $\vec{a}_0$ and angular velocity $\vec{\omega}$ (and angular acceleration $\dot{\vec{\omega}}$), a particle of mass $m$ with velocity $\vec{v}'$ and position $\vec{r}'$ in that frame experiences the following fictitious forces:
\begin{align*}
\vec{F}_{\text{trans}} &= -m\vec{a}_0 \quad \text{(translational)} \\
\vec{F}_{\text{centrifugal}} &= -m\vec{\omega} \times (\vec{\omega} \times \vec{r}') \quad \text{(centrifugal)} \\
\vec{F}_{\text{Coriolis}} &= -2m\vec{\omega} \times \vec{v}' \quad \text{(Coriolis)} \\
\vec{F}_{\text{Euler}} &= -m\dot{\vec{\omega}} \times \vec{r}' \quad \text{(Euler, if $\dot{\vec{\omega}} \neq 0$)}
\end{align*}
The equation of motion in the non-inertial frame becomes:
\[
\sum \vec{F}_{\text{real}} + \vec{F}_{\text{fictitious}} = m\vec{a}'.
\]
\end{propriete}

\begin{popfigure}
\begin{tikzpicture}[scale=1.1]
% Rotating frame: Earth or turntable
\draw[thick, popblue] (0,0) circle (3);
\draw[->, thick, popdark] (0,0) -- (0,3.5) node[above] {$\vec{\omega}$};

% Particle at position r' on x-axis
\fill[popred] (2.5,0) circle (0.08) node[below right] {$m$};
\draw[thick, poporange] (0,0) -- (2.5,0) node[midway, above] {$r'$};

% Velocity v' radial outward
\draw[->, thick, popdark] (2.5,0) -- (3.5,0) node[right] {$\vec{v}'$ (radial out)};

% Centrifugal force: radial outward
\draw[->, thick, popgreen] (2.5,0) -- (4.0,0) node[right] {$\vec{F}_{\text{cf}} = m\omega^2 r'$};

% Coriolis force: -2m ω × v' ; ω=ωk, v'=v'i => ω×v' = ωv' j (tangential +y). Fictitious = -2mωv' j (down, -y)
\draw[->, thick, poppurple] (2.5,0) -- (2.5,-1) node[below] {$\vec{F}_{\text{Cor}}$ (opposite rotation)};

% Inertial frame view (small inset)
\begin{scope}[xshift=8cm, yshift=1cm, scale=0.6]
\draw[thick, popblue] (0,0) circle (2);
\draw[->, thick, popdark] (0,0) -- (0,2.5) node[above] {$\vec{\omega}$};
\fill[popred] (2,0) circle (0.06);
\draw[->, thick, popdark] (2,0) -- (2,1) node[right] {$\vec{v}_{\text{inertial}}$ (tangential)};
\draw[->, thick, poporange] (2,0) -- (1,0) node[left] {$\vec{F}_{\text{centripetal}}$};
\draw[dashed, thin] (0,0) -- (2,0);
\node at (0,-2.5) {Inertial Frame: Real centripetal force only};
\end{scope}

\node at (0,-3.8) {Rotating Frame: Real force + Centrifugal + Coriolis = $m\vec{a}'$};
\end{tikzpicture}
\legende{In a rotating frame (e.g., Earth, merry-go-round), fictitious centrifugal and Coriolis forces appear. The Coriolis force deflects moving objects to the right in the Northern Hemisphere (e.g., trade winds, ocean currents affecting Cameroon's coast).}
\end{popfigure}

\begin{exemplebox}[Apparent Weight in a Lift]
A person of mass $m = \SI{70}{kg}$ stands on a bathroom scale (calibrated in Newtons) in a lift. Take $g = \SI{9.81}{m/s^2}$.
\begin{enumerate}
    \item Lift at rest or constant velocity: $a=0 \Rightarrow R = mg = \SI{686.7}{N}$.
    \item Lift accelerating upward at $\SI{2}{m/s^2}$: $R - mg = ma \Rightarrow R = m(g+a) = 70 \times 11.81 = \SI{826.7}{N}$.
    \item Lift accelerating downward at $\SI{2}{m/s^2}$: $mg - R = ma \Rightarrow R = m(g-a) = 70 \times 7.81 = \SI{546.7}{N}$.
    \item Lift in free fall ($a=g$): $R = 0$ (weightlessness).
\end{enumerate}
\textbf{Two equivalent views:}
\begin{itemize}
    \item \textit{Inertial frame (ground):} Net real force $R - mg = ma$.
    \item \textit{Non-inertial frame (lift):} Add fictitious force $-ma$. Equilibrium: $R - mg - ma = 0 \Rightarrow R = m(g+a)$.
\end{itemize}
\end{exemplebox}

\courssub{Extensions: Advanced Formulations (Beyond GCE Syllabus)}

While $\vec{F} = m\vec{a}$ is the workhorse of GCE Mechanics, physicists and engineers use more powerful formulations for complex systems (e.g., robots, satellites, quantum fields).

\begin{aretenir}[Extra for GCE: Lagrangian and Hamiltonian Mechanics]
These are \textbf{not required for the GCE Advanced Level examination} but are the standard language of university physics and engineering.
\begin{itemize}
    \item \textbf{Lagrangian Mechanics:} Defines the \cle{Lagrangian} $L = T - V$ (Kinetic minus Potential energy). The equations of motion are the \cle{Euler-Lagrange equations}:
    \[
    \frac{d}{dt}\left(\frac{\partial L}{\partial \dot{q}_i}\right) - \frac{\partial L}{\partial q_i} = 0,
    \]
    where $q_i$ are \cle{generalised coordinates} (any convenient parameters, e.g., angles for a double pendulum). This automatically handles constraints (strings, hinges) without calculating constraint forces (tensions, reactions).
    \item \textbf{Hamiltonian Mechanics:} Defines the \cle{Hamiltonian} $H = T + V$ (Total energy) in terms of coordinates $q_i$ and conjugate momenta $p_i = \partial L/\partial \dot{q}_i$. Hamilton's equations:
    \[
    \dot{q}_i = \frac{\partial H}{\partial p_i}, \quad \dot{p}_i = -\frac{\partial H}{\partial q_i}.
    \]
    This formulation is the direct gateway to Quantum Mechanics (Heisenberg/Schrödinger equations) and Statistical Mechanics.
\end{itemize}
\textbf{Why learn them?} They reveal symmetries and conservation laws (Noether's Theorem) and simplify problems with many constraints.
\end{aretenir}

\courssub{Revision Checklist and Exam Problem-Solving Strategy}

Success in GCE Mechanics questions on Newton's Laws requires a systematic, rigorous approach. Examiners award marks for clear diagrams, correct resolution of forces, and correct application of $\sum F = ma$.

\begin{methode}[GCE Exam Technique: Solving Newton's Laws Problems]
\begin{enumerate}
    \item \textbf{Read carefully:} Identify the system (particle, connected particles, rigid body). Note "smooth" (no friction) vs "rough" ($\mu$ given).
    \item \textbf{Draw a clear Free-Body Diagram (FBD) for \textbf{each} body:} Isolate the body. Show \textbf{all} forces: Weight ($mg$), Normal Reaction ($R$), Tension ($T$), Friction ($F_f \le \mu R$), Applied forces. \textbf{Do not} show $ma$ on the FBD; it is the result, not a force.
    \item \textbf{Choose coordinate axes:} Usually parallel/perpendicular to acceleration or slope. State your choice clearly (e.g., "Taking up the plane as positive").
    \item \textbf{Resolve forces:} Write components for each force along chosen axes.
    \item \textbf{Apply $\sum F = ma$ component-wise:} Write $\sum F_x = ma_x$, $\sum F_y = ma_y$. For connected particles, write an equation for each particle (and pulley if massive).
    \item \textbf{Solve the simultaneous equations:} Find unknowns ($a$, $T$, $R$, $\mu$, etc.).
    \item \textbf{Check consistency:} Does the direction of friction oppose relative motion? Is $F_f \le \mu R$ satisfied? Is the string taut ($T > 0$)?
    \item \textbf{State final answer with units:} e.g., "Acceleration = $\SI{2.45}{m/s^2}$ up the plane; Tension = $\SI{122.5}{N}$."
\end{enumerate}
\end{methode}

\begin{exoresolu}[Connected Particles on an Inclined Plane (GCE Style)]
A block $A$ of mass $\SI{4}{kg}$ rests on a rough plane inclined at $30^\circ$ to the horizontal. It is connected by a light inextensible string passing over a smooth pulley at the top of the plane to a hanging block $B$ of mass $\SI{3}{kg}$. The coefficient of friction between $A$ and the plane is $\mu = 0.25$. The system is released from rest. Find:
\begin{enumerate}
    \item the acceleration of the system,
    \item the tension in the string,
    \item the distance moved by $A$ in the first $\SI{2}{s}$.
\end{enumerate}
\tcblower
\textbf{Step 1: Free-Body Diagrams}
\begin{itemize}
    \item Block $A$ (on plane): Weight $4g$ vertically down. Normal reaction $R$ perpendicular to plane. Friction $F$ down the plane (opposing motion of $A$ up the plane). Tension $T$ up the plane.
    \item Block $B$ (hanging): Weight $3g$ down. Tension $T$ up.
\end{itemize}

\textbf{Step 2: Axes}
For $A$: $x$-axis up the plane, $y$-axis perpendicular to plane.
For $B$: $y$-axis vertically downward (direction of likely acceleration).

\textbf{Step 3: Equations} (using $g = \SI{9.81}{m/s^2}$)
\textbf{Block $A$ ($x$-direction):} $T - 4g\sin30^\circ - F = 4a$.
$y$-direction: $R - 4g\cos30^\circ = 0 \Rightarrow R = 4g\cos30^\circ = 4 \times 9.81 \times \frac{\sqrt{3}}{2} \approx \SI{33.95}{N}$.
$F = \mu R = 0.25 \times 33.95 \approx \SI{8.49}{N}$.
Equation (1): $T - 4(9.81)(0.5) - 8.49 = 4a \Rightarrow T - 19.62 - 8.49 = 4a \Rightarrow T - 28.11 = 4a$.

\textbf{Block $B$ ($y$-direction):} $3g - T = 3a \Rightarrow 29.43 - T = 3a$. \quad (Equation 2)

\textbf{Step 4: Solve}
Add (1) and (2): $(T - 28.11) + (29.43 - T) = 4a + 3a \Rightarrow 1.32 = 7a \Rightarrow a = \frac{1.32}{7} \approx \SI{0.1886}{m/s^2}$.
From (2): $T = 29.43 - 3(0.1886) \approx \SI{28.86}{N}$.

\textbf{Step 5: Kinematics for distance}
$u=0$, $t=2$, $a=0.1886$. $s = ut + \frac{1}{2}at^2 = 0 + 0.5 \times 0.1886 \times 4 = \SI{0.377}{m}$.

\textbf{Answers:} (i) $a \approx \SI{0.189}{m/s^2}$ (A moves up, B moves down). (ii) $T \approx \SI{28.9}{N}$. (iii) $s \approx \SI{0.377}{m}$.
\end{exoresolu}

\begin{attention}[Common GCE Traps]
\begin{itemize}
    \item \textbf{Sign errors:} Define positive direction \textbf{before} writing equations. Stick to it.
    \item \textbf{Friction direction:} Always opposes \textbf{relative motion} (or tendency). If unsure, assume a direction; if $F$ comes out negative, it acts the other way.
    \item \textbf{Normal reaction on a slope:} $R = mg\cos\theta$, \textbf{not} $mg$. Only $mg\cos\theta$ is perpendicular to the plane.
    \item \textbf{Connected particles:} Acceleration magnitude $a$ is the \textbf{same} for both (if string inextensible). Tension $T$ is the \textbf{same} throughout (if string light, pulley smooth).
    \item \textbf{Units:} Always give final answers in SI units ($\mathrm{m}$, $\mathrm{kg}$, $\mathrm{s}$, $\mathrm{N}$) unless asked otherwise.
    \item \textbf{$g$ value:} Use $g = \SI{9.81}{m/s^2}$ unless the question specifies $g = \SI{10}{m/s^2}$ or $g = \SI{9.8}{m/s^2}$.
\end{itemize}
\end{attention}

\begin{exercice}[Consolidation Exercise]
\begin{enumerate}
    \item A $\SI{1000}{kg}$ car accelerates from rest to $\SI{27}{m/s}$ ($ \approx \SI{100}{km/h}$) in $\SI{8}{s}$ on a horizontal road. The total resistance force is constant at $\SI{500}{N}$. Find the driving force of the engine.
    \item A particle of mass $\SI{0.5}{kg}$ slides down a rough plane inclined at $40^\circ$. The coefficient of friction is $0.3$. Find its acceleration.
    \item \textbf{Conceptual:} Explain why an astronaut in the International Space Station (orbiting at $\approx \SI{400}{km}$ altitude) feels weightless, even though $g \approx \SI{8.7}{m/s^2}$ there. (Hint: Non-inertial frame / free fall).
    \item \textbf{Relativity check:} Calculate the Lorentz factor $\gamma$ for an electron accelerated through a potential difference of $\SI{1}{MV}$ (rest energy $m_0c^2 = \SI{0.511}{MeV}$). Is a Newtonian treatment valid?
\end{enumerate}
\end{exercice}

\begin{corrige}[Exercise 4]
\begin{enumerate}
    \item $a = \frac{27}{8} = \SI{3.375}{m/s^2}$. $\sum F = ma \Rightarrow D - 500 = 1000 \times 3.375 \Rightarrow D = 3375 + 500 = \SI{3875}{N}$.
    \item $x$-axis down slope: $mg\sin40^\circ - \mu mg\cos40^\circ = ma$. $a = g(\sin40^\circ - 0.3\cos40^\circ) = 9.81(0.6428 - 0.3 \times 0.7660) = 9.81(0.6428 - 0.2298) = 9.81 \times 0.413 \approx \SI{4.05}{m/s^2}$.
    \item The ISS and astronaut are in continuous free fall (centripetal acceleration $a_c = g$). In the non-inertial frame of the ISS, the fictitious force $-mg$ exactly cancels the real weight $mg$. Contact force (apparent weight) = 0.
    \item Kinetic energy $K = eV = \SI{1}{MeV}$. Total energy $E = K + m_0c^2 = 1.511\ \mathrm{MeV}$. $\gamma = E/m_0c^2 = 1.511/0.511 \approx 2.96$. Since $\gamma \gg 1$, $v \approx 0.94c$. Newtonian mechanics fails; relativistic formulas required.
\end{enumerate}
\end{corrige}

\begin{aretenir}[Chapter Summary: Newton's Laws]
\begin{itemize}
    \item \textbf{1st Law:} Defines inertial frames. $\sum \vec{F} = 0 \Leftrightarrow \vec{a} = 0$. Inertia = mass.
    \item \textbf{2nd Law:} $\sum \vec{F} = \frac{d\vec{p}}{dt} = m\vec{a}$ (constant mass). Vector equation $\Rightarrow$ 3 scalar equations.
    \item \textbf{3rd Law:} Action-reaction pairs: same magnitude, opposite direction, \textbf{different bodies}, same nature.
    \item \textbf{Applications:} Impulse ($\Delta \vec{p} = \int \vec{F}\,dt$) for collisions/safety; Equilibrium ($\sum \vec{F}=0, \sum \vec{\tau}=0$) for structures.
    \item \textbf{Limitations:} Relativistic speeds ($v \to c$), Quantum scales ($x \sim \hbar/p$), Non-inertial frames (add fictitious forces).
    \item \textbf{Exam Key:} FBD $\rightarrow$ Axes $\rightarrow$ Resolve $\rightarrow$ $\sum F = ma$ $\rightarrow$ Solve $\rightarrow$ Check.
\end{itemize}
\end{aretenir}

\newpage

\coursec{Integration activity}

\coursec{Applications, Limitations and Extensions}

\courssub{Integrative Activity: The Mini-Roller Coaster}

\begin{aretenir}[Aims of the activity]
This integrative activity brings together the whole chapter: Newton's three laws, free-body diagrams, circular motion and energy reasoning. By the end of the activity you should be able to:
\begin{itemize}
\item apply the \cle{First Law} (inertia) to explain why a marble tends to continue in a straight line and why a centripetal force is needed to bend its path;
\item apply the \cle{Second Law}, $\sum \vec{F} = m\vec{a}$, to a body moving in a vertical circle, at the top and at the bottom of a loop;
\item identify \cle{action--reaction pairs} (Third Law) between the marble, the track and the Earth;
\item use \cle{conservation of mechanical energy} to determine the minimum release height for a complete loop;
\item draw clear \cle{free-body diagrams} and interpret the \emph{normal reaction} as the physical quantity that engineers monitor for safety;
\item discuss the \cle{limitations of the Newtonian model} (friction, air resistance, rolling, point-mass approximation) and how they affect a real design.
\end{itemize}
\end{aretenir}

\begin{aretenir}[Skills mobilised]
Modelling a real situation, choosing a system, drawing free-body diagrams, combining $F=ma$ with circular-motion kinematics ($a_c = v^2/R$), applying energy conservation, communicating results in a short technical report.
\end{aretenir}

\courssub{Situation}

The science club of your school in Buea has been invited to the regional science fair in Limbe. The club decides to present a \textbf{mini-roller coaster}: a track built from flexible plastic tubing and stiff cardboard, fixed to a plywood board. A small steel marble is released from rest at the top of a ramp of adjustable height $h$, runs down the ramp, along a short horizontal section, and then through a \textbf{vertical circular loop} of radius $R$ before leaving the track.

The club wants the marble to complete the loop \emph{without losing contact with the track}, and wants to justify the design to the jury using Newton's laws. They also want to discuss what would change if, instead of a marble, the loop carried a \emph{human passenger} in a seat with a seat belt, like the looping rides found in amusement parks.

\textbf{Data.} Take $g = \SI{9.81}{m.s^{-2}}$. The loop has radius $R = \SI{0.20}{m}$. The marble has mass $m = \SI{0.020}{kg}$ and is modelled as a \emph{point particle sliding without friction} (a first model; friction and rolling are discussed at the end). The marble is released from rest at a height $h$ measured vertically above the \emph{bottom} of the loop. The top of the loop is therefore at height $2R$ above the bottom.

\begin{popfigure}
\begin{tikzpicture}[scale=1.1, >=stealth]
% Ground line
\draw[thick, popmuted] (-0.5,0) -- (8,0);
% Ramp: starts at (0, h) where h=2.6cm approx 0.5m scaled? Let's scale 1cm = 0.1m? No, just schematic.
% Let's set R = 1cm for drawing? No, R=0.9cm in code. Let's make it cleaner.
% Loop center at (5, 1.5), radius 1.5. Bottom at (5,0). Top at (5,3).
% Ramp from (0,3.5) to (3.5,0).
\draw[very thick, popblue] (0,3.5) .. controls (1.5,3.0) and (2.5,1.0) .. (3.5,0);
% Horizontal entry
\draw[very thick, popblue] (3.5,0) -- (3.5,0); % just point
% Loop
\draw[very thick, popblue] (5,1.5) circle (1.5);
% Horizontal exit
\draw[very thick, popblue] (6.5,0) -- (8,0);
% Marble at start
\fill[poporange] (0.2,3.52) circle (0.1);
% Height h
\draw[<->, popdark] (7.5,0) -- (7.5,3.5) node[midway, right, text=popdark] {$h$};
\draw[dashed, gray] (0,3.5) -- (7.5,3.5);
% Height 2R
\draw[<->, popdark] (6.8,0) -- (6.8,3.0) node[midway, right, text=popdark] {$2R$};
\draw[dashed, gray] (5,3.0) -- (6.8,3.0);
% Loop center and R label
\fill[popdark] (5,1.5) circle (0.04);
\draw[<->, popdark] (5,1.5) -- (6.5,1.5) node[midway, above, text=popdark] {$R$};
% Labels
\node[text=popink, align=center] at (1.0,3.9) {Release from rest};
\node[text=popink, align=center] at (5,-0.8) {Vertical loop of radius $R$};
\end{tikzpicture}
\legende{The mini-roller coaster: release at height $h$, vertical loop of radius $R$.}
\end{popfigure}

The club's questions to you, the physics team, are the tasks below.

\courssub{Tasks}

\begin{enumerate}
\item \textbf{Energy analysis.} Using conservation of mechanical energy between the release point and the top of the loop, show that the speed $v$ of the marble at the top of the loop satisfies $v^2 = 2g(h - 2R)$.

\item \textbf{Free-body diagram at the top.} Draw the free-body diagram of the marble at the \emph{top} of the loop (weight and normal reaction of the track). Apply Newton's Second Law along the vertical (centripetal) direction to show that
\[ N + mg = \frac{mv^2}{R}, \]
where $N$ is the magnitude of the normal reaction of the track on the marble.

\item \textbf{Contact condition.} Explain, using Newton's First Law, why a force directed towards the centre is necessary for the marble to follow the circle. Deduce that the marble just keeps contact with the track when $N = 0$, and show that the minimum speed at the top is $v_{\min} = \sqrt{gR}$.

\item \textbf{Minimum height.} Combine the results of questions 1 and 3 to show that the minimum release height is $h_{\min} = \dfrac{5R}{2}$. Compute its numerical value for $R = \SI{0.20}{m}$.

\item \textbf{Bottom of the loop.} Draw the free-body diagram at the \emph{bottom} of the loop and show that the normal reaction there is $N_{\text{bottom}} = mg + \dfrac{mv_b^2}{R}$, where $v_b$ is the speed at the bottom. For $h = h_{\min}$, show that $N_{\text{bottom}} = 6mg$ and compute its value in newtons.

\item \textbf{Action--reaction pairs.} Identify all the Third-Law pairs acting when the marble is at the top of the loop (marble--track, marble--Earth). Explain why the weight of the marble and the normal reaction of the track on the marble do \emph{not} form an action--reaction pair.

\item \textbf{Human passenger.} A passenger of mass $M = \SI{60}{kg}$ goes through a loop of radius $R = \SI{8}{m}$ at a fairground ride, with speed $v = \SI{12}{m.s^{-1}}$ at the top. Compute the force exerted by the seat (or seat belt) on the passenger at the top. Comment on the direction of this force and on what the passenger feels (``weightlessness'' when $v = \sqrt{gR}$).

\item \textbf{Limitations of the model.} Discuss at least three reasons why the real marble in the cardboard track needs a release height \emph{greater} than $\frac{5R}{2}$ (friction, air resistance, rolling of the marble, deformation of the track). Explain what the club should do experimentally to find the true minimum height.

\item \textbf{Report.} Write a half-page summary for the jury, including your diagrams, the value of $h_{\min}$, and one sentence on the domain of validity of Newton's laws.
\end{enumerate}

\courssub{Model Answer}

\textbf{1. Energy analysis.} We choose the bottom of the loop as the reference level (zero) for gravitational potential energy.
\begin{itemize}
\item At the release point: the marble is at rest ($v=0$) at height $h$. Mechanical energy $E_i = mgh$.
\item At the top of the loop: height $2R$, speed $v$. Mechanical energy $E_f = \frac{1}{2}mv^2 + mg(2R)$.
\end{itemize}
Since the track is assumed frictionless and the marble slides without rolling, mechanical energy is conserved: $E_i = E_f$.
\[ mgh = \frac{1}{2}mv^2 + 2mgR \]
Dividing by $m$ and multiplying by 2:
\[ v^2 = 2g(h - 2R). \]

\textbf{2. Free-body diagram at the top.} At the top of the loop, the centre of the circular path is \emph{below} the marble. Two forces act on the marble, both directed vertically downwards (towards the centre):
\begin{itemize}
\item the weight $\vec{W} = m\vec{g}$ (action of the Earth on the marble);
\item the normal reaction $\vec{N}$ of the track on the marble (the track is above the marble, so it pushes downward).
\end{itemize}
The acceleration is centripetal, directed downwards, with magnitude $a_c = v^2/R$.

\begin{popfigure}
\begin{tikzpicture}[scale=1.0, >=stealth]
% Top of loop
\draw[thick, popblue] (0,0) circle (1.5);
% Marble at top (0, 1.5)
\fill[poporange] (0,1.5) circle (0.12);
% Weight
\draw[->, very thick, popdark] (0,1.5) -- (0,0.3) node[midway, right] {$m\vec{g}$};
% Normal reaction
\draw[->, very thick, poppurple] (0,1.5) -- (0,2.7) node[midway, right] {$\vec{N}$};
% Centre
\fill[popdark] (0,0) circle (0.03);
\draw[dashed, gray] (0,1.5) -- (0,0) node[midway, left] {$R$};
% Acceleration vector (optional, distinct style)
\draw[->, very thick, popgreen] (0,1.5) -- (0,-0.2) node[midway, left] {$\vec{a}_c$};
\node[text=popink, align=center] at (0,-2.2) {Top of the loop: FBD};
\end{tikzpicture}
\legende{Free-body diagram at the top of the loop. Both forces point toward the centre.}
\end{popfigure}

Applying Newton's Second Law ($\sum F = ma$) in the downward (centripetal) direction:
\[ N + mg = m\frac{v^2}{R}. \]

\textbf{3. Contact condition.} Newton's First Law (Law of Inertia) states that a body continues in a state of uniform rectilinear motion unless acted upon by a net external force. To make the marble follow a circular path, a net force directed towards the centre (centripetal force) is continuously required to change the direction of the velocity vector.

If the speed is too low, the required centripetal force $mv^2/R$ becomes smaller than the weight $mg$. Since the track can only \emph{push} (not pull), the normal reaction $N$ cannot become negative. The limiting case for maintaining contact is $N = 0$. At this limit, gravity alone provides the centripetal force:
\[ mg = \frac{m v_{\min}^2}{R} \quad\Longrightarrow\quad v_{\min} = \sqrt{gR}. \]

\textbf{4. Minimum height.} Substitute $v^2 = gR$ (from the contact condition) into the energy equation (Question 1):
\[ gR = 2g(h_{\min} - 2R) \]
\[ h_{\min} - 2R = \frac{R}{2} \quad\Longrightarrow\quad h_{\min} = \frac{5R}{2}. \]
For $R = \SI{0.20}{m}$:
\[ h_{\min} = \frac{5 \times 0.20}{2} = \SI{0.50}{m}. \]

\textbf{5. Bottom of the loop.} At the bottom, the centre of the circle is \emph{above} the marble. The centripetal acceleration points upward.
\begin{itemize}
\item Weight $m\vec{g}$ acts downward.
\item Normal reaction $\vec{N}_{\text{bottom}}$ acts upward (track pushes up on marble).
\end{itemize}

\begin{popfigure}
\begin{tikzpicture}[scale=1.0, >=stealth]
% Bottom of loop
\draw[thick, popblue] (0,0) circle (1.5);
% Marble at bottom (0, -1.5)
\fill[poporange] (0,-1.5) circle (0.12);
% Weight
\draw[->, very thick, popdark] (0,-1.5) -- (0,-2.7) node[midway, right] {$m\vec{g}$};
% Normal reaction
\draw[->, very thick, poppurple] (0,-1.5) -- (0,-0.3) node[midway, right] {$\vec{N}_{\text{bottom}}$};
% Centre
\fill[popdark] (0,0) circle (0.03);
\draw[dashed, gray] (0,-1.5) -- (0,0) node[midway, left] {$R$};
% Acceleration
\draw[->, very thick, popgreen] (0,-1.5) -- (0,0.2) node[midway, left] {$\vec{a}_c$};
\node[text=popink, align=center] at (0,-3.2) {Bottom of the loop: FBD};
\end{tikzpicture}
\legende{Free-body diagram at the bottom of the loop. $\vec{N}_{\text{bottom}}$ points toward the centre.}
\end{popfigure}

Applying Newton's Second Law in the upward (centripetal) direction:
\[ N_{\text{bottom}} - mg = \frac{m v_b^2}{R} \quad\Longrightarrow\quad N_{\text{bottom}} = mg + \frac{m v_b^2}{R}. \]
To find $v_b$ for $h = h_{\min} = 5R/2$, use energy conservation between release and bottom:
\[ mgh_{\min} = \frac{1}{2}m v_b^2 \quad\Longrightarrow\quad v_b^2 = 2gh_{\min} = 2g\left(\frac{5R}{2}\right) = 5gR. \]
Substitute into the expression for $N_{\text{bottom}}$:
\[ N_{\text{bottom}} = mg + \frac{m(5gR)}{R} = 6mg. \]
Numerical value:
\[ N_{\text{bottom}} = 6 \times \SI{0.020}{kg} \times \SI{9.81}{m.s^{-2}} = \SI{1.1772}{N} \approx \SI{1.18}{N}. \]
The track at the bottom must withstand a force six times the weight of the marble; this is the critical design point for structural strength.

\textbf{6. Action--reaction pairs.} When the marble is at the top of the loop, the following Third-Law pairs exist:
\begin{itemize}
\item \textbf{Marble--Track pair:} The track exerts a downward force $\vec{N}$ on the marble. The marble exerts an \emph{upward} force $-\vec{N}$ on the track. These forces are equal in magnitude, opposite in direction, and act on \emph{different bodies}.
\item \textbf{Marble--Earth pair:} The Earth exerts a downward gravitational force $m\vec{g}$ (weight) on the marble. The marble exerts an \emph{upward} gravitational force $-m\vec{g}$ on the Earth. These forces are equal in magnitude, opposite in direction, and act on \emph{different bodies}.
\end{itemize}
The weight $m\vec{g}$ and the normal reaction $\vec{N}$ \textbf{do not form an action--reaction pair} because:
\begin{enumerate}
\item They both act on the \emph{same body} (the marble).
\item They are not necessarily equal in magnitude (here $N = mv^2/R - mg$, which is zero at the minimum speed).
\item They are of different physical nature (gravitational vs. contact electromagnetic).
\end{enumerate}

\textbf{7. Human passenger.} At the top of the loop, the forces on the passenger (mass $M = \SI{60}{kg}$) are weight $M\vec{g}$ (down) and the force $\vec{F}$ from the seat/belt (down, towards centre). Speed $v = \SI{12}{m.s^{-1}}$, radius $R = \SI{8}{m}$.
Newton's Second Law (downward positive):
\[ F + Mg = \frac{M v^2}{R} \]
\[ F = M\left(\frac{v^2}{R} - g\right) = 60\left(\frac{144}{8} - 9.81\right) = 60 \times (18 - 9.81) = 60 \times 8.19 = \SI{491.4}{N} \approx \SI{491}{N}. \]
The seat (or seat belt) must push \emph{downward} on the passenger with a force of about \SI{491}{N}. By Newton's Third Law, the passenger pushes \emph{upward} on the seat with the same force; this is the "apparent weight" felt by the passenger.

If the speed were exactly $v = \sqrt{gR} = \sqrt{9.81 \times 8} \approx \SI{8.86}{m.s^{-1}}$, then $F = 0$. The passenger would feel \emph{weightless} (no contact force from the seat), with gravity alone providing the centripetal force. For speeds below this, a seat belt or harness (providing a downward pull, or an upward push from a shoulder harness) is essential to keep the passenger on the circular path.

\textbf{8. Limitations of the model.} The ideal prediction $h_{\min} = 5R/2 = \SI{0.50}{m}$ assumes a point mass sliding without friction. In reality, the marble needs a greater release height because:
\begin{itemize}
\item \textbf{Friction} between the marble and the tubing converts mechanical energy into thermal energy, reducing the speed at the top.
\item \textbf{Air resistance} (drag) opposes the motion, dissipating energy, especially at higher speeds on the descent.
\item \textbf{Rolling without slipping:} The marble rotates. Part of the potential energy is converted into rotational kinetic energy $\frac{1}{2}I\omega^2$. For a solid sphere ($I = \frac{2}{5}mr^2$), the total kinetic energy is $\frac{7}{10}mv^2$, so the translational speed $v$ at a given height is lower than for a sliding block. The theoretical minimum height for a rolling sphere is $h_{\min} = \frac{27}{10}R = 2.7R$ (compared to $2.5R$).
\item \textbf{Track deformation and vibrations:} The cardboard track flexes and vibrates, absorbing energy that is not recovered.
\end{itemize}
\textbf{Experimental determination:} The club should release the marble from a height clearly above the predicted minimum (e.g., $\SI{0.60}{m}$) and verify it completes the loop. Then lower the release height in small steps (e.g., $\SI{0.01}{m}$), performing several trials at each height (e.g., 5 trials) to account for variability. The true minimum height is the lowest height at which the marble successfully completes the loop in all (or most) trials. The difference between this experimental value and $\SI{0.50}{m}$ quantifies the total energy losses.

\textbf{9. Report (sample summary for the jury).}
\emph{``Our mini-roller coaster model applies Newton's laws to a marble of mass \SI{20}{g} in a vertical loop of radius \SI{0.20}{m}. Modelling the marble as a point particle sliding without friction, conservation of energy and Newton's second law yield a minimum release height $h_{\min} = 5R/2 = \SI{0.50}{m}$ for the marble to complete the loop without losing contact. At the top of the loop, the centripetal force is provided by the sum of the weight and the normal reaction from the track ($N + mg = mv^2/R$); contact is maintained as long as $N \ge 0$, giving $v_{\min} = \sqrt{gR}$. At the bottom, the normal reaction peaks at $6mg \approx \SI{1.18}{N}$, dictating the required track strength. Action--reaction pairs are identified between the marble--track and marble--Earth; weight and normal reaction are not a Third-Law pair as they act on the same body. Real-world effects (friction, air drag, rolling inertia, track deformation) increase the required height; we determined the true minimum experimentally by incremental trials. Newton's laws provide an excellent model here because speeds are non-relativistic and the system is macroscopic.''}

\begin{attention}[Common errors to avoid in the examination]
\begin{itemize}
\item \textbf{Inventing a ``centrifugal force'' on the free-body diagram.} In an inertial reference frame (ground frame), only real forces (weight, normal reaction, tension, friction) appear. Centrifugal force is a fictitious force that appears only in the rotating non-inertial frame of the marble.
\item \textbf{Setting $v=0$ at the top as the condition for completing the loop.} The marble must have a minimum speed $v_{\min} = \sqrt{gR}$ at the top to maintain contact with the track. Zero speed would mean the marble falls straight down.
\item \textbf{Pairing weight and normal reaction as action and reaction.} They act on the \emph{same body} (the marble). Action--reaction pairs \emph{always} act on two \emph{different} bodies.
\item \textbf{Forgetting the direction of the normal reaction.} At the top of an \emph{inside} loop, the track is above the marble, so it pushes \emph{down}. At the bottom, the track is below, so it pushes \emph{up}. Always draw the FBD first.
\item \textbf{Unit inconsistency.} Ensure $R$ is in metres, $g$ in $\mathrm{m\,s^{-2}}$, mass in $\mathrm{kg}$ to get force in newtons and energy in joules.
\end{itemize}
\end{attention}

\newpage

\coursec{Methods and worked examples}

\coursec{Newton's Laws of Motion}

\begin{savaistu}[Historical Context]
Sir Isaac Newton (1643–1727) published his three laws of motion in \emph{Philosophiæ Naturalis Principia Mathematica} (1687). These laws form the foundation of classical mechanics and remain the primary tool for analysing motion at everyday speeds and scales. In Cameroon, from the motion of vehicles on the Yaoundé–Douala highway to the launch of satellites, Newton's laws provide the quantitative framework engineers and physicists rely on.
\end{savaistu}

\courssub{Inertia and Newton's First Law (Law of Inertia)}

\begin{definition}[Inertia and Inertial Mass]
\emph{Inertia} is the inherent property of a body to resist any change in its state of rest or uniform motion in a straight line. The quantitative measure of inertia is the \emph{inertial mass} $m$ (SI unit: kilogram, kg). The greater the mass, the greater the force required to produce a given acceleration.
\end{definition}

\begin{definition}[Inertial Frame of Reference]
An \emph{inertial frame of reference} is a frame in which a body not subject to any net external force moves with constant velocity (which includes being at rest). Any frame moving at constant velocity relative to an inertial frame is also inertial. The Earth's surface is approximately inertial for most classroom mechanics; a lift accelerating upward is \emph{not} inertial.
\end{definition}

\begin{propriete}[Newton's First Law (Law of Inertia)]
If the vector sum of all forces acting on a particle is zero, the particle remains at rest or continues to move with constant velocity in a straight line.
\[
\sum \vec{F} = \vec{0} \quad\Longleftrightarrow\quad \vec{a} = \vec{0} \quad\Longleftrightarrow\quad \vec{v} = \text{constant}.
\]
This law defines what a force \emph{is}: an agent that changes the state of motion. It also identifies the inertial frames in which the Second Law holds.
\end{propriete}

\begin{methode}[Drawing a Free-Body Diagram and Testing Equilibrium]
\begin{enumerate}
    \item \textbf{Identify the body} of interest and isolate it mentally from its surroundings.
    \item \textbf{List all forces} acting \emph{on that body only}: weight ($\vec{W}=m\vec{g}$), normal reactions ($\vec{R}$), tensions ($\vec{T}$), friction ($\vec{F}$), applied pushes/pulls, etc. \textbf{Never} include forces exerted \emph{by} the body on other objects.
    \item \textbf{Draw a clear diagram}: represent the body as a dot or a small box. Draw each force as an arrow starting from the dot, pointing in the correct direction, labelled with its symbol and magnitude (if known).
    \item \textbf{Choose convenient axes} (usually horizontal/vertical or parallel/perpendicular to an incline). Resolve each force into components along these axes.
    \item \textbf{Apply Newton's First Law} (equilibrium condition): the vector sum of forces is zero.
    \[
    \sum \vec{F} = \vec{0} \quad\Longleftrightarrow\quad 
    \begin{cases}
    \sum F_x = 0 \\
    \sum F_y = 0
    \end{cases}
    \]
    \item \textbf{Solve the resulting equations} for the unknown magnitudes. Check that the solutions are physically plausible (e.g., tensions $\ge 0$, friction $\le \mu R$).
\end{enumerate}
\end{methode}

\begin{exoresolu}[A Particle Held in Equilibrium by Two Strings]
A particle of mass \SI{2}{kg} is suspended by two light inextensible strings. String $A$ is horizontal and attached to a vertical wall; string $B$ is attached to the ceiling and makes an angle of $30^\circ$ with the horizontal. The particle rests in equilibrium. Find the tension in each string. (Take $g = \SI{10}{m.s^{-2}}$).
\tcblower
\textbf{Solution.}
\begin{enumerate}
    \item \textbf{Free-body diagram:} The particle is subject to three forces:
    \begin{itemize}
        \item Weight $\vec{W} = mg = 2 \times 10 = \SI{20}{N}$, vertically downwards.
        \item Tension $\vec{T}_A$ in string $A$, horizontally towards the wall (to the left).
        \item Tension $\vec{T}_B$ in string $B$, at $30^\circ$ above the horizontal (to the right and up).
    \end{itemize}
    \begin{popfigure}
    \begin{tikzpicture}[scale=1.2]
        \draw[thick] (0,0) circle (0.15);
        \draw[->, thick, popblue] (0,0) -- (-1.5,0) node[left] {$T_A$};
        \draw[->, thick, poporange] (0,0) -- (1.3,0.75) node[above right] {$T_B$};
        \draw[->, thick, popgreen] (0,0) -- (0,-1.2) node[below] {$W=20\,\mathrm{N}$};
        \draw[dashed, thin] (0,0) -- (1.5,0);
        \draw (0.3,0) arc (0:30:0.3) node[midway, right] {$30^\circ$};
    \end{tikzpicture}
    \legende{Free-body diagram of the particle}
    \end{popfigure}
    \item \textbf{Resolve horizontally ($x$-axis positive to the right):}
    \[
    \sum F_x = 0 \;\Longrightarrow\; -T_A + T_B\cos 30^\circ = 0 \quad\Rightarrow\quad T_A = T_B\cos 30^\circ = \frac{\sqrt{3}}{2}T_B. \tag{1}
    \]
    \item \textbf{Resolve vertically ($y$-axis positive upward):}
    \[
    \sum F_y = 0 \;\Longrightarrow\; T_B\sin 30^\circ - W = 0 \quad\Rightarrow\quad T_B \times \frac{1}{2} = 20 \;\Longrightarrow\; T_B = \SI{40}{N}.
    \]
    \item \textbf{Substitute into (1):}
    \[
    T_A = \frac{\sqrt{3}}{2} \times 40 = 20\sqrt{3} \approx \SI{34.6}{N}.
    \]
    \item \textbf{Answer:} Tension in string $A$ is $\boldsymbol{20\sqrt{3}\,\mathrm{N}}$ (or $\SI{34.6}{N}$), tension in string $B$ is $\boldsymbol{40\,\mathrm{N}}$.
\end{enumerate}
\begin{attention}[Common Mistake]
Do not confuse the angle with the vertical. Here the angle is given with the horizontal, so the vertical component is $T_B\sin 30^\circ$. Always read the angle description carefully and draw the angle on your diagram.
\end{attention}
\end{exoresolu}

\begin{exercice}[Equilibrium on an Incline]
A block of mass \SI{4}{kg} rests in equilibrium on a smooth plane inclined at $30^\circ$ to the horizontal, held by a force $P$ acting parallel to the plane up the slope. Find $P$ and the normal reaction $R$. ($g=\SI{10}{m.s^{-2}}$).
\end{exercice}
\begin{corrige}[Exercise 1]
$P = mg\sin30^\circ = 4\times10\times0.5 = \SI{20}{N}$; $R = mg\cos30^\circ = 40\times\sqrt{3}/2 = 20\sqrt{3} \approx \SI{34.6}{N}$.
\end{corrige}

\courssub{Force, Acceleration and Newton's Second Law}

\begin{definition}[Linear Momentum]
The \emph{linear momentum} $\vec{p}$ of a particle of mass $m$ moving with velocity $\vec{v}$ is $\vec{p} = m\vec{v}$. It is a vector quantity (SI unit: $\mathrm{kg\,m\,s^{-1}}$ or $\mathrm{N\,s}$).
\end{definition}

\begin{propriete}[Newton's Second Law (Momentum Form)]
The rate of change of linear momentum of a body is directly proportional to the resultant external force acting on it, and takes place in the direction of that force.
\[
\vec{F}_{\text{net}} = \frac{d\vec{p}}{dt}.
\]
For a body of \emph{constant mass}, $\frac{d\vec{p}}{dt} = m\frac{d\vec{v}}{dt} = m\vec{a}$, giving the familiar form:
\[
\vec{F}_{\text{net}} = m\vec{a}.
\]
\textbf{Examinable:} You must be able to state the law in both forms and explain when $F=ma$ is valid (constant mass).
\end{propriete}

\begin{definition}[The Newton]
One \emph{newton} (\SI{1}{N}) is the force that gives a mass of \SI{1}{kg} an acceleration of \SI{1}{m.s^{-2}}. $1\,\mathrm{N} = 1\,\mathrm{kg\,m\,s^{-2}}$.
\end{definition}

\begin{methode}[Solving $F=ma$ Problems in a Straight Line]
\begin{enumerate}
    \item \textbf{Draw the free-body diagram} of the moving body (or each body if connected).
    \item \textbf{Choose the positive direction} along the line of motion (usually the direction of acceleration or the expected motion).
    \item \textbf{Write Newton's Second Law} for the body (or each body):
    \[
    \sum F_{\text{(in positive direction)}} = m a
    \]
    where $a$ is the acceleration (positive if in the chosen direction).
    \item \textbf{If friction is present}, use $F = \mu R$ for limiting friction (or $F \le \mu R$ if not slipping). The normal reaction $R$ is found from equilibrium perpendicular to the motion.
    \item \textbf{Solve the equation(s)} for the unknown ($a$, $F$, $\mu$, etc.). If multiple bodies are connected, they share the same magnitude of acceleration (and tension if the string is light and inextensible).
    \item \textbf{Use kinematics equations} ($v=u+at$, $s=ut+\frac{1}{2}at^2$, $v^2=u^2+2as$) if distances, times or velocities are required.
\end{enumerate}
\end{methode}

\begin{exoresolu}[Block on a Rough Horizontal Plane with an Angled Push]
A block of mass \SI{5}{kg} rests on a rough horizontal table. The coefficient of friction between the block and the table is $0.4$. A force of \SI{50}{N} is applied to the block at an angle of $30^\circ$ above the horizontal. Find the acceleration of the block. (Take $g = \SI{10}{m.s^{-2}}$).
\tcblower
\textbf{Solution.}
\begin{enumerate}
    \item \textbf{Free-body diagram and axes:} Choose $x$ horizontal (forward), $y$ vertical (upward).
    \begin{popfigure}
    \begin{tikzpicture}[scale=1]
        \draw[thick, fill=popblueL] (-1,-0.5) rectangle (1,0.5);
        \draw[->, thick, poporange] (0,0) -- (1.5,0.866) node[above right] {$P=50\,\mathrm{N}$};
        \draw[->, thick, popgreen] (0,0) -- (0,-1) node[below] {$W=50\,\mathrm{N}$};
        \draw[->, thick, popblue] (0,0) -- (0,1) node[above] {$R$};
        \draw[->, thick, poppurple] (0,0) -- (-1.5,0) node[left] {$F$};
        \draw[dashed, thin] (0,0) -- (1.5,0);
        \draw (0.3,0) arc (0:30:0.3) node[midway, right] {$30^\circ$};
    \end{tikzpicture}
    \legende{Forces on the block}
    \end{popfigure}
    Forces: Weight $W=mg=50\,\mathrm{N}\downarrow$; Applied force $P=50\,\mathrm{N}$ at $30^\circ$; Normal reaction $R\uparrow$; Friction $F$ opposing motion (assumed forward $\Rightarrow$ friction backward).
    \item \textbf{Vertical equilibrium (no vertical acceleration):}
    \[
    \sum F_y = 0 \;\Longrightarrow\; R + P\sin 30^\circ - W = 0
    \]
    \[
    R + 50 \times \frac{1}{2} - 50 = 0 \;\Longrightarrow\; R = 50 - 25 = \SI{25}{N}.
    \]
    \item \textbf{Friction magnitude (limiting, since the block moves):}
    \[
    F = \mu R = 0.4 \times 25 = \SI{10}{N}.
    \]
    \item \textbf{Horizontal equation (Newton's 2nd Law):}
    \[
    \sum F_x = m a \;\Longrightarrow\; P\cos 30^\circ - F = 5 a
    \]
    \[
    50 \times \frac{\sqrt{3}}{2} - 10 = 5a \;\Longrightarrow\; 25\sqrt{3} - 10 = 5a
    \]
    \[
    a = \frac{25\sqrt{3} - 10}{5} = 5\sqrt{3} - 2 \approx \SI{6.66}{m.s^{-2}}.
    \]
    \item \textbf{Answer:} Acceleration is $\boldsymbol{(5\sqrt{3}-2)\,\mathrm{m.s^{-2}}}$ (approx $\SI{6.66}{m.s^{-2}}$).
\end{enumerate}
\begin{attention}[Sign of Friction]
Always draw friction opposing the \emph{actual or impending} motion. If the calculated acceleration comes out negative, your assumed direction of motion was wrong; reverse friction and re-solve.
\end{attention}
\end{exoresolu}

\begin{exercice}[Braking Distance]
A car of mass \SI{1200}{kg} travels at \SI{25}{m.s^{-1}} on a horizontal road. The driver applies the brakes, causing a constant total resistive force of \SI{6000}{N}. Find (a) the deceleration, (b) the stopping distance. ($g=\SI{10}{m.s^{-2}}$).
\end{exercice}
\begin{corrige}[Exercise 2]
(a) $-6000 = 1200 a \Rightarrow a = -\SI{5}{m.s^{-2}}$. (b) $0 = 25^2 + 2(-5)s \Rightarrow s = \SI{62.5}{m}$.
\end{corrige}

\courssub{Newton's Second Law on an Inclined Plane (Two Dimensions)}

\begin{methode}[Motion on a Rough Inclined Plane]
\begin{enumerate}
    \item \textbf{Draw the free-body diagram} with axes: $x$ parallel to the plane (positive up or down the slope), $y$ perpendicular to the plane (positive outward).
    \item \textbf{Resolve the weight:} $W_x = mg\sin\theta$ (down the slope), $W_y = mg\cos\theta$ (into the plane).
    \item \textbf{Perpendicular equilibrium:} $R = mg\cos\theta$ (plus/minus any other force components perpendicular to the plane).
    \item \textbf{Friction:} $F = \mu R = \mu mg\cos\theta$, acting opposite to the direction of motion (or impending motion) along the plane.
    \item \textbf{Parallel equation (Newton's 2nd Law):}
    \[
    \sum F_x = m a \quad\Rightarrow\quad \text{(forces up the slope)} - \text{(forces down the slope)} = m a
    \]
    \item \textbf{Solve for $a$} (or $\mu$, $\theta$, etc.). If the body is initially at rest, check whether the component of weight down the slope exceeds the maximum static friction $\mu_s R$ before assuming motion.
\end{enumerate}
\end{methode}

\begin{exoresolu}[Block Sliding Down a Rough Incline]
A block of mass \SI{3}{kg} is placed on a rough plane inclined at $25^\circ$ to the horizontal. The coefficient of friction is $0.3$. The block is released from rest. Find: (a) the acceleration of the block; (b) the distance travelled in the first \SI{2}{s}. (Take $g = \SI{10}{m.s^{-2}}$, $\sin 25^\circ \approx 0.423$, $\cos 25^\circ \approx 0.906$).
\tcblower
\textbf{Solution.}
\begin{enumerate}
    \item \textbf{Axes:} $x$ down the slope (direction of motion), $y$ perpendicular outward.
    \item \textbf{Weight components:}
    \[
    W_x = mg\sin 25^\circ = 3 \times 10 \times 0.423 = \SI{12.69}{N},\quad
    W_y = mg\cos 25^\circ = 30 \times 0.906 = \SI{27.18}{N}.
    \]
    \item \textbf{Normal reaction:} $R = W_y = \SI{27.18}{N}$.
    \item \textbf{Friction (opposing motion, so up the slope):}
    \[
    F = \mu R = 0.3 \times 27.18 = \SI{8.154}{N}.
    \]
    \item \textbf{Newton's 2nd Law down the slope:}
    \[
    \sum F_x = m a \;\Longrightarrow\; W_x - F = 3 a
    \]
    \[
    12.69 - 8.154 = 3a \;\Longrightarrow\; 4.536 = 3a \;\Longrightarrow\; a = \SI{1.512}{m.s^{-2}}.
    \]
    \item \textbf{Distance in 2 s (from rest, $u=0$):}
    \[
    s = ut + \frac{1}{2}at^2 = 0 + \frac{1}{2} \times 1.512 \times 2^2 = \SI{3.024}{m}.
    \]
    \item \textbf{Answers:} (a) $\boldsymbol{1.51\,\mathrm{m.s^{-2}}}$ (b) $\boldsymbol{3.02\,\mathrm{m}}$.
\end{enumerate}
\begin{aretenir}[Inclined Plane Shortcut]
For a block sliding down a rough incline with no other forces:
\[
a = g(\sin\theta - \mu\cos\theta).
\]
Derive it once, then use it to save time in exams.
\end{aretenir}
\end{exoresolu}

\begin{exoresolu}[Block Projected Up a Rough Incline]
A block of mass \SI{2}{kg} is projected up a rough plane inclined at $30^\circ$ with an initial speed of \SI{10}{m.s^{-1}}. The coefficient of friction is $0.2$. Find (a) the deceleration while moving up; (b) the distance travelled before coming to rest; (c) the acceleration on the way down. ($g=\SI{10}{m.s^{-2}}$).
\tcblower
\textbf{Solution.}
\begin{enumerate}
    \item \textbf{Up the plane (positive up):} $W_x = mg\sin30^\circ = 10\,\mathrm{N}$ down; $F = \mu mg\cos30^\circ = 0.2\times20\times\sqrt{3}/2 = 2\sqrt{3}\approx\SI{3.46}{N}$ down.
    \[
    -10 - 3.46 = 2a \Rightarrow a = -6.73\,\mathrm{m.s^{-2}} \quad\text{(deceleration } 6.73\,\mathrm{m.s^{-2}}\text{)}.
    \]
    \item \textbf{Distance up:} $0 = 10^2 + 2(-6.73)s \Rightarrow s = \frac{100}{13.46} \approx \SI{7.43}{m}$.
    \item \textbf{Down the plane (positive down):} $W_x$ down, $F$ up.
    \[
    10 - 3.46 = 2a \Rightarrow a = \SI{3.27}{m.s^{-2}}.
    \]
\end{enumerate}
\end{exoresolu}

\courssub{Action–Reaction and Newton's Third Law}

\begin{propriete}[Newton's Third Law (Action–Reaction)]
If body $A$ exerts a force $\vec{F}_{AB}$ on body $B$, then body $B$ simultaneously exerts a force $\vec{F}_{BA}$ on body $A$ such that
\[
\vec{F}_{AB} = -\vec{F}_{BA}.
\]
The two forces:
\begin{itemize}
    \item have equal magnitude,
    \item act in opposite directions,
    \item act on \emph{different} bodies,
    \item are of the \emph{same nature} (both gravitational, both contact, both tension, etc.).
\end{itemize}
\textbf{Examinable:} You must be able to identify action–reaction pairs in a given situation and explain why they do \emph{not} cancel each other (they act on different bodies).
\end{propriete}

\begin{experience}[Identifying Action–Reaction Pairs]
Consider a book resting on a table.
\begin{itemize}
    \item \textbf{Pair 1 (Gravitational):} Earth pulls book down ($\vec{W}$); book pulls Earth up ($-\vec{W}$).
    \item \textbf{Pair 2 (Contact):} Table pushes book up ($\vec{R}$); book pushes table down ($-\vec{R}$).
\end{itemize}
$\vec{W}$ and $\vec{R}$ on the book are \emph{not} an action–reaction pair (they act on the same body and are of different nature). They happen to be equal and opposite because the book is in equilibrium (First Law), not because of the Third Law.
\end{experience}

\begin{attention}[Third Law Trap]
In a free-body diagram of a \emph{single} body, you never draw both forces of an action–reaction pair. You draw only the force acting \emph{on that body}. The reaction force appears on the free-body diagram of the \emph{other} body.
\end{attention}

\courssub{Connected Particles: Pulleys and Newton's Third Law}

\begin{methode}[Systems of Connected Particles (Light Inextensible Strings, Smooth Pulleys)]
\begin{enumerate}
    \item \textbf{Draw a separate free-body diagram for each particle.}
    \item \textbf{Identify the constraints:}
    \begin{itemize}
        \item The string is \emph{light} $\Rightarrow$ tension $T$ is the same throughout.
        \item The string is \emph{inextensible} $\Rightarrow$ all particles have the same magnitude of acceleration $a$.
        \item The pulley is \emph{smooth} $\Rightarrow$ tension is unchanged across it.
    \end{itemize}
    \item \textbf{Choose a positive direction for each particle} consistent with the string (e.g., upward for the hanging mass, up the slope for the mass on the incline).
    \item \textbf{Write $\sum F = ma$ for each particle} using the same $T$ and $a$.
    \item \textbf{Solve the simultaneous equations} for $a$ and $T$ (and possibly $\mu$ or $m$).
    \item \textbf{Check:} If $a$ is negative, the system accelerates in the opposite direction; tensions must be positive.
\end{enumerate}
\end{methode}

\begin{exoresolu}[Two Masses on a Rough Table Connected Over a Pulley]
A mass $A$ of \SI{4}{kg} rests on a rough horizontal table ($\mu = 0.25$). It is connected by a light inextensible string passing over a smooth pulley at the edge of the table to a hanging mass $B$ of \SI{3}{kg}. The system is released from rest. Find (a) the acceleration of the system, (b) the tension in the string. (Take $g = \SI{10}{m.s^{-2}}$).
\tcblower
\textbf{Solution.}
\begin{enumerate}
    \item \textbf{Free-body diagrams:}
    \begin{itemize}
        \item Mass $A$ (on table): Weight $W_A=40\,\mathrm{N}\downarrow$, Normal $R\uparrow$, Tension $T\rightarrow$ (to the right), Friction $F\leftarrow$ (opposing motion).
        \item Mass $B$ (hanging): Weight $W_B=30\,\mathrm{N}\downarrow$, Tension $T\uparrow$.
    \end{itemize}
    \item \textbf{Assume motion:} $B$ falls, $A$ moves right $\Rightarrow$ $a$ is downward for $B$, rightward for $A$.
    \item \textbf{Equations:}
    \begin{align*}
    \text{For } A \text{ (horizontal):} &\quad T - F = 4a \tag{1} \\
    \text{Vertical on } A: &\quad R = W_A = 40\,\mathrm{N} \;\Rightarrow\; F = \mu R = 0.25 \times 40 = \SI{10}{N}. \\
    \text{For } B \text{ (vertical down positive):} &\quad 30 - T = 3a \tag{2}
    \end{align*}
    \item \textbf{Solve (1) and (2):}
    From (1): $T = 4a + 10$.
    Substitute into (2): $30 - (4a + 10) = 3a \;\Rightarrow\; 20 = 7a \;\Rightarrow\; a = \frac{20}{7} \approx \SI{2.86}{m.s^{-2}}$.
    Then $T = 4 \times \frac{20}{7} + 10 = \frac{80}{7} + \frac{70}{7} = \frac{150}{7} \approx \SI{21.4}{N}$.
    \item \textbf{Answers:} (a) $\boldsymbol{\frac{20}{7}\,\mathrm{m.s^{-2}}}$ (b) $\boldsymbol{\frac{150}{7}\,\mathrm{N}}$.
\end{enumerate}
\begin{attention}[Direction of Friction]
If you are unsure which way the system moves, assume a direction. If $a$ comes out negative, the motion is opposite; friction reverses. The magnitude of $a$ and $T$ will be correct.
\end{attention}
\end{exoresolu}

\begin{exoresolu}[Two Blocks on Two Inclines Connected Over a Pulley]
Two blocks $P$ (mass \SI{2}{kg}) and $Q$ (mass \SI{5}{kg}) are connected by a light inextensible string passing over a smooth pulley. $P$ rests on a smooth plane inclined at $30^\circ$; $Q$ rests on a rough plane inclined at $45^\circ$ ($\mu = 0.2$). The string is parallel to each plane. The system is released from rest. Find (a) the acceleration of the system, (b) the tension in the string. (Take $g = \SI{10}{m.s^{-2}}$, $\sin 30^\circ=0.5$, $\cos 30^\circ=\sqrt{3}/2$, $\sin 45^\circ=\cos 45^\circ=\sqrt{2}/2$).
\tcblower
\textbf{Solution.}
\begin{enumerate}
    \item \textbf{System diagram and assumption:} $Q$ is heavier and on a steeper incline; assume $Q$ moves down its plane, $P$ moves up its plane. Acceleration magnitude $a$ same for both.
    \item \textbf{Free-body diagrams:}
    \begin{itemize}
        \item $P$ (smooth, $30^\circ$): Weight $2g$, Normal $R_P$, Tension $T$ up the plane. Axis: $x$ up the plane.
        \item $Q$ (rough, $45^\circ$): Weight $5g$, Normal $R_Q$, Tension $T$ up the plane, Friction $F$ up the plane (opposing motion down). Axis: $x$ down the plane.
    \end{itemize}
    \item \textbf{Equations for $P$ (up positive):}
    \[
    T - 2g\sin 30^\circ = 2a \;\Longrightarrow\; T - 20 \times 0.5 = 2a \;\Longrightarrow\; T - 10 = 2a. \tag{1}
    \]
    \item \textbf{Equations for $Q$ (down positive):}
    \[
    \text{Perpendicular: } R_Q = 5g\cos 45^\circ = 50 \times \frac{\sqrt{2}}{2} = 25\sqrt{2}\,\mathrm{N}.
    \]
    \[
    \text{Friction: } F = \mu R_Q = 0.2 \times 25\sqrt{2} = 5\sqrt{2}\,\mathrm{N}.
    \]
    \[
    \text{Parallel: } 5g\sin 45^\circ - T - F = 5a
    \]
    \[
    50 \times \frac{\sqrt{2}}{2} - T - 5\sqrt{2} = 5a \;\Longrightarrow\; 25\sqrt{2} - 5\sqrt{2} - T = 5a
    \]
    \[
    20\sqrt{2} - T = 5a. \tag{2}
    \]
    \item \textbf{Solve (1) and (2):} Add (1) and (2):
    \[
    (T - 10) + (20\sqrt{2} - T) = 2a + 5a \;\Longrightarrow\; 20\sqrt{2} - 10 = 7a
    \]
    \[
    a = \frac{20\sqrt{2} - 10}{7} \approx \frac{28.28 - 10}{7} = \frac{18.28}{7} \approx \SI{2.61}{m.s^{-2}}.
    \]
    From (1): $T = 2a + 10 = 2 \times 2.61 + 10 \approx \SI{15.22}{N}$.
    Exact: $T = 2\left(\frac{20\sqrt{2} - 10}{7}\right) + 10 = \frac{40\sqrt{2} - 20 + 70}{7} = \frac{40\sqrt{2} + 50}{7}\,\mathrm{N}$.
    \item \textbf{Answers:} (a) $\boldsymbol{\frac{20\sqrt{2}-10}{7}\,\mathrm{m.s^{-2}}}$ (b) $\boldsymbol{\frac{40\sqrt{2}+50}{7}\,\mathrm{N}}$.
\end{enumerate}
\begin{aretenir}[Exam Technique]
In multi-body problems, always define the positive direction for each body \emph{before} writing equations. Write the constraint ($a$ same magnitude) explicitly. Solve the system by elimination (add/subtract equations to eliminate $T$ first).
\end{aretenir}
\end{exoresolu}

\courssub{Applications of Newton's Laws in Daily Life}

\begin{experience}[Motion in a Lift – Apparent Weight]
A person of mass $m$ stands on a weighing scale (which measures the normal reaction $R$) inside a lift accelerating vertically with acceleration $a$ (upward positive). The forces on the person are weight $mg\downarrow$ and normal reaction $R\uparrow$. Newton's Second Law in the inertial (ground) frame gives:
\[
R - mg = ma \quad\Rightarrow\quad R = m(g + a).
\]
This $R$ is the \emph{apparent weight}.
\begin{itemize}
    \item Lift accelerating upward ($a>0$): $R > mg$ (feels heavier).
    \item Lift accelerating downward ($a<0$): $R < mg$ (feels lighter).
    \item Constant velocity ($a=0$): $R = mg$ (true weight).
    \item Free fall ($a = -g$): $R = 0$ (weightlessness).
\end{itemize}
\end{experience}

\begin{exoresolu}[Reading on a Scale in an Accelerating Lift]
A man of mass \SI{70}{kg} stands on a bathroom scale (calibrated in newtons) inside a lift. Find the scale reading when: (a) the lift accelerates upward at \SI{2}{m.s^{-2}}; (b) the lift accelerates downward at \SI{3}{m.s^{-2}}; (c) the lift moves downward at constant velocity \SI{4}{m.s^{-1}}; (d) the lift cable snaps and the lift falls freely. (Take $g = \SI{10}{m.s^{-2}}$).
\tcblower
\textbf{Solution.}
Use $R = m(g + a)$ with upward positive.
\begin{enumerate}
    \item (a) $a = +\SI{2}{m.s^{-2}}$: $R = 70(10 + 2) = \SI{840}{N}$.
    \item (b) $a = -\SI{3}{m.s^{-2}}$: $R = 70(10 - 3) = \SI{490}{N}$.
    \item (c) Constant velocity $\Rightarrow a = 0$: $R = 70 \times 10 = \SI{700}{N}$ (true weight).
    \item (d) Free fall $\Rightarrow a = -g = -\SI{10}{m.s^{-2}}$: $R = 70(10 - 10) = \SI{0}{N}$.
\end{enumerate}
\textbf{Answers:} (a) $\boldsymbol{840\,\mathrm{N}}$ (b) $\boldsymbol{490\,\mathrm{N}}$ (c) $\boldsymbol{700\,\mathrm{N}}$ (d) $\boldsymbol{0\,\mathrm{N}}$.
\begin{savaistu}[Weightlessness in Orbit]
Astronauts in the International Space Station experience weightlessness not because gravity is zero (it's about $90\%$ of surface gravity at that altitude), but because they are in continuous free fall ($a = -g$), so the normal reaction from the spacecraft floor is zero.
\end{savaistu}
\end{exoresolu}

\begin{experience}[Vehicle Dynamics – Braking and Turning]
\begin{itemize}
    \item \textbf{Braking on a horizontal road:} Maximum deceleration without skidding is $a_{\max} = \mu g$, where $\mu$ is the coefficient of static friction between tyres and road. For a typical dry road $\mu \approx 0.7$, giving $a_{\max} \approx 7\,\mathrm{m.s^{-2}}$.
    \item \textbf{Banked curves:} On a frictionless banked curve of angle $\theta$ and radius $r$, the ideal speed for no side friction is $v = \sqrt{rg\tan\theta}$. With friction, there is a range of safe speeds.
    \item \textbf{Overturning:} A vehicle of height $h$ (centre of mass) and track width $2d$ will overturn before skidding if $\frac{v^2}{rg} > \frac{d}{h}$.
\end{itemize}
\end{experience}

\begin{exoresolu}[Banked Curve with Friction]
A car travels on a circular bend of radius \SI{50}{m} banked at $15^\circ$. The coefficient of friction between tyres and road is $0.3$. Find the maximum speed without skidding outward. ($g=\SI{10}{m.s^{-2}}$, $\tan15^\circ\approx0.268$, $\sin15^\circ\approx0.259$, $\cos15^\circ\approx0.966$).
\tcblower
\textbf{Solution.}
Forces: Weight $mg\downarrow$, Normal $R\perp$ plane, Friction $F$ down the plane (preventing outward skid).
Resolve vertically: $R\cos15^\circ - F\sin15^\circ = mg$.
Resolve horizontally (toward centre): $R\sin15^\circ + F\cos15^\circ = \frac{mv^2}{r}$.
At maximum speed, $F = \mu R = 0.3R$.
Substitute $F$:
$R(\cos15^\circ - 0.3\sin15^\circ) = mg \Rightarrow R = \frac{mg}{0.966 - 0.3\times0.259} = \frac{mg}{0.888}$.
Horizontal: $R(\sin15^\circ + 0.3\cos15^\circ) = \frac{mv^2}{50}$.
$\frac{mg}{0.888}(0.259 + 0.3\times0.966) = \frac{mv^2}{50}$.
$v^2 = 50 \times 10 \times \frac{0.259 + 0.290}{0.888} = 500 \times \frac{0.549}{0.888} \approx 309$.
$v \approx \SI{17.6}{m.s^{-1}}$ ($\approx \SI{63.4}{km.h^{-1}}$).
\end{exoresolu}

\courssub{Limitations and Extensions of Newton's Laws}

\begin{propriete}[Limitations of Newton's Laws]
Newton's laws are supremely accurate for macroscopic objects at everyday speeds, but they fail in the following regimes:
\begin{itemize}
    \item \textbf{High speeds (Relativistic limit):} When $v \gtrsim 0.1c$ ($c = \SI{3\times10^8}{m.s^{-1}}$), relativistic mechanics must be used. The correct form is $\vec{F} = \frac{d}{dt}(\gamma m\vec{v})$ where $\gamma = (1-v^2/c^2)^{-1/2}$. Mass is not constant; momentum is not $mv$.
    \item \textbf{Quantum scale:} At atomic/subatomic scales ($\lesssim \SI{10^{-10}}{m}$), quantum mechanics replaces deterministic trajectories with probabilities. The concept of a definite force acting on a definite trajectory breaks down.
    \item \textbf{Non-inertial frames:} Newton's laws in their simple form ($\vec{F}=m\vec{a}$) do not hold in accelerating frames unless fictitious (pseudo) forces are added. For a frame with acceleration $\vec{a}_{\text{frame}}$, add $-m\vec{a}_{\text{frame}}$ to the real forces.
    \item \textbf{Variable mass systems:} $F = ma$ is \textbf{not} directly applicable if mass changes (e.g., rocket losing fuel, raindrop gaining mass). Use the momentum form: $\vec{F}_{\text{ext}} = \frac{d\vec{p}}{dt} = \frac{d}{dt}(m\vec{v})$.
\end{itemize}
\end{propriete}

\begin{methode}[Variable Mass Systems – Rocket Equation (1D)]
\begin{enumerate}
    \item \textbf{System:} Rocket of instantaneous mass $m$, velocity $v$. In time $dt$, mass $dm$ (negative) is ejected at relative speed $u$ backward.
    \item \textbf{Momentum conservation (no external forces):}
    \[
    mv = (m+dm)(v+dv) + (-dm)(v-u).
    \]
    \item \textbf{Simplify:} $mv = mv + m\,dv + v\,dm + dm\,dv - v\,dm + u\,dm$.
    Neglect $dm\,dv$ (second order): $0 = m\,dv + u\,dm \;\Rightarrow\; m\,dv = -u\,dm$.
    \item \textbf{Integrate:} $\int_{v_0}^{v} dv = -u \int_{m_0}^{m} \frac{dm}{m} \;\Rightarrow\; v - v_0 = u \ln\left(\frac{m_0}{m}\right)$.
    \item \textbf{With external force (e.g., gravity):} $m\frac{dv}{dt} = F_{\text{ext}} + u\frac{dm}{dt}$ (thrust $= u\frac{dm}{dt}$, note $\frac{dm}{dt}<0$).
\end{enumerate}
\end{methode}

\begin{exoresolu}[Rocket Launch – Variable Mass]
A rocket of initial mass \SI{1000}{kg} (including \SI{800}{kg} of fuel) burns fuel at a constant rate of \SI{50}{kg.s^{-1}}. The exhaust gases are ejected at a constant speed of \SI{2000}{m.s^{-1}} relative to the rocket. Ignoring air resistance and assuming $g = \SI{10}{m.s^{-2}}$ constant, find: (a) the thrust force; (b) the acceleration at lift-off; (c) the velocity at burnout (when all fuel is exhausted), assuming vertical flight.
\tcblower
\textbf{Solution.}
\begin{enumerate}
    \item \textbf{Thrust:} $T = u \left|\frac{dm}{dt}\right| = 2000 \times 50 = \SI{100000}{N} = \SI{100}{kN}$.
    \item \textbf{At lift-off ($t=0$, $m_0 = 1000\,\mathrm{kg}$):} Net upward force $= T - m_0g = 100000 - 10000 = \SI{90000}{N}$.
    \[
    a_0 = \frac{90000}{1000} = \SI{90}{m.s^{-2}}.
    \]
    \item \textbf{Burnout time:} $t_b = \frac{800}{50} = \SI{16}{s}$. Mass at burnout $m_b = 1000 - 800 = \SI{200}{kg}$.
    \item \textbf{Velocity at burnout (with gravity):} Equation of motion: $m\frac{dv}{dt} = T - mg = u\frac{dm}{dt} - mg$ (since $\frac{dm}{dt} = -50$).
    \[
    \frac{dv}{dt} = \frac{u}{m}\frac{dm}{dt} - g.
    \]
    Integrate from $t=0$ to $t_b$:
    \[
    \int_0^{v_b} dv = u \int_{m_0}^{m_b} \frac{dm}{m} - \int_0^{t_b} g\,dt
    \]
    \[
    v_b = u \ln\left(\frac{m_0}{m_b}\right) - g t_b = 2000 \ln\left(\frac{1000}{200}\right) - 10 \times 16
    \]
    \[
    v_b = 2000 \ln 5 - 160 \approx 2000 \times 1.6094 - 160 = 3218.8 - 160 = \SI{3058.8}{m.s^{-1}}.
    \]
    \item \textbf{Answers:} (a) $\boldsymbol{100\,\mathrm{kN}}$ (b) $\boldsymbol{90\,\mathrm{m.s^{-2}}}$ (c) $\boldsymbol{\approx 3060\,\mathrm{m.s^{-1}}}$.
\end{enumerate}
\begin{attention}[Variable Mass Trap]
Never write $F = m\frac{dv}{dt}$ for a rocket losing mass. The correct form is $F_{\text{ext}} = \frac{d}{dt}(mv) = m\frac{dv}{dt} + v\frac{dm}{dt}$. The thrust arises from the momentum carried away by the ejected mass. Use the derived rocket equation or start from momentum conservation.
\end{attention}
\end{exoresolu}

\begin{experience}[Raindrop Gaining Mass]
A raindrop falls through a cloud, gaining mass by condensation at a rate $\frac{dm}{dt} = km$ (proportional to its current mass). If air resistance is negligible, the equation of motion is $mg = \frac{d}{dt}(mv) = m\frac{dv}{dt} + v\frac{dm}{dt}$. With $\frac{dm}{dt}=km$, this becomes $g = \frac{dv}{dt} + kv$. The terminal velocity (when $\frac{dv}{dt}=0$) is $v_{\text{term}} = g/k$.
\end{experience}

\courssub{Composite Problem: Incline, Pulley, and Friction}

\begin{methode}[General Strategy for Complex Multi-Body Mechanics Problems]
\begin{enumerate}
    \item \textbf{Read the whole problem} and visualise the system. Identify all bodies, constraints (strings, pulleys, contact), and given data.
    \item \textbf{Draw a clear system diagram} showing the geometry (angles, pulley positions).
    \item \textbf{Draw a separate free-body diagram for each body} with its own coordinate axes (aligned with its acceleration).
    \item \textbf{Write the constraint equations} (e.g., $a_A = a_B$, $T_1 = T_2$ for smooth pulley).
    \item \textbf{Write Newton's 2nd Law for each body} in its chosen coordinate directions.
    \item \textbf{Count equations and unknowns.} You should have as many independent scalar equations as unknowns ($a$, $T$, $R$, $F$, etc.).
    \item \textbf{Solve algebraically first}, then substitute numbers. Keep $g$ symbolic until the end if possible.
    \item \textbf{Verify:} Check signs (does the assumed motion direction hold?), check limiting cases (e.g., $\mu=0$, $\theta=0$).
\end{enumerate}
\end{methode}

\begin{exoresolu}[Three-Body System: Table, Pulley, and Hanging Masses]
A mass $A$ (\SI{2}{kg}) on a rough horizontal table ($\mu=0.2$) is connected by a light string to mass $B$ (\SI{3}{kg}) hanging vertically over a smooth pulley at the table edge. A second light string connects $B$ to another hanging mass $C$ (\SI{1}{kg}) over a second smooth pulley. The system is released from rest. Find the acceleration and the tensions $T_1$ (between A and B) and $T_2$ (between B and C). ($g=\SI{10}{m.s^{-2}}$).
\tcblower
\textbf{Solution.}
\begin{enumerate}
    \item \textbf{Assume motion:} $B$ and $C$ fall, $A$ moves right. Common acceleration $a$.
    \item \textbf{Free-body diagrams and equations:}
    \begin{itemize}
        \item $A$ (horizontal): $T_1 - F = 2a$; $R = 20\,\mathrm{N}$, $F = \mu R = 4\,\mathrm{N}$.
        $\Rightarrow T_1 - 4 = 2a$. \hfill (1)
        \item $B$ (vertical down): $30 + T_2 - T_1 = 3a$. \hfill (2)
        \item $C$ (vertical down): $10 - T_2 = 1a$. \hfill (3)
    \end{itemize}
    \item \textbf{Solve:} From (3): $T_2 = 10 - a$.
    Substitute into (2): $30 + 10 - a - T_1 = 3a \Rightarrow 40 - T_1 = 4a \Rightarrow T_1 = 40 - 4a$.
    Substitute into (1): $40 - 4a - 4 = 2a \Rightarrow 36 = 6a \Rightarrow a = \SI{6}{m.s^{-2}}$.
    Then $T_1 = 40 - 24 = \SI{16}{N}$, $T_2 = 10 - 6 = \SI{4}{N}$.
    \item \textbf{Check:} All tensions positive, $a$ positive $\Rightarrow$ assumption correct.
\end{enumerate}
\end{exoresolu}

\begin{aretenir}[Summary of Key Formulas for GCE Revision]
\begin{center}
\begin{tabular}{|l|l|}
\hline
\textbf{Situation} & \textbf{Key Equation(s)} \\
\hline
Equilibrium (1st Law) & $\sum F_x = 0,\; \sum F_y = 0$ \\
\hline
Constant mass (2nd Law) & $\sum \vec{F} = m\vec{a}$ \\
\hline
Inclined plane (rough, sliding down) & $a = g(\sin\theta - \mu\cos\theta)$ \\
\hline
Connected particles (light string, smooth pulley) & Same $a$, same $T$ \\
\hline
Lift (apparent weight) & $R = m(g+a)$ \\
\hline
Rocket (no external force) & $v = v_0 + u\ln(m_0/m)$ \\
\hline
Rocket (with gravity) & $v = v_0 + u\ln(m_0/m) - gt$ \\
\hline
Friction (limiting) & $F = \mu R$ \\
\hline
\end{tabular}
\end{center}
\end{aretenir}

\begin{exercice}[GCE-Style Mixed Problems]
\begin{enumerate}
    \item A block of mass \SI{6}{kg} is pulled up a rough plane inclined at $30^\circ$ by a force of \SI{50}{N} parallel to the plane. The coefficient of friction is $0.25$. Find the acceleration. ($g=\SI{10}{m.s^{-2}}$).
    \item Two particles of masses \SI{5}{kg} and \SI{3}{kg} are connected by a light inextensible string passing over a smooth pulley. The \SI{5}{kg} mass rests on a rough horizontal table ($\mu=0.2$), the \SI{3}{kg} mass hangs freely. Find the tension and the distance moved in the first \SI{2}{s} from rest.
    \item A rocket of initial mass \SI{500}{kg} ejects gas at \SI{10}{kg.s^{-1}} with a relative speed of \SI{1500}{m.s^{-1}}. Find the thrust and the initial acceleration if launched vertically ($g=\SI{10}{m.s^{-2}}$).
    \item A man of mass \SI{80}{kg} stands in a lift. The lift accelerates upward at \SI{1.5}{m.s^{-2}} for \SI{3}{s}, then moves at constant velocity for \SI{5}{s}, then decelerates at \SI{2}{m.s^{-2}} to rest. Sketch a graph of the scale reading against time.
\end{enumerate}
\end{exercice}

\begin{corrige}[Exercise 3]
\begin{enumerate}
    \item $R = 60\cos30^\circ = 30\sqrt{3}\approx\SI{51.96}{N}$, $F = 0.25R \approx \SI{12.99}{N}$. Up plane: $50 - 60\sin30^\circ - 12.99 = 6a \Rightarrow 50 - 30 - 12.99 = 6a \Rightarrow a \approx \SI{1.17}{m.s^{-2}}$.
    \item For $5\,\mathrm{kg}$: $T - F = 5a$, $F = 0.2\times50 = 10\,\mathrm{N}$. For $3\,\mathrm{kg}$: $30 - T = 3a$. Solving: $30 - 10 = 8a \Rightarrow a = \SI{2.5}{m.s^{-2}}$, $T = 30 - 7.5 = \SI{22.5}{N}$. Distance: $s = \frac{1}{2}\times2.5\times4 = \SI{5}{m}$.
    \item Thrust $= 1500 \times 10 = \SI{15000}{N}$. Weight $= 5000\,\mathrm{N}$. Net force $= 10000\,\mathrm{N}$. $a_0 = \frac{10000}{500} = \SI{20}{m.s^{-2}}$.
    \item Scale reading $R = m(g+a)$: $0\le t<3$: $R=80(11.5)=\SI{920}{N}$; $3\le t<8$: $R=80(10)=\SI{800}{N}$; $8\le t\le 9$: $R=80(8)=\SI{640}{N}$; $t>9$: $R=\SI{800}{N}$. Graph: horizontal segments at 920, 800, 640, 800 N.
\end{enumerate}
\end{corrige}

\newpage

\coursec{Exercises}

\courssub{I check my knowledge}

\begin{exercice}[Multiple choice: Newton's laws]
For each statement, choose the single correct answer.
\begin{enumerate}
\item A body moves with constant velocity on a smooth horizontal surface. The resultant force acting on it is:
\begin{enumerate}
\item equal to its weight;
\item zero;
\item directed forwards;
\item directed backwards.
\end{enumerate}
\item A force of \SI{20}{N} gives a body an acceleration of \SI{4}{m.s^{-2}}. The mass of the body is:
\begin{enumerate}
\item \SI{80}{kg};
\item \SI{5}{kg};
\item \SI{0.2}{kg};
\item \SI{24}{kg}.
\end{enumerate}
\item A book rests on a table. The reaction partner of the \emph{weight} of the book is:
\begin{enumerate}
\item the normal contact force of the table on the book;
\item the force of the book on the table;
\item the gravitational pull of the book on the Earth;
\item the friction between book and table.
\end{enumerate}
\item Newton's laws hold exactly:
\begin{enumerate}
\item in every reference frame;
\item only in inertial reference frames;
\item only in frames attached to the Earth;
\item only for bodies at rest.
\end{enumerate}
\end{enumerate}
\end{exercice}

\begin{exercice}[True or false? Justify]
State whether each statement is true or false, and justify your answer in one or two sentences.
\begin{enumerate}
\item ``A force is needed to keep a body moving at constant velocity.''
\item ``Action and reaction forces act on the same body.''
\item ``If the resultant force on a body doubles, its acceleration doubles, provided its mass is unchanged.''
\item ``A passenger in a car that brakes suddenly is thrown forward because a force pushes him forward.''
\end{enumerate}
\end{exercice}

\begin{exercice}[Definitions and statements]
\begin{enumerate}
\item State Newton's first law of motion and define the term \emph{inertia}.
\item State Newton's second law in words and write its vector form, defining each symbol and its SI unit.
\item State Newton's third law and give the three properties that an action--reaction pair must satisfy.
\item Define an \emph{inertial reference frame} and give one example of a frame that is (approximately) inertial and one that is not.
\end{enumerate}
\end{exercice}

\begin{exercice}[Direct calculations]
Take $g = \SI{9.81}{m.s^{-2}}$.
\begin{enumerate}
\item Find the resultant force needed to give a mass of \SI{800}{g} an acceleration of \SI{2.5}{m.s^{-2}}.
\item A resultant force of \SI{60}{N} acts on a body of mass \SI{15}{kg} initially at rest. Find its velocity after \SI{4}{s}.
\item Calculate the weight of a student of mass \SI{55}{kg}.
\item A lift of mass \SI{600}{kg} accelerates upwards at \SI{1.2}{m.s^{-2}}. Find the tension in the cable, neglecting friction.
\end{enumerate}
\end{exercice}

\courssub{I practise}

\begin{exercice}[Braking car (GCE Paper 2 style)]
A car of mass \SI{1200}{kg} travelling at \SI{30}{m.s^{-1}} brakes uniformly and comes to rest in \SI{5}{s}.
\begin{enumerate}
\item Calculate the deceleration of the car.
\item Calculate the average braking force.
\item Calculate the distance travelled while braking.
\item State two factors that would increase the braking distance in practice.
\end{enumerate}
\end{exercice}

\begin{exercice}[Atwood machine]
Two blocks of masses \SI{3}{kg} and \SI{5}{kg} are connected by a light inextensible string passing over a smooth fixed pulley. The system is released from rest. Take $g = \SI{9.81}{m.s^{-2}}$.
\begin{enumerate}
\item Draw a diagram showing the forces acting on each block.
\item Write down the equation of motion for each block.
\item Find the acceleration of the system.
\item Find the tension in the string.
\item Find the velocity of the \SI{5}{kg} block after it has descended \SI{1}{m}.
\end{enumerate}
\end{exercice}

\begin{exercice}[Rocket propulsion]
A rocket of initial mass \SI{500}{kg} ejects gas at a steady rate of \SI{2}{kg.s^{-1}} with exhaust velocity \SI{800}{m.s^{-1}} relative to the rocket. Ignore gravity and air resistance.
\begin{enumerate}
\item Show that the thrust on the rocket is \SI{1600}{N}.
\item Determine the acceleration of the rocket at $t = 0$.
\item Determine the mass of the rocket at $t = \SI{10}{s}$ and hence its acceleration at that instant.
\item Explain why the acceleration increases with time even though the thrust is constant.
\end{enumerate}
\end{exercice}

\begin{exercice}[Forces in a turning car]
A passenger of mass \SI{60}{kg} sits in a car that takes a sharp circular bend of radius \SI{40}{m} at a constant speed of \SI{15}{m.s^{-1}}.
\begin{enumerate}
\item Calculate the centripetal acceleration of the passenger.
\item Calculate the resultant force on the passenger and state what provides it.
\item Explain why the passenger \emph{feels} pushed outwards.
\item Identify the real forces acting on the passenger, and name the fictitious force that appears in the rotating frame of the car.
\end{enumerate}
\end{exercice}

\courssub{I prepare for the GCE}

\begin{exercice}[Connected bodies on a slope (GCE Advanced Level style)]
A block $A$ of mass \SI{4}{kg} rests on a rough plane inclined at $30^{\circ}$ to the horizontal. It is connected by a light inextensible string, passing over a smooth pulley at the top of the plane, to a block $B$ of mass \SI{6}{kg} hanging freely. The coefficient of friction between $A$ and the plane is $0.2$. Take $g = \SI{9.81}{m.s^{-2}}$.
\begin{enumerate}
\item Draw a clear diagram showing all the forces acting on each block. \hfill (3 marks)
\item Calculate the normal reaction on $A$ and the limiting frictional force. \hfill (3 marks)
\item Show that the system moves and find the acceleration of the blocks. \hfill (5 marks)
\item Find the tension in the string. \hfill (2 marks)
\item After $B$ has fallen \SI{2}{m}, the string breaks. Describe the subsequent motion of $A$, justifying your answer. \hfill (3 marks)
\end{enumerate}
\hfill (Total: 16 marks)
\end{exercice}

\begin{exercice}[Lift, reaction and apparent weight (GCE Advanced Level style)]
A man of mass \SI{70}{kg} stands on a weighing scale inside a lift. Take $g = \SI{9.81}{m.s^{-2}}$.
\begin{enumerate}
\item State Newton's second law and explain how it applies to the man in the accelerating lift. \hfill (3 marks)
\item Find the reading of the scale (in newtons) when the lift:
\begin{enumerate}
\item moves upwards with constant velocity; \hfill (2 marks)
\item accelerates upwards at \SI{1.5}{m.s^{-2}}; \hfill (3 marks)
\item accelerates downwards at \SI{1.5}{m.s^{-2}}. \hfill (3 marks)
\end{enumerate}
\item Explain, using Newton's laws, why the man feels ``heavier'' in case (ii) and ``lighter'' in case (iii). \hfill (3 marks)
\item The cable snaps and the lift falls freely. What is the reading of the scale? Comment on the term \emph{weightlessness}. \hfill (3 marks)
\end{enumerate}
\hfill (Total: 17 marks)
\end{exercice}

\begin{exercice}[Essay: scope and limits of Newton's laws (GCE Advanced Level style)]
\begin{enumerate}
\item Critically discuss the statement: ``Newton's third law implies that forces always cancel, so nothing can ever accelerate.'' Provide a correct explanation, supported by a clearly identified action--reaction pair in a concrete example (for instance a person walking, or a pirogue pushed from a river bank). \hfill (6 marks)
\item Explain, with one example each, why Newton's laws cease to give accurate results:
\begin{enumerate}
\item at speeds close to the speed of light;
\item at the scale of atoms and subatomic particles.
\end{enumerate}
Name the theories that extend or replace Newtonian mechanics in these domains. \hfill (6 marks)
\item ``The Earth is not a perfect inertial frame, yet we apply Newton's laws successfully in everyday life in Cameroon.'' Discuss this statement, giving one situation where the non-inertial character of the Earth matters and one where it can safely be ignored. \hfill (4 marks)
\end{enumerate}
\hfill (Total: 16 marks)
\end{exercice}

\newpage

\coursec{Solutions to the exercises}

\begin{corrige}[Exercise 1 --- Multiple choice: Newton's laws]
\begin{enumerate}
\item \textbf{Answer: (b) zero.} Constant velocity means zero acceleration. By Newton's first law (equivalently the second law with $\vec{a}=\vec{0}$), the resultant force is zero: the weight is balanced by the normal reaction of the surface.
\item \textbf{Answer: (b) \SI{5}{kg}.} From $F=ma$: $m=\dfrac{F}{a}=\dfrac{20}{4}=\SI{5}{kg}$.
\item \textbf{Answer: (c) the gravitational pull of the book on the Earth.} An action--reaction pair consists of two forces of the \emph{same physical nature} acting on \emph{different bodies}. The weight is the Earth's gravitational pull on the book; its partner is the book's gravitational pull on the Earth. The normal contact force is \emph{not} the partner of the weight: both act on the book, and they need not even be equal if the book accelerates vertically.
\item \textbf{Answer: (b) only in inertial reference frames.} Newton's laws are stated relative to inertial (Galilean) frames; in accelerating frames, fictitious (inertial) forces must be introduced to keep using them.
\end{enumerate}
\end{corrige}

\begin{corrige}[Exercise 2 --- True or false? Justify]
\begin{enumerate}
\item \textbf{False.} By Newton's first law, a body in uniform rectilinear motion needs \emph{no} resultant force to maintain its motion. A resultant force is only needed to \emph{change} the velocity. (If an engine force is present, it merely balances friction: the \emph{resultant} is still zero.)
\item \textbf{False.} Action and reaction act on \emph{different} bodies: the action on body $B$ by $A$, the reaction on body $A$ by $B$. That is precisely why they can never cancel each other --- cancellation only makes sense for forces acting on the \emph{same} body.
\item \textbf{True.} From $\vec{F}=m\vec{a}$ with $m$ constant, $a=\dfrac{F}{m}$, so $a$ is directly proportional to $F$: doubling $F$ doubles $a$.
\item \textbf{False.} No forward force acts on the passenger. By inertia, his body tends to continue at the car's original velocity while the car slows down; the \emph{relative} motion forwards is a consequence of Newton's first law, not of a pushing force. (This is why seat belts are compulsory.)
\end{enumerate}
\end{corrige}

\begin{corrige}[Exercise 3 --- Definitions and statements]
\begin{enumerate}
\item \textbf{Newton's first law:} In an inertial reference frame, a body remains at rest or in uniform rectilinear motion unless acted upon by a non-zero resultant external force. \textbf{Inertia} is the natural tendency of a body to resist any change in its state of motion; mass is the quantitative measure of inertia.
\item \textbf{Newton's second law:} In an inertial frame, the resultant external force acting on a body equals the product of its mass and its acceleration (equivalently, the rate of change of its momentum):
\[ \vec{F}=m\vec{a}, \qquad\text{i.e. } \vec{F}=\dfrac{d\vec{p}}{dt},\ \vec{p}=m\vec{v}. \]
$\vec{F}$: resultant force, in newtons ($\mathrm{N}$); $m$: mass, in kilograms ($\mathrm{kg}$); $\vec{a}$: acceleration, in $\mathrm{m\,s^{-2}}$.
\item \textbf{Newton's third law:} If body $A$ exerts a force $\vec{F}_{A/B}$ on body $B$, then body $B$ exerts a force $\vec{F}_{B/A}$ on body $A$ such that $\vec{F}_{B/A}=-\vec{F}_{A/B}$. The pair must satisfy: (i) equal magnitudes; (ii) opposite directions along the same line of action; (iii) the two forces act on \emph{different} bodies and are of the same physical nature.
\item An \textbf{inertial reference frame} is a frame in which Newton's first law holds: a body free of any resultant force moves with constant velocity. \emph{Approximately inertial:} a frame fixed to the Earth's surface (for everyday, short-duration motions), or better, a frame fixed relative to distant stars. \emph{Non-inertial:} a car accelerating or turning, a rotating merry-go-round, a lift during acceleration.
\end{enumerate}
\end{corrige}

\begin{corrige}[Exercise 4 --- Direct calculations]
\begin{enumerate}
\item Convert the mass: $m=\SI{800}{g}=\SI{0.8}{kg}$. Then
\[ F=ma=0.8\times 2.5=\SI{2}{N}. \]
\item $a=\dfrac{F}{m}=\dfrac{60}{15}=\SI{4}{m.s^{-2}}$. Starting from rest with constant acceleration:
\[ v=u+at=0+4\times 4=\SI{16}{m.s^{-1}}. \]
\item $W=mg=55\times 9.81=\SI{539.55}{N}\approx\SI{540}{N}$.
\item Forces on the lift: tension $T$ upwards, weight $mg$ downwards. Taking upwards as positive, Newton's second law gives
\[ T-mg=ma \;\Longrightarrow\; T=m(g+a)=600\times(9.81+1.2)=600\times 11.01=\SI{6606}{N}. \]
The tension is about $\SI{6.61}{kN}$, greater than the weight, as expected for upward acceleration.
\end{enumerate}
\end{corrige}

\begin{corrige}[Exercise 5 --- Braking car (GCE Paper 2 style)]
\begin{enumerate}
\item Uniform deceleration: $v=u+at$ with $v=0$, $u=\SI{30}{m.s^{-1}}$, $t=\SI{5}{s}$:
\[ a=\dfrac{0-30}{5}=-\SI{6}{m.s^{-2}}, \]
so the deceleration is $\SI{6}{m.s^{-2}}$.
\item Magnitude of the average braking force:
\[ F=m|a|=1200\times 6=\SI{7200}{N}. \]
The force is directed opposite to the motion.
\item Distance travelled:
\[ s=ut+\tfrac12 at^2=30\times 5-\tfrac12\times 6\times 25=150-75=\SI{75}{m}. \]
(Check: $v^2=u^2+2as \Rightarrow s=\dfrac{900}{12}=\SI{75}{m}$.)
\item Any two of: greater initial speed (braking distance grows like $u^2$); wet, muddy or oily road reducing friction (common on unpaved roads in the rainy season); worn tyres or brakes; larger mass/loading; downhill slope; longer driver reaction time increasing total stopping distance.
\end{enumerate}
\end{corrige}

\begin{corrige}[Exercise 6 --- Atwood machine]
\begin{enumerate}
\item Diagram: each block carries its weight ($3g$ and $5g$ downwards) and the tension $T$ upwards; the \SI{5}{kg} block accelerates downwards, the \SI{3}{kg} block upwards.
\begin{popfigure}
\begin{center}
\begin{tikzpicture}[scale=0.9,>=stealth]
\draw[fill=gray!20] (0,0) circle (0.4);
\draw[fill=black] (0,0) circle (0.04);
\draw[thick] (-0.4,0) -- (-0.4,-1.5);
\draw[thick] (0.4,0) -- (0.4,-2.2);
\draw[fill=popblueL] (-0.75,-2.1) rectangle (-0.05,-1.5);
\node at (-0.4,-1.8) {\small $3$ kg};
\draw[fill=popblueL] (0.0,-3.0) rectangle (0.8,-2.2);
\node at (0.4,-2.6) {\small $5$ kg};
\draw[->,thick,popdark] (-0.4,-1.5) -- (-0.4,-0.8) node[right]{\small $T$};
\draw[->,thick,popdark] (-0.4,-2.1) -- (-0.4,-2.9) node[right]{\small $3g$};
\draw[->,thick,popdark] (0.4,-2.2) -- (0.4,-1.4) node[right]{\small $T$};
\draw[->,thick,popdark] (0.4,-3.0) -- (0.4,-3.9) node[right]{\small $5g$};
\draw[->,thick,poporange] (1.2,-2.6) -- (1.2,-3.4) node[right]{\small $a$};
\draw[->,thick,poporange] (-1.2,-1.8) -- (-1.2,-1.0) node[left]{\small $a$};
\end{tikzpicture}
\end{center}
\legende{Forces on the two blocks of the Atwood machine.}
\end{popfigure}
\item Taking the direction of motion of each block as positive:
\[ \text{Block } (5\ \mathrm{kg}):\ 5g-T=5a; \qquad \text{Block } (3\ \mathrm{kg}):\ T-3g=3a. \]
\item Adding the two equations eliminates $T$:
\[ 2g=8a \;\Longrightarrow\; a=\dfrac{2\times 9.81}{8}=\SI{2.4525}{m.s^{-2}}\approx\SI{2.45}{m.s^{-2}}. \]
\item Substituting back: $T=3g+3a=3(9.81+2.4525)=3\times 12.2625=\SI{36.79}{N}\approx\SI{36.8}{N}$.
\item Constant acceleration from rest over $s=\SI{1}{m}$:
\[ v^2=u^2+2as=0+2\times 2.4525\times 1=4.905 \;\Longrightarrow\; v\approx\SI{2.21}{m.s^{-1}}. \]
\end{enumerate}
\end{corrige}

\begin{corrige}[Exercise 7 --- Rocket propulsion]
\begin{enumerate}
\item Thrust is the rate of change of momentum of the ejected gas:
\[ F=\dfrac{dm}{dt}\,v_{\text{ex}}=2\times 800=\SI{1600}{N}. \]
By Newton's third law, the gas pushed backwards exerts an equal forward thrust on the rocket.
\item At $t=0$, $m_0=\SI{500}{kg}$:
\[ a_0=\dfrac{F}{m_0}=\dfrac{1600}{500}=\SI{3.2}{m.s^{-2}}. \]
\item Mass decreases at \SI{2}{kg.s^{-1}}: $m(10)=500-2\times 10=\SI{480}{kg}$. Hence
\[ a(10)=\dfrac{1600}{480}=\SI{3.33}{m.s^{-2}}. \]
\item The thrust $F$ is constant, but $a=\dfrac{F}{m(t)}$ with $m(t)=500-2t$ decreasing as fuel is burnt and ejected. A constant force acting on a decreasing mass produces an increasing acceleration. (This is why real rockets accelerate so strongly near the end of a stage's burn.)
\end{enumerate}
\end{corrige}

\begin{corrige}[Exercise 8 --- Forces in a turning car]
\begin{enumerate}
\item Centripetal acceleration:
\[ a_c=\dfrac{v^2}{r}=\dfrac{15^2}{40}=\dfrac{225}{40}=\SI{5.625}{m.s^{-2}}\approx\SI{5.63}{m.s^{-2}}, \]
directed towards the centre of the circle.
\item Resultant force:
\[ F=ma_c=60\times 5.625=\SI{337.5}{N}, \]
directed towards the centre. It is provided by the friction between the seat and the passenger (and the push of the door or seat belt if the turn is sharp).
\item The passenger's body, by inertia, tends to continue in a straight line (Newton's first law) while the car turns inwards. Relative to the car, the body is pressed against the outer side; this \emph{feeling} of being pushed outwards is the effect of inertia, not of a real outward force.
\item Real forces on the passenger: his weight $mg$ (downwards), the normal reaction of the seat (upwards), and the friction/contact force of the seat and door directed towards the centre of the bend. In the rotating frame of the car, a \textbf{fictitious centrifugal force} $m\dfrac{v^2}{r}$ directed outwards must be added to ``explain'' the equilibrium in that non-inertial frame.
\end{enumerate}
\end{corrige}

\begin{corrige}[Exercise 9 --- Connected bodies on a slope (GCE Advanced Level style)]
\textbf{Indicative mark scheme:} (a) 3 marks: two force diagrams, all forces named; (b) 1 mark for $R$, 2 for $F_{\lim}$; (c) 2 marks for both equations of motion, 1 for the comparison showing motion, 2 for $a$; (d) 2 marks; (e) 1 mark for deceleration, 2 for justification.

\begin{enumerate}
\item On $A$: weight $4g$ vertically down, normal reaction $R$ perpendicular to the plane, friction $F$ acting \emph{down} the plane (opposing the impending motion up the plane), tension $T$ up the plane. On $B$: weight $6g$ down, tension $T$ up.
\begin{popfigure}
\begin{center}
\begin{tikzpicture}[scale=0.95,>=stealth]
\draw[thick] (0,0) -- (6,0);
\draw[thick] (0,0) -- (5.2,3.0);
\draw[fill=popblueL,rotate around={30:(2.6,1.5)}] (2.2,1.5) rectangle (3.0,2.3);
\node[rotate=30] at (2.6,1.95) {\small $A$};
\draw[fill=gray!20] (5.2,3.0) circle (0.25);
\draw[thick] (3.0,2.25) -- (5.2,3.25);
\draw[thick] (5.45,3.0) -- (5.45,1.4);
\draw[fill=popblueL] (5.05,0.8) rectangle (5.85,1.4);
\node at (5.45,1.1) {\small $B$};
\draw[->,thick,popdark] (2.6,1.9) -- (2.6,0.4) node[below left]{\small $4g$};
\draw[->,thick,popdark,rotate around={30:(2.6,1.9)}] (2.6,1.9) -- (2.6,3.0) node[above]{\small $R$};
\draw[->,thick,poporange,rotate around={30:(2.6,1.9)}] (2.6,1.9) -- (3.7,1.9) node[above right=-2pt]{\small $T$};
\draw[->,thick,poporange,rotate around={30:(2.6,1.9)}] (2.6,1.9) -- (1.6,1.9) node[below left=-2pt]{\small $F$};
\draw[->,thick,popdark] (5.45,1.4) -- (5.45,2.3) node[right]{\small $T$};
\draw[->,thick,popdark] (5.45,0.8) -- (5.45,-0.1) node[right]{\small $6g$};
\draw (1.1,0) arc (0:30:1.1) node[midway,right=1pt]{\small $30^\circ$};
\end{tikzpicture}
\end{center}
\legende{Forces on block $A$ (on the plane) and block $B$ (hanging).}
\end{popfigure}
\item Perpendicular to the plane there is no acceleration, so
\[ R=4g\cos 30^{\circ}=4\times 9.81\times 0.8660=\SI{33.98}{N}\approx\SI{34.0}{N}. \]
Limiting friction:
\[ F_{\lim}=\mu R=0.2\times 33.98=\SI{6.80}{N}. \]
\item Compare the driving and resisting forces along the direction of possible motion. Component of $A$'s weight down the plane:
\[ 4g\sin 30^{\circ}=4\times 9.81\times 0.5=\SI{19.62}{N}. \]
Total force resisting the descent of $B$: $19.62+6.80=\SI{26.42}{N}$. Weight of $B$: $6g=\SI{58.86}{N}$. Since $58.86>26.42$, the system \emph{does} move, with $B$ descending and $A$ moving up the plane. \hfill$\blacksquare$

Equations of motion (string inextensible: same $a$ for both):
\begin{align*}
B:\quad 6g-T &= 6a,\\
A:\quad T-4g\sin 30^{\circ}-F_{\lim} &= 4a.
\end{align*}
Adding: $6g-4g\sin 30^{\circ}-F_{\lim}=10a$:
\[ a=\dfrac{58.86-19.62-6.80}{10}=\dfrac{32.44}{10}=\SI{3.24}{m.s^{-2}}. \]
\item From $B$'s equation:
\[ T=6g-6a=6(9.81-3.244)=6\times 6.566=\SI{39.4}{N}. \]
\item When the string breaks, the tension vanishes. Along the plane, $A$ (moving up) is now subject to $4g\sin30^{\circ}+F_{\lim}=26.42$ $\mathrm{N}$ directed \emph{down} the plane, giving a deceleration
\[ a'=\dfrac{26.42}{4}=\SI{6.61}{m.s^{-2}}. \]
So $A$ decelerates up the plane, comes instantaneously to rest, and then --- since $4g\sin30^{\circ}=\SI{19.62}{N}>\mu R=\SI{6.80}{N}$ --- it accelerates back \emph{down} the plane with
\[ a''=\dfrac{19.62-6.80}{4}=\SI{3.21}{m.s^{-2}}. \]
(Examiner's expectation: the candidate must note that friction reverses direction when the motion reverses.)
\end{enumerate}
\end{corrige}

\begin{corrige}[Exercise 10 --- Lift, reaction and apparent weight (GCE Advanced Level style)]
\textbf{Indicative mark scheme:} (a) 1 mark for the statement, 2 for correct application; (b) 2, 3, 3 marks with method, substitution, answer; (c) 3 marks for reasoning with $N$; (d) 1 mark for the value, 2 for the comment.

\begin{enumerate}
\item \textbf{Newton's second law:} in an inertial frame, the resultant force on a body equals the product of its mass and acceleration: $\vec{F}=m\vec{a}$. The man is carried by the lift, so he has the same acceleration $a$ as the lift. Two real forces act on him: his weight $mg$ (downwards) and the normal reaction $N$ of the scale (upwards). Taking upwards positive: $N-mg=ma$. The scale reads $N$, the \emph{apparent weight}.
\item \begin{enumerate}
\item Constant velocity: $a=0$, so $N=mg=70\times 9.81=\SI{686.7}{N}$.
\item Upward acceleration $a=\SI{1.5}{m.s^{-2}}$:
\[ N=m(g+a)=70\times(9.81+1.5)=70\times 11.31=\SI{791.7}{N}. \]
\item Downward acceleration: taking upwards positive, $a=-1.5$:
\[ N=m(g-1.5)=70\times 8.31=\SI{581.7}{N}. \]
\end{enumerate}
\item The sensation of ``heaviness'' is the compression of the body against the scale, i.e. the magnitude of $N$. In (ii) the scale must both support the weight \emph{and} provide the upward acceleration, so $N>mg$: the man feels heavier. In (iii) part of the weight is ``used'' to produce the downward acceleration, so $N<mg$: he feels lighter. The true weight $mg$ never changes.
\item If the cable snaps, the lift and the man fall freely with $a=g$. Then
\[ N=m(g-g)=0. \]
The scale reads \textbf{zero}. This is \emph{apparent weightlessness}: the man still has weight $mg$ (gravity still acts --- that is what accelerates him!), but no contact force is exerted on him because he and the scale fall together. ``Weightlessness'' therefore means absence of \emph{support reaction}, not absence of gravity.
\end{enumerate}
\end{corrige}

\begin{corrige}[Exercise 11 --- Essay: scope and limits of Newton's laws (GCE Advanced Level style)]
\textbf{Examiner's expectations:} precise statements, correctly identified pairs, correct naming of theories, and a balanced, structured argument. Award marks for physical insight, not length.

\begin{enumerate}
\item The statement is \textbf{false}, and rests on a confusion about the third law. An action--reaction pair acts on \emph{two different bodies}: the two forces can therefore never cancel, since cancellation (a zero resultant) is only meaningful for forces acting on the \emph{same} body. Whether a body accelerates depends solely on the resultant of the forces acting \emph{on that body} (Newton's second law).

\emph{Concrete example --- walking:} the walker pushes the ground backwards with his foot (action on the Earth); the ground pushes the walker forwards (reaction on the walker). It is this reaction --- friction of the ground on the foot --- that is the external force accelerating the walker forwards. The two forces are equal and opposite, but one acts on the Earth and the other on the walker, so the walker accelerates. (Similarly, a pirogue pushed from the bank: the pole pushes the bank, the bank pushes the pirogue forwards.)

\item \begin{enumerate}
\item \textbf{Speeds close to $c$:} Newtonian mechanics assumes mass is constant and that velocities add simply. At speeds comparable to the speed of light, momentum is no longer $m\vec{v}$ but $p=\dfrac{mv}{\sqrt{1-v^2/c^2}}$, and $F=ma$ fails to predict observed motion (e.g. electrons in particle accelerators require energies far exceeding the Newtonian prediction). The extending theory is \textbf{Einstein's special relativity}, which reduces to Newton's laws when $v\ll c$.
\item \textbf{Atomic and subatomic scales:} particles such as electrons do not follow definite trajectories; their behaviour is governed by probabilities and wave functions, and energy is quantised. Newton's laws cannot, for instance, explain the stability of atoms or the discrete spectra of hydrogen. The replacing theory is \textbf{quantum mechanics}, whose predictions agree with Newtonian mechanics for macroscopic bodies (correspondence principle).
\end{enumerate}

\item The Earth rotates on its axis and orbits the Sun, so a frame fixed to its surface is accelerating and is only \emph{approximately} inertial; strictly, fictitious forces (centrifugal, Coriolis) should be added.

\emph{Where it matters:} for large-scale or long-duration motions --- e.g. the deflection of trade winds and ocean currents, the rotation of the plane of a Foucault pendulum, or long-range artillery --- the Coriolis effect due to the Earth's rotation is measurable and must be included.

\emph{Where it can be ignored:} for everyday, short-duration, small-scale motions --- a lorry braking on the Douala--Yaound\'e road, a falling mango, a lift in a building --- the correction due to the Earth's rotation is utterly negligible compared with other uncertainties, so Newton's laws applied in the Earth frame give excellent results. This is why Newtonian mechanics remains the working tool of engineers in Cameroon and worldwide.
\end{enumerate}
\end{corrige}

\newpage

\coursec{Summary sheet}

\begin{popfigure}
\begin{tikzpicture}[
  central/.style={circle, draw=popdark, fill=popblueL, thick, minimum size=1.4cm, text width=2.2cm, align=center, font=\bfseries},
  branch/.style={rectangle, draw=popdark, rounded corners, thick, minimum width=2.5cm, minimum height=1cm, text width=2.3cm, align=center, font=\small},
  line/.style={thick, -}
]
\node[central] (center) {Newton's Laws of Motion};
% Branches
\node[branch, fill=poporange] (b1) at (4,0) {First Law (Inertia)};
\node[branch, fill=popgreen] (b2) at (2,3.46) {Second Law $\sum\vec F=m\vec a$};
\node[branch, fill=poppurple] (b3) at (-2,3.46) {Third Law (Action--Reaction)};
\node[branch, fill=poppink] (b4) at (-4,0) {Applications (FBD, Equilibrium)};
\node[branch, fill=popgold] (b5) at (-2,-3.46) {Limitations (High speed, Quantum)};
\node[branch, fill=popblue] (b6) at (2,-3.46) {Extensions (Impulse, Momentum)};
\draw[line, poporange] (center) -- (b1);
\draw[line, popgreen] (center) -- (b2);
\draw[line, poppurple] (center) -- (b3);
\draw[line, poppink] (center) -- (b4);
\draw[line, popgold] (center) -- (b5);
\draw[line, popblue] (center) -- (b6);
\end{tikzpicture}
\legende{Mind map of Newton's Laws of Motion: the three laws at the centre, with their applications, limitations and extensions.}
\end{popfigure}

\begin{bilan}
\textbf{Key Definitions}
\begin{itemize}
\item \cle{Inertia}: the inherent property of a body to resist any change in its state of rest or of uniform motion in a straight line. Mass is the quantitative measure of inertia.
\item \cle{Force}: a vector interaction (push or pull) that can change the state of motion or the momentum of a body; measured in newtons (N), with $1\ \mathrm{N}=1\ \mathrm{kg\,m\,s^{-2}}$.
\item \cle{Mass}: quantitative measure of inertia; a scalar quantity with SI unit the kilogram (kg). It is constant for a given body in classical mechanics.
\item \cle{Acceleration}: rate of change of velocity with time; $\vec a = \dfrac{d\vec v}{dt}$. It is a vector, with SI unit $\mathrm{m\,s^{-2}}$.
\item \cle{Momentum} (linear momentum): $\vec p = m\vec v$; a vector quantity, unit $\mathrm{kg\,m\,s^{-1}}$, conserved in any isolated system.
\item \cle{Impulse}: $\vec I = \displaystyle\int_{t_1}^{t_2} \vec F\,dt = \Delta\vec p$; for a constant force, $\vec I = \vec F\,\Delta t$. It equals the change in momentum and is the area under the $F$--$t$ graph.
\item \cle{Action--Reaction pair}: two forces $\vec F_{12}$ and $\vec F_{21}$ of the same nature, equal in magnitude, opposite in direction, acting along the same line but on \emph{different} bodies.
\item \cle{Free-body diagram (FBD)}: a sketch showing \emph{all} the external forces acting on a single isolated body, and nothing else.
\item \cle{Equilibrium}: a body is in translational equilibrium when $\sum\vec F = \vec 0$; it is then at rest or moving with constant velocity (First Law).
\item \cle{Weight}: $\vec W = m\vec g$, the gravitational pull of the Earth on the body, where $\vec g$ is the gravitational field strength ($g \approx 9.81\ \mathrm{m\,s^{-2}}$ near the Earth's surface).
\item \cle{Normal reaction}: the contact force exerted by a surface on a body, perpendicular to the surface; it adjusts to prevent interpenetration.
\item \cle{Friction}: force opposing relative motion (kinetic friction, $F_f = \mu_k R$) or impending motion (static friction, $F_f \le \mu_s R$), where $R$ is the normal reaction.
\item \cle{Tension}: pulling force transmitted along a light, inextensible string; it has the same magnitude throughout the string if the string is light and any pulley is smooth.
\end{itemize}

\textbf{Laws / Formulas to Know}
\begin{itemize}
\item \textbf{Newton's First Law (Law of Inertia)}: a body remains at rest or in uniform motion in a straight line unless acted upon by a resultant external force:
\[ \sum \vec F = \vec 0 \;\Longleftrightarrow\; \vec v = \text{constant} \quad (\text{including } \vec v = \vec 0). \]
\item \textbf{Newton's Second Law}: the resultant force equals the rate of change of momentum:
\[ \sum \vec F = \frac{d\vec p}{dt}. \]
For a body of constant mass this reduces to $\sum \vec F = m\vec a$.
\item \textbf{Newton's Third Law}: if body 1 exerts a force $\vec F_{12}$ on body 2, then body 2 exerts $\vec F_{21}$ on body 1 with
\[ \vec F_{12} = -\vec F_{21}; \]
the two forces act on \emph{different} bodies and therefore never cancel on a single free-body diagram.
\item \textbf{Impulse--Momentum Theorem}: $\vec I = \Delta\vec p = m\vec v - m\vec u$; essential for collisions, impacts and variable forces.
\item \textbf{Conservation of Linear Momentum}: in an isolated system ($\sum\vec F_{\text{ext}} = \vec 0$),
\[ \sum\vec p_{\text{before}} = \sum\vec p_{\text{after}}. \]
\item \textbf{Work--Energy Theorem}: $W_{\text{net}} = \Delta K = \dfrac{1}{2}m(v^{2}-u^{2})$.
\item \textbf{Centripetal Force} (uniform circular motion): $F_c = \dfrac{mv^{2}}{r} = mr\omega^{2}$, directed towards the centre of the circle; it is the \emph{resultant} of real forces (tension, friction, reaction, weight component).
\item \textbf{Newton's Law of Gravitation}: $F = G\dfrac{m_1 m_2}{r^{2}}$; useful background for understanding weight and the variation of $g$.
\end{itemize}

\textbf{Problem-Solving Methods}
\begin{enumerate}
\item \textbf{Read and visualise}: identify the bodies involved, the motion described, and any constraints (strings, pulleys, surfaces, lifts).
\item \textbf{Draw free-body diagrams}: for \emph{each} body separately, sketch all external forces with labels (weight $mg$, normal reaction $R$, friction $F_f$, tension $T$, applied forces).
\item \textbf{Choose coordinate axes}: align one axis with the direction of motion or acceleration (e.g. along an inclined plane); resolve every other force into components.
\item \textbf{Write Newton's Second Law in components}: $\sum F_x = ma_x$, $\sum F_y = ma_y$. For equilibrium, set $a_x = a_y = 0$.
\item \textbf{Apply constraints}: for connected particles joined by a light inextensible string, the accelerations have equal magnitude and the tension is uniform throughout the string.
\item \textbf{Solve the resulting equations}: use simultaneous equations for connected bodies; keep all quantities in SI units.
\item \textbf{Check the answer}: verify dimensions, sign and direction, and physical plausibility (e.g. required friction must not exceed $\mu_s R$; a tension cannot be negative).
\item \textbf{For impulse problems}: compute $\Delta\vec p$ from the velocities before and after, or find the area under the $F$--$t$ graph; take care with signs when a body rebounds.
\end{enumerate}

\textbf{Common Mistakes (Traps to Avoid)}
\begin{itemize}
\item Confusing \emph{mass} (kg) with \emph{weight} (N): $W = mg$ is a force; never add a mass to a force.
\item Treating action--reaction forces as cancelling on a single free-body diagram: they act on \emph{different} bodies.
\item Believing a force is needed to \emph{maintain} uniform motion: the First Law says a resultant force is needed only to \emph{change} the velocity.
\item Forgetting to resolve forces into components when the motion is not along the force direction (especially on inclined planes, where the weight gives $mg\sin\theta$ down the plane and $mg\cos\theta$ into the plane).
\item Assuming the tension is the same on both sides of a pulley that is massive or rough.
\item Using $v^{2} = u^{2} + 2as$ or $s = ut + \tfrac12 at^{2}$ when the acceleration is not constant.
\item Getting the direction of friction wrong: it always opposes relative motion or the \emph{tendency} to move.
\item Mixing scalar and vector quantities (e.g. equating speed with velocity, or adding momenta without regard to direction and sign).
\item Omitting units in final answers; GCE examiners penalise missing or incorrect units.
\end{itemize}

\textbf{GCE Examination Tips}
\begin{itemize}
\item Always \emph{state the law} you are applying before writing equations (e.g. ``By Newton's Second Law applied to the particle\ldots'').
\item Define every symbol at first use: ``Let $T$ be the tension in the string (in N)''.
\item Use vector notation ($\vec F$, $\vec a$) when the question asks for vector form; otherwise clear component form with a stated positive direction is acceptable.
\item In equilibrium questions, write explicitly $\sum F_x = 0$ and $\sum F_y = 0$; do not merely say ``the forces balance''.
\item For momentum and impulse questions, draw a clear before/after diagram with velocity arrows and a marked positive direction.
\item Practise past GCE papers: typical contexts include blocks on rough inclined planes, lift (apparent weight) problems, connected particles over pulleys, and collisions.
\item Time management: allocate about 12--15 minutes per mechanics question; do not spend too long on a single part.
\item If stuck, write down the relevant formulas and the given data; method marks are awarded for a correct approach even without the final answer.
\end{itemize}
\end{bilan}

\findecours

\end{document}
